\section{Theta lifts between definite unitary groups}\label{sec:pf:tlift}

\providecommand{\tliftS}{\mathcal{S}}
\providecommand{\tliftPol}{\mathcal{P}}
\providecommand{\tliftPD}{\mathcal{P}_{\A}}
\providecommand{\tliftAE}{\A_E}
\providecommand{\tliftdia}{\Delta}
\providecommand{\tliftK}{K_2}
\providecommand{\tliftKfr}{K_{\mathrm{fr}}}
\providecommand{\tliftht}{\mathsf{a}}
\providecommand{\tliftlift}{\Theta}
\providecommand{\tliftZ}{Z}
\providecommand{\tliftM}{M}
\providecommand{\tliftLW}{\mathrm{LW}}
\providecommand{\tliftLP}{\mathrm{LP}}
\providecommand{\tliftInt}{K^{\mathrm{int}}}
\providecommand{\tliftGinf}{G_\infty}
\providecommand{\tliftGS}{G_S}
\providecommand{\tliftGoff}{G^{S}}
\providecommand{\tliftvac}{\varphi^{\circ}}
\providecommand{\tliftsymp}{\langle\cdot,\cdot\rangle}
\providecommand{\tliftV}{V}
\providecommand{\tliftPP}{\mathbf{P}}
\providecommand{\tliftE}{\widehat{E}}

This section proves \cref{thm:galois:tlift}: the statement $\sfname{TLiftK}$ of
\cref{def:galois:tliftk}, an existence theorem for global theta lifts from a totally definite unitary
group $\Uni(X)$ of rank $n$ to totally definite unitary groups $\Uni(W)$ of every rank $d$ beyond a
threshold $d_0$, with control of the archimedean weights and of the local components at the split
places. It is used in step (iv) of the proof of \cref{thm:galois:reduction}(b), to pass from odd to even
rank.

\subsubsection*{The idea}
Let $\omega$ be a Weil representation of $\Uni(X)(\A)\times\Uni(W)(\A)$ on the Schwartz--Bruhat space of
$M_{n\times d}(\A)$ (with $\A$ the adeles of the totally real field $F=E^+$), split by Hecke
characters $\chi_X,\chi_W$ \cite{Weil1964,Kudla1994}, and let $\theta(g,h;\Phi)$ be the associated
theta series. For an automorphic form $\varphi$ on $\Uni(X)$ the theta lift
\[
  \tliftlift(\Phi,\varphi)(h)=\int_{\Uni(X)(F)\backslash\Uni(X)(\A)}\theta(g,h;\Phi)\,\overline{\varphi(g)}\,dg
\]
is an automorphic form on $\Uni(W)$, and it depends equivariantly on $(\Phi,\varphi)$. Three things have
to be shown.
\begin{enumerate}[label=(\arabic*),nosep]
\item \emph{Archimedean matching.} In the Fock model of $\omega$ at the infinite places (where both
  groups are compact) one finds a polynomial Fock vector $f_0$ which generates, under $\Uni(X)(F_\infty)$,
  the same irreducible type as a vector $\varphi$ of the given automorphic representation $U$, and whose
  $\Uni(W)(F_\infty)$-type has the prescribed highest weight $\lambda'$. Here the joint highest weight
  vectors of the pair $(\GL_n,\GL_d)$ acting on polynomials on $M_{n\times d}$ are used
  \cite{KashiwaraVergne1978,Howe1989}.
\item \emph{Non-vanishing.} For $d\ge d_0$ the doubling integral
  $\int_{\Uni(X)(\A)}\langle\omega(g)\Phi,\Phi\rangle\,\overline{\langle\varphi(\cdot\,g),\varphi\rangle}\,dg$ converges
  absolutely, and it vanishes if $\tliftlift(\Phi,\varphi)=0$: this is the Rallis inner product
  argument \cite{Rallis1984,PiatetskiShapiroRallis1987,Li1992}, deduced from the Siegel--Weil formula on the
  doubled group $\Uni(X\oplus(-X))$ \cite{Weil1965,KudlaRallis1988,Ichino2007}. On the other hand, the
  doubling integral is \emph{not} zero for $\Phi=f_0\otimes\bigotimes_v\phi_v$ with suitable ``bump
  vectors'' $\phi_v$ at finitely many places: the integrand vanishes off a small open subgroup at these
  places, the archimedean integral is computed by Schur orthogonality, and the contribution of the
  remaining places is controlled by a tail estimate. The threshold $d_0$ comes from an integrability
  exponent of a height function on the doubled group.
\item \emph{Reading off the lift.} A nonzero lift has a nonzero equivariant projection to an
  automorphic representation $U'$ of $\Uni(W)$; $U'$ has the weights determined by $f_0$, and at a split
  place the equivariance of the lift produces a nonzero equivariant map from the Schwartz space to
  antilinear maps between the local components, i.e.\ the type II local theta correspondence
  \cite{Minguez2008}.
\end{enumerate}
The formal proof \cite{LeanSource2026} follows this plan. It proves everything it uses about the Weil
representation, the Siegel--Weil formula and the Fock model; we state these results precisely in
\cref{pf:tlift:tools}, compare them with their published forms in \cref{pf:tlift:classical}, and convert the
arguments that are specific to this proof.

\subsubsection*{Organisation}
\Cref{pf:tlift:setting} defines the objects in the statement and restates \cref{thm:galois:tlift}.
\Cref{pf:tlift:tools} introduces Schwartz spaces, Heisenberg operators, Weil representations, theta
series and the polynomial Fock space, and states the analytic and representation-theoretic tools
(including the Siegel--Weil formula). \Cref{pf:tlift:classical} gives the published forms of the Weil
representation and of the Siegel--Weil formula, the differences with the formal statements, and the structure of
their formal proofs. \Cref{pf:tlift:dblgrp} treats the doubled group: compactness of the standard maximal compact
subgroup, the Iwasawa decomposition, heights and the integrability exponent. \Cref{pf:tlift:rallis}
proves the Rallis inner product statement. \Cref{pf:tlift:nonvan} proves the non-vanishing of the
doubling integral. \Cref{pf:tlift:real} constructs the archimedean data, \cref{pf:tlift:autsec} shows that
lifts are automorphic forms, \cref{pf:tlift:read} projects and reads off the lift, and
\cref{pf:tlift:proof} assembles the proof.

% ---------------------------------------------------------------------------
\subsection{Setting and statement}\label{pf:tlift:setting}

Throughout, $E$ is a CM field, $F=E^+$ its maximal totally real subfield, and $x\mapsto\bar x$ the
complex conjugation of $E$; it is also used for the induced automorphisms of $E\otimes_FR$ for
$F$-algebras $R$, and on matrices entrywise. We write $\A=\A_F$, $\A_f$ for the finite adeles,
$F_\infty=F\otimes_\Q\R$, and $\tliftAE=\A\otimes_FE$. For a finite place $v$ of $F$ (a nonzero prime of
$\OO_F$) let $F_v$, $\OO_v$ and $\Nm v$ be the completion, its valuation ring and the norm. For an
infinite place $w$ of $F$ we fix below an embedding $\tau_w\colon E\to\C$ above $w$; then
$E\otimes_FF_w\cong\C$ via $\tau_w$, and we write $\tau_w$ also for the induced maps $\tliftAE\to\C$ and on
matrices. If a finite place $v$ splits in $E$, an $F$-embedding $t\colon E\to F_v$ determines the place
$w_t$ of $E$ above $v$ at which $t$ is continuous, an isomorphism $F_v\cong E_{w_t}$, and a projection
$\pi_t\colon\tliftAE\to F_v$. Let $\eta$ be the quadratic character of $\A^\times$ which is trivial
exactly on $F^\times\Nm(\tliftAE^\times)$.

\begin{definition}[Good CM data]\label{pf:tlift:good}
\emph{Base data} for $E$ consist of a character $\psi\colon\A\to\C^\times$; for each infinite place $w$ of
$F$ an embedding $\tau_w\colon E\to\C$ above $w$; a finite set $P$ of finite places of $F$ with an
$F$-embedding $t_v\colon E\to F_v$ for each $v\in P$; and some further entries (two integer-valued
functions on the infinite places and one more finite place with an embedding), which play no role in
this section. The base data are \emph{good} if:
\begin{enumerate}[label=(G\arabic*),nosep]
\item $E/F$ is a CM extension: $F$ is totally real, $E$ is totally complex, $[E:F]=2$, and
  $\tau(\bar x)=\overline{\tau(x)}$ for every $\tau\colon E\to\C$ and $x\in E$;
\item $E$ contains a square root of $-D$ for some integer $D>0$;
\item $E/F$ is unramified at every finite place;
\item $\psi$ is \emph{standard}: continuous, trivial on $F$, and $\psi(a)=\exp(-2\pi i\Tr_{F_\infty/\R}a)$
  for every $a\in\A$ whose finite components vanish;
\item every finite place $v$ of $F$ which is inert in $E$ satisfies $v\nmid2$ and $\Nm v>3$, and
  $\psi_v=\psi|_{F_v}$ is unramified at $v$, i.e.\ $\{x\in F_v:\psi_v(x\OO_v)=1\}=\OO_v$;
\item $\psi_v$ is unramified at every $v\in P$.
\end{enumerate}
\formal{UpperTheory.TP.W.GoodK}
\end{definition}

For $J\in M_m(E)$ and $x,y\in E^m$ (or in $\tliftAE^m$) we put $h_J(x,y)={}^tx\,J\,\bar y$.

\begin{definition}[Hermitian and skew-hermitian spaces]\label{pf:tlift:spaces}
Fix base data. A \emph{space of dimension $m$} is a tuple $Y=(J,L,\vartheta,b)$ with $J\in M_m(E)$, an
$\OO_E$-submodule $L\subset E^m$, matrices $\vartheta_w\in M_m(\C)$ for the infinite places $w$ of $F$, and
$b_v\in\GL_m(F_v)$ for the finite places $v$. It is a \emph{hermitian space of signature $(r,s)$} if
\begin{enumerate}[label=(H\arabic*),nosep]
\item ${}^t\bar J=J$;
\item for every $w$ the hermitian matrix $\tau_w(J)$ has signature $(r,s)$: $r+s=m$ and
  ${}^tA\,\tau_w(J)\,\bar A=\mathrm{diag}(1_r,-1_s)$ for some $A\in M_m(\C)$;
\item $L$ is a lattice in $E^m$;
\item $L$ is self-dual: $L=\{y\in E^m:h_J(y,x)\in\OO_E\text{ for all }x\in L\}$;
\item $\vartheta_w$ is a Cartan involution of $\tau_w(J)$: $\vartheta_w^2=1$,
  ${}^t\vartheta_w\,\tau_w(J)\,\bar\vartheta_w=\tau_w(J)$ and $\Re\,h_{\tau_w(J)}(\vartheta_wz,z)>0$ for $z\ne0$;
\item for $v\in P$, $b_v$ is a basis of the $\OO_v$-lattice $t_v(L\otimes_{\OO_F}\OO_v)\subset F_v^m$.
\end{enumerate}
It is a \emph{skew-hermitian space of signature $(r,s)$} if (H3) and (H6) hold and
\begin{enumerate}[label=(S\arabic*),nosep]
\item ${}^t\bar J=-J$;
\item[(S2)] $i\,\tau_w(J)$ has signature $(r,s)$ for every $w$;
\item[(S4)] $L=\{y:\Tr_{E/\Q}(o\,h_J(y,x))\in\Z\text{ for all }x\in L,\ o\in\OO_E\}$;
\item[(S5)] $\vartheta_w$ is a Cartan involution of $i\,\tau_w(J)$.
\end{enumerate}
A space is \emph{definite} if $\min(r,s)=0$.
\formal{GlobalTheory.IndefiniteRowSideFullHecke.WeilTensorModels.Spc}
\end{definition}

If $Y$ is definite, then $\vartheta_w=\pm1$: an involution preserving a definite form has no eigenvector
$z$ of eigenvalue $\mp1$ by the positivity in (H5). The space $(1_n,L,1,b)$ with $L$ self-dual for the
standard form is a positive definite hermitian space; this is the space $X$ of the statement below.

\subsubsection*{Unitary groups and automorphic representations}
For a nondegenerate $J\in M_m(E)$ and an $F$-algebra $R$ let
$\Uni_J(R)=\{g\in\GL_m(R\otimes_FE):{}^tg\,J\,\bar g=J\}$, and for a space $Y$ write
$\Uni(Y)=\Uni_{J_Y}$. The group $\Uni_J(\A)$ is a closed subgroup of $\GL_m(\tliftAE)$, locally compact,
second countable and Hausdorff, and
$\Uni_J(F)=\Uni_J(\A)\cap\GL_m(E)$.

The statement uses the \emph{standard} model $\Uni_n(\A)=\{g\in\GL_n(\A_E):{}^tg\,\bar g=1\}$, where $\A_E$ is
the adele ring of $E$. It contains the subgroups $\Uni_n(E)$ of rational points, $\Uni_n(E_\infty)$ of
elements with trivial finite components, $\Uni_n(\A_f)$ of elements with trivial infinite components,
for each infinite place $w$ of $F$ the subgroup $\Uni_n(F_w)$ of elements of $\Uni_n(E_\infty)$ which are
trivial at the infinite places other than $w$, with the isomorphism $k\mapsto\tau_w(k)$ onto the compact
group $\Uni(n)\subset\GL_n(\C)$, and for each finite place $w$ of $E$ the subgroup $\Uni_n(F_v)$,
$v=w|_F$, of elements of $\Uni_n(\A_f)$ with trivial components at the places of $E$ not above $v$, with
the projection $c_w\colon\Uni_n(F_v)\to\GL_n(E_w)$. If $v$ splits, $c_w$ is an isomorphism of
topological groups. We write $R(x)\varphi=\varphi(\,\cdot\,x)$ for right translation.

\begin{definition}[Automorphic forms and representations; weights; local components]\label{pf:tlift:aut}
Let $n\ge1$.
\begin{enumerate}[label=(\alph*),nosep]
\item An \emph{automorphic form} on $\Uni_n$ is a continuous function $\varphi$ on $\Uni_n(\A)$ with
  $\varphi(\gamma g)=\varphi(g)$ for $\gamma\in\Uni_n(E)$, such that the elements $k\in\Uni_n(\A_f)$ with
  $R(k)\varphi=\varphi$ form an open subset of $\Uni_n(\A_f)$, and such that the right translates
  $R(k)\varphi$, $k\in\Uni_n(E_\infty)$, span a finite-dimensional space.
\item An \emph{automorphic representation} of $\Uni_n$ is a nonzero space $U$ of automorphic forms which
  is stable under all $R(x)$, $x\in\Uni_n(\A)$, and has no invariant subspaces other than $0$ and $U$.
\item For a dominant weight $\lambda\in\Z^n$ ($\lambda_1\ge\dots\ge\lambda_n$) let $V_\lambda$ be the
  irreducible algebraic representation of $\GL_n(\C)$ of highest weight $\lambda$, realised as a
  space of polynomial functions on $M_n(\C)$ (the dual Weyl module of the dominant weight
  $\lambda-\lambda_n$), on which $g$ acts by left substitution twisted by $\det(g)^{\lambda_n}$.
  $U$ has \emph{weight $\lambda$ at $\tau_w$} if there is a
  nonzero linear map $j\colon V_\lambda\to U$ with $j(V_\lambda(\tau_w(k))x)=R(k)j(x)$ for
  $k\in\Uni_n(F_w)$.
\item Let $w$ be a finite place of $E$ and $v=w|_F$. A smooth irreducible representation $\tau$ of
  $\GL_n(E_w)$ is a \emph{local component} of $U$ at $w$ if there is a nonzero linear map
  $f\colon\tau\to U$ with $f(\tau(c_w(x))y)=R(x)f(y)$ for $x\in\Uni_n(F_v)$.
\end{enumerate}
\formal{GlobalTheory.UnitaryBC.IsAutRepU}
\end{definition}

\begin{definition}[Splitting characters]\label{pf:tlift:chars}
A \emph{unitary Hecke character} of $E$ is a continuous homomorphism $\chi\colon\tliftAE^\times\to S^1$
trivial on $E^\times$. For $m\in\N$ we write $\chi|_{\A^\times}=\eta^m$ if $\chi(a)=1$ for
$a\in F^\times\Nm(\tliftAE^\times)$ and $\chi(a)=(-1)^m$ for the other $a\in\A^\times$. A family
$(m(w))_w$ of integers indexed by the infinite places of $F$ is the family of \emph{archimedean exponents}
of $\chi$ if
\[
  \chi(z)=\prod_w\Bigl(\frac{\tau_w(z)}{|\tau_w(z)|}\Bigr)^{m(w)}
\]
for every $z\in\tliftAE^\times$ whose finite components are $1$. If $v$ splits and $t\colon E\to F_v$, we
write $\chi_t$ for the character of $F_v^\times$ obtained from the component of $\chi$ at $w_t$ and the
isomorphism $F_v\cong E_{w_t}$.
\formal{GlobalTheory.IndefiniteRowSideFullHecke.WeilTensorModels.HasExp}
\end{definition}

For a split place $v$, $t\colon E\to F_v$, characters $\nu,\beta$ of $F_v^\times$ and
$\Phi\colon M_{n\times d}(F_v)\to\C$, $g\in\GL_n(F_v)$, $h\in\GL_d(F_v)$, put
\begin{equation}\label{eq:tlift:omA}
  \bigl(\omega^{\nu,\beta}(g,h)\Phi\bigr)(x)=\nu(\det g)^{-1}\beta(\det h)\,|\det g|_v^{-d/2}\,
  |\det h|_v^{-n/2}\,\Phi\bigl(g^{-1}x\,{}^th^{-1}\bigr),
\end{equation}
where $|\cdot|_v$ is the normalised absolute value. Let
$\tliftS(M_{n\times d}(F_v))$ be the space of locally constant compactly supported functions.

\begin{definition}[Lifts]\label{pf:tlift:islift}
Let $\lambda\in\Z^n$, $\lambda'\in\Z^d$ be dominant and $m_X,m_W\in\Z$. The relation
$\tliftLW(\lambda,m_X,m_W,\lambda')$ means
\begin{equation}\label{eq:tlift:liftwt}
  2\lambda'_k=\begin{cases}2\lambda_k-d-m_W+n+m_X,&k\le n,\\ n+m_X,&n<k\le d.\end{cases}
\end{equation}
For a split place $v$, $t\colon E\to F_v$, unitary Hecke characters $\chi_X,\chi_W$ and smooth irreducible
representations $\tau$ of $\GL_n(E_{w_t})$ and $\tau'$ of $\GL_d(E_{w_t})$, the relation
$\tliftLP_t(\chi_X,\chi_W,\tau,\tau')$ means: there is a nonzero linear map $\beta$ from
$\tliftS(M_{n\times d}(F_v))$ to the antilinear maps $\tau\to\tau'$ such that
\begin{equation}\label{eq:tlift:locpair}
  \beta\bigl(\omega^{\chi_{W,t}^{-1},\chi_{X,t}}(g,h)\Phi\bigr)\circ\tau(g)=\tau'(h)\circ\beta(\Phi)
  \qquad(g\in\GL_n(F_v),\ h\in\GL_d(F_v),\ \Phi\in\tliftS),
\end{equation}
with $\GL_n(F_v)\cong\GL_n(E_{w_t})$. Given base data, an automorphic representation $U$ of $\Uni_n$
with weights $\lambda_w$ at $\tau_w$, and $\chi_X,\chi_W,m_X,m_W$, a space $U'$ of functions on
$\Uni_d(\A)$ is a \emph{lift} of $U$ if $U'$ is an automorphic representation of $\Uni_d$; for every
infinite place $w$ it has a weight $\lambda'_w$ at $\tau_w$ with
$\tliftLW(\lambda_w,m_X(w),m_W(w),\lambda'_w)$; and $\tliftLP_t(\chi_X,\chi_W,\tau,\tau')$ holds for every
split place $v$, every $t\colon E\to F_v$ and all local components $\tau$ of $U$ and $\tau'$ of $U'$ at
$w_t$. \formal{UpperTheory.TP.IsLift}
\end{definition}

In terms of $\lambda_w=\mu_w+a_w$ with $2a_w=d+m_W(w)$ and $2b_w=n+m_X(w)$, \eqref{eq:tlift:liftwt}
says $\lambda'_w=(\mu_w,0,\dots,0)+b_w$, the Kashiwara--Vergne correspondence of highest weights
\cite{KashiwaraVergne1978}. Relation \eqref{eq:tlift:locpair} is the statement that $(\tau,\tau')$ occurs
in the type II local theta correspondence for the Weil representation \eqref{eq:tlift:omA} of
$\GL_n\times\GL_d$ \cite{Minguez2008}.

\begin{theorem}[\cref{thm:galois:tlift}]\label{pf:tlift:main}
Let $E$ be a CM field with good base data $(\psi,(\tau_w),P,\dots)$, let $n\ge1$ and let $X$ be an
$n$-dimensional positive definite hermitian space with $J_X=1_n$. There is $d_0\in\N$ with the following
property. Let $d\ge d_0$, let $W$ be a $d$-dimensional skew-hermitian space of signature $(d,0)$ with
$J_W=\delta\cdot1_d$ for some $\delta\in E$, let $\chi_X,\chi_W$ be unitary Hecke characters with
$\chi_X|_{\A^\times}=\eta^n$, $\chi_W|_{\A^\times}=\eta^d$ and archimedean exponents $m_X,m_W$, and let $U$ be
an automorphic representation of $\Uni_n$ with weights $\lambda_w$ at $\tau_w$ such that
\begin{equation}\label{eq:tlift:dom}
  d+m_W(w)\le2\lambda_{w,k}\qquad\text{for all infinite places $w$ and all $k$.}
\end{equation}
Then there is a lift $U'$ of $U$ (\cref{pf:tlift:islift}).
\formallabel{UpperTheory.TP.K.tLiftK}\formallabel{UpperTheory.TP.tLift_holds_k}
\end{theorem}

\begin{remark}[On the formal statement]\label{pf:tlift:remstmt}
\begin{enumerate}[label=(\alph*),nosep]
\item The second tag is the formal proof; its statement is a verbatim copy of $\sfname{TLiftK}$.
\item The threshold $d_0$ depends only on $E$, the base data, $n$ and $X$. In the proof
  $d_0=\max(N,2n)$ for an integer $N>2t_0$, where $t_0$ is an integrability exponent of a height function
  on the doubled group of $X$ (\cref{pf:tlift:hgtfst}). The formal statement gives no explicit value; tracing the
  constants of the proof shows that $d_0=5+16n^2[F:\Q]$ is admissible (\cref{pf:tlift:remdzero}).
\item The archimedean exponents are imposed on all archimedean ideles $z$ (those with finite components $1$);
  since $E_\infty^\times\cong(\C^\times)^{[F:\Q]}$ is connected, this is equivalent to imposing them near $1$. The relation
  $\tliftLP$ is required for all pairs of local components (no uniqueness of local components is
  asserted or needed). No cuspidality is involved: every automorphic form on the compact quotient
  $\Uni_n(E)\backslash\Uni_n(\A)$ is admitted.
\item Of the base data only $\psi$, the embeddings $\tau_w$, the set $P$ and the embeddings $t_v$ enter;
  $P$ enters only through (G6) and (H6).
\end{enumerate}
\end{remark}


% ---------------------------------------------------------------------------
\subsection{Weil representations, theta series and the tools}\label{pf:tlift:tools}

In this subsection $(\psi,(\tau_w),P,\dots)$ are good base data, $X$ is a hermitian space of dimension $n$
and $W$ a skew-hermitian space of dimension $d$ (\cref{pf:tlift:spaces}), and $\kappa=\{1,\dots,n\}\times
\{1,\dots,d\}$. For a ring $R$ we identify $R^\kappa$ with $M_{n\times d}(R)$.

\subsubsection*{Schwartz spaces and Heisenberg operators}
Let $\tliftS(F_\infty^\kappa)$ be the Schwartz space of $F_\infty^\kappa\cong\R^{[F:\Q]nd}$ and, for a finite
place $v$, $\tliftS(F_v^\kappa)$ the space of locally constant compactly supported functions. The
\emph{Schwartz--Bruhat space} $\tliftS(\A^\kappa)$ is the span of the functions $f\otimes\phi_f$,
$x\mapsto f(x_\infty)\phi_f(x_f)$, with $f\in\tliftS(F_\infty^\kappa)$ and $\phi_f$ locally constant with compact
support on $\A_f^\kappa$. For $f\in\tliftS(F_\infty^\kappa)$ and a family $\phi=(\phi_v)_v$ with
$\phi_v\in\tliftS(F_v^\kappa)$ and $\phi_v=\one_{\OO_v^\kappa}$ for almost all $v$, the \emph{pure tensor}
$f\otimes\phi$ is the function $x\mapsto f(x_\infty)\prod_v\phi_v(x_v)$; it lies in $\tliftS(\A^\kappa)$.

Let $h(u,u')=\sum_{i,j,k,l}u_{ij}\,(J_X)_{ik}(J_W)_{jl}\,\bar u'_{kl}$ for $u,u'\in M_{n\times d}(\tliftAE)$ and
\[
  \langle u,u'\rangle=-\Tr_{\tliftAE/\A}h(u,u').
\]
This is an alternating $\A$-bilinear form, and $(g,h)\cdot u=g\,u\,{}^th$ defines an action of
$\Uni(X)(\A)\times\Uni(W)(\A)$ preserving it. \emph{Darboux coordinates} are an $\A$-linear isomorphism
$e=(e_1,e_2)\colon M_{n\times d}(\tliftAE)\to\A^\kappa\times\A^\kappa$ with
$\langle u,u'\rangle=e_1(u)\cdot e_2(u')-e_1(u')\cdot e_2(u)$. The \emph{Heisenberg operators} are
\begin{equation}\label{eq:tlift:schr}
  (\rho(u)\Phi)(x)=\psi\bigl(x\cdot e_2(u)+\tfrac12\,e_1(u)\cdot e_2(u)\bigr)\,\Phi(x+e_1(u))\qquad(x\in\A^\kappa).
\end{equation}
An adele, a matrix of adeles or a group element $g$ is \emph{supported at} a finite place $v$ (resp.\ at
infinity) if it (resp.\ $g-1$) is killed by every adele whose $v$-component (resp.\ archimedean
component) vanishes. An operator $T$ on functions on $\A^\kappa$ \emph{acts at $v$} by $\phi\mapsto\phi_1$
($\phi,\phi_1$ functions on $F_v^\kappa$) if $T(\phi\otimes r)=\phi_1\otimes r$ for every function $r$ on
$\A^\kappa$ which is invariant under translations supported at $v$ and such that $\phi\otimes r$ is a
Schwartz--Bruhat function; here $(\phi\otimes r)(x)=\phi(x_v)r(x)$. Acting at infinity is defined in the
same way.

\begin{definition}[Good models]\label{pf:tlift:model}
A \emph{model} for $(X,W)$ is $M=(e,(\tliftvac_v)_v,\mu,D)$ with Darboux coordinates $e$, functions
$\tliftvac_v$ on $F_v^\kappa$ for the finite places $v$, a measure $\mu$ on $\Uni(X)(\A)$ and a set
$D\subset\Uni(X)(\A)$. It is \emph{good} if
\begin{enumerate}[label=(M\arabic*),nosep]
\item $e$ maps $M_{n\times d}(E)$ into $F^\kappa\times F^\kappa$;
\item for every $v$, $\tliftvac_v$ is a \emph{vacuum vector}: $\tliftvac_v\in\tliftS(F_v^\kappa)$, $\tliftvac_v\ne0$, and
  for some $\vartheta\in\OO_E$ with $\vartheta+\bar\vartheta=1$, every $\gamma$ in the group $\Lambda_v$ generated
  by the matrices $a\,x\,{}^ty$ ($a\in\A$ supported at $v$ with $a_v\in\OO_v$, $x\in L_X$, $y\in L_W$)
  satisfies: $\rho(\gamma)$ acts at $v$ by
  $\tliftvac_v\mapsto\psi\bigl(\tfrac12\Tr_{\tliftAE/\A}(\vartheta\,h(\gamma,\gamma))\bigr)\tliftvac_v$;
\item $\tliftvac_v=\one_{\OO_v^\kappa}$ for almost all $v$;
\item $\mu$ is a Haar measure and $D$ is a measurable fundamental domain for the left action of
  $\Uni(X)(F)$.
\end{enumerate}
\formal{GlobalTheory.IndefiniteRowSideFullHecke.WeilTensorModels.Model.IsGood}
\end{definition}

\subsubsection*{Weil representations}
The formal notion of a Weil representation with a splitting is the following characterisation of
Kudla's splittings \cite{Kudla1994,HarrisKudlaSweet1996} through the action of the Siegel parabolic
subgroup of a doubled group in the Schr\"odinger model. A linear form $\ell$ on $\tliftS(\A^{\kappa'})$ is
\emph{admissible} if for every locally constant compactly supported $\phi_f$ on $\A_f^{\kappa'}$ the form
$f\mapsto\ell(f\otimes\phi_f)$ is continuous on $\tliftS(F_\infty^{\kappa'})$. For $t\in\GL_m(\tliftAE)$ we put
$|\det t|_{\tliftAE}=|\Nm_{\tliftAE/\A}\det t|_\A$.

\begin{definition}[Weil representations]\label{pf:tlift:weil}
Let $M$ be a model and $\chi_X,\chi_W$ unitary Hecke characters. A homomorphism $\omega$ from
$\Uni(X)(\A)\times\Uni(W)(\A)$ to the linear endomorphisms of the space of functions on $\A^\kappa$ is a
\emph{Weil representation} for $(M,\psi,\chi_X,\chi_W)$ if:
\begin{enumerate}[label=(W\arabic*),nosep]
\item $\omega$ preserves $\tliftS(\A^\kappa)$ and is unitary on it for the $L^2$-norm of every Haar measure;
  $\ell\circ\omega(x)$ is admissible and $x\mapsto\ell(\omega(x)\Phi)$ is continuous for every admissible $\ell$
  and $\Phi\in\tliftS(\A^\kappa)$; and $\omega(g,h)\rho(u)\Phi=\rho(g\,u\,{}^th)\,\omega(g,h)\Phi$;
\item (splitting of $\Uni(W)$ by $\chi_X$) there are Darboux coordinates $e^{(2)}$ on
  $M_{n\times2d}(\tliftAE)$, for the form built from $J_X$ and $\mathrm{diag}(J_W,-J_W)$, with
  \[e^{(2)}(u_1\,|\,u_2)=\bigl((e_1(u_1),e_1(u_2)),(e_2(u_1),-e_2(u_2))\bigr),\]
  and a representation
  $\omega^{(2)}$ of $\Uni_{\mathrm{diag}(J_W,-J_W)}(\A)$ which satisfies (W1) for the action $u\mapsto u\,{}^tg$
  and the \emph{Siegel rule}: if $\ell$ is admissible and $\ell(\rho^{(2)}(z\,{}^t\tliftdia)\Phi)=\ell(\Phi)$
  for all $z\in M_{n\times d}(\tliftAE)$, where $\tliftdia={1_d\choose 1_d}$, then
  \[
    \ell(\omega^{(2)}(p)\Phi)=\chi_X(\det t)\,|\det t|_{\tliftAE}^{n/2}\,\ell(\Phi)\qquad\text{whenever }p\tliftdia=\tliftdia\,t;
  \]
  moreover $\omega^{(2)}(\mathrm{diag}(h,1))(\Phi_1\otimes\Phi_2)=\omega(1,h)\Phi_1\otimes\Phi_2$;
\item (splitting of $\Uni(X)$ by $\chi_W$) the same with the roles of $X$ and $W$ exchanged, using the
  coordinates $u\mapsto e({}^tu)$; and $\omega$ is continuous in the sense of (W1).
\end{enumerate}
A \emph{Weil pair} $(M,\chi_X,\chi_W,\omega)$ for $(X,W)$ consists of a good model $M$, unitary Hecke
characters with $\chi_X|_{\A^\times}=\eta^n$ and $\chi_W|_{\A^\times}=\eta^d$, and a Weil representation $\omega$ for
$(M,\psi,\chi_X,\chi_W)$.
\formal{GlobalTheory.WeilM.IsUWeilPair}
\end{definition}

For a Weil pair the \emph{theta series} and the \emph{theta lift} of a function $\varphi$ on
$\Uni(X)(\A)$ are
\begin{equation}\label{eq:tlift:theta}
  \theta(g,h;\Phi)=\sum_{\xi\in F^\kappa}(\omega(g,h)\Phi)(\xi),\qquad
  \tliftlift_M(\Phi,\varphi)(h)=\int_D\theta(g,h;\Phi)\,\overline{\varphi(g)}\,d\mu(g).
\end{equation}
We also use the matrix coefficients, for a Haar measure $\nu$ on $\A^\kappa$,
\begin{equation}\label{eq:tlift:coef}
  \tliftZ_\Phi(g)=\int_{\A^\kappa}(\omega(g,1)\Phi)(x)\,\overline{\Phi(x)}\,d\nu(x),\qquad
  \tliftM_\varphi(g)=\int_D\varphi(yg)\,\overline{\varphi(y)}\,d\mu(y).
\end{equation}

\subsubsection*{The polynomial Fock space}
For $u$ supported at infinity (an \emph{archimedean} $u$) and $f\in\tliftS(F_\infty^\kappa)$, let $D_uf$ be the
function such that $\frac{d}{dt}\big|_{t=0}\rho(tu)(f\otimes r)=D_uf\otimes r$ pointwise, for all
admissible complementary factors $r$. Let $\vartheta_{X,w},\vartheta_{W,w}$ be the Cartan involutions of $X$ and $W$.

\begin{definition}[Polynomial Fock space]\label{pf:tlift:pol}
$f_0\in\tliftS(F_\infty^\kappa)$ is a \emph{Gaussian} if there are signs $\varepsilon_w\in\{\pm1\}$ such that
$D_{u'}f_0=i\,D_uf_0$ for all archimedean $u,u'$ with
$\tau_w(u')=\varepsilon_w\,i\,\vartheta_{X,w}\,\tau_w(u)\,{}^t\vartheta_{W,w}$ for every $w$. The \emph{polynomial Fock
space} $\tliftPol=\tliftPol_M$ is the smallest subspace of $\tliftS(F_\infty^\kappa)$ containing the Gaussians
and stable under all $D_u$. A family $\phi=(\phi_v)_v$ with $\phi_v\in\tliftS(F_v^\kappa)$ for all $v$ and
$\phi_v=\tliftvac_v$ for almost all $v$ is called \emph{finite data}, and $\tliftPD=\tliftPD(M)$ is the span of the
pure tensors $f\otimes\phi$ with $f\in\tliftPol$ and finite data $\phi$.
\formal{GlobalTheory.ThK.PolData}
\end{definition}

\subsubsection*{The doubled group}
For a space $Y$ of dimension $m$ let $G_2(Y)=\Uni_{\mathrm{diag}(J_Y,-J_Y)}$, the unitary group of
$Y\oplus(-Y)$, and $\tliftdia={1_m\choose1_m}$. The \emph{Siegel parabolic} $P$ consists of the $p$ with
$p\tliftdia=\tliftdia\,t_p$ for some $t_p\in\GL_m$, the \emph{Levi component} of $p$. Put $\iota(g,g')=\mathrm{diag}(g,g')$
for $g,g'\in\Uni(Y)(\A)$; thus $\iota(g,g)\in P(\A)$ with $t_{\iota(g,g)}=g$. Let $\widehat L_Y\subset\tliftAE^m$ be
the closure of $L_Y$.

\begin{definition}[Standard maximal compact subgroup, heights]\label{pf:tlift:stab2}
The \emph{standard maximal compact subgroup} $\tliftK\subset G_2(Y)(\A)$ consists of the $k$ such that
$\tau_w(k)$ commutes with $\mathrm{diag}(\vartheta_w,-\vartheta_w)$ for every infinite $w$, and both $k$ and $k^{-1}$
map $\widehat L_Y\oplus\widehat L_Y$ into itself. A \emph{height} for $\tliftK$ is a function $\tliftht$ on
$G_2(Y)(\A)$ such that $G_2(Y)(\A)=P(\A)\tliftK$ and $\tliftht(pk)=|\det t_p|_{\tliftAE}$ for $p\in P(\A)$,
$k\in\tliftK$. \formal{GlobalTheory.ThK.Dbl.IsHgt}
\end{definition}

\subsubsection*{Doubling data, sections and Eisenstein series}
Let $\kappa_2=\kappa\sqcup\kappa$. For functions $\Phi_1,\Phi_2$ on $\A^\kappa$ let $\Phi_1\boxtimes\bar\Phi_2$ be the
function $(x,y)\mapsto\Phi_1(x)\overline{\Phi_2(y)}$ on $\A^{\kappa_2}$.

\begin{definition}[Doubling data]\label{pf:tlift:dbl}
Let $A$ be a space of dimension $m'$ and $Y$ a space of dimension $m$, one hermitian and one
skew-hermitian, and $\chi$ a unitary Hecke character. \emph{Doubling data} for $(A,Y,\chi)$ consist of
Darboux coordinates $e^{(2)}$ on $M_{m'\times2m}(\tliftAE)$ (for $J_A$ and $\mathrm{diag}(J_Y,-J_Y)$) mapping
rational matrices to rational vectors, a Weil representation $\omega_2$ of $G_2(Y)(\A)$ on functions on
$\A^{\kappa_2}$ ($\kappa_2$ the index set of these coordinates) satisfying (W1) and the Siegel rule of (W2) with
the character $\chi$ and the Lagrangian $\tliftdia$, a representation $S$ of $\Uni(A)(\A)$ satisfying (W1) for
the linear action $u\mapsto gu$, and an admissible linear form $\ell$ on $\tliftS(\A^{\kappa_2})$ which is
invariant under $S$ and under the Heisenberg operators $\rho_2(z\,{}^t\tliftdia)$, and such that
$\sum_{y\in M_{m'\times m}(E)}\ell\bigl(\rho_2(y\,|\,{-y})\Psi\bigr)=\sum_{\xi\in F^{\kappa_2}}\Psi(\xi)$ for
$\Psi\in\tliftS(\A^{\kappa_2})$. We write $\theta_2(x,h;\Psi)=\sum_\xi(\omega_2(x)S(h)\Psi)(\xi)$.
\formal{GlobalTheory.ThK.Dbl.IsDbl}
\end{definition}

\begin{definition}[Sections, Eisenstein series, continuation]\label{pf:tlift:eiscont}
A continuous function $f$ on $G_2(Y)(\A)$ is a \emph{section} of exponent $s_1$ for $(\tliftK,\chi)$ if
$f(px)=\chi(\det t_p)\,|\det t_p|_{\tliftAE}^{s_1+m/2}f(x)$ for $p\in P(\A)$ and the right translates
$f(\,\cdot\,k)$, $k\in\tliftK$, span a finite-dimensional space. For a height $\tliftht$ and $s\in\C$ put
$f_s=\tliftht^{\,s-s_1}f$, and let
\begin{equation}\label{eq:tlift:eis}
  E(x,s;f)=\sum_{\gamma\in P(F)\backslash G_2(Y)(F)}f_s(\gamma x).
\end{equation}
A \emph{continuation} of $E(\cdot,\cdot;f)$ on an open set $\mathcal U\subset\C$ is a function $\tliftE(s,x)$ on
$\mathcal U\times G_2(Y)(\A)$ such that: $\mathcal U$ is connected and contains a punctured neighbourhood of
every point of $\C$; there is $t$ with $\{\Re s>t\}\subset\mathcal U$ and \eqref{eq:tlift:eis} converging to
$\tliftE(s,x)$ there; $\tliftE$ is holomorphic in $s$ and jointly continuous; and $\tliftE$ is of moderate
growth in $x$, locally uniformly in $s$. \formal{GlobalTheory.ThK.Dbl.IsEisCont}
\end{definition}

\subsubsection*{The tool package}
The following results are proved in the development. For those which are classical we give the
references and do not reproduce the formal proofs.

\begin{theorem}[Weil representations]\label{pf:tlift:weilex}
Let the base data be good, $X$ hermitian and $W$ skew-hermitian (of any signatures), $M$ a good model,
and $\chi_X,\chi_W$ unitary Hecke characters with $\chi_X|_{\A^\times}=\eta^n$ and $\chi_W|_{\A^\times}=\eta^d$. Then
there is a Weil representation for $(M,\psi,\chi_X,\chi_W)$.
\formallabel{PositivityTheory.ThKW.S.weilEx_k}\formallabel{PositivityTheory.ThKW.S.exists_isWeil_k}
\end{theorem}

This is the existence of the Weil representation of the dual pair $(\Uni(X),\Uni(W))$ over a number field
\cite{Weil1964}, with the splitting of the metaplectic cover over the unitary groups attached to the
characters $\chi_X,\chi_W$ \cite{Kudla1994,HarrisKudlaSweet1996}. The published form, the differences with the
formal statement and the formal proof are given in \cref{pf:tlift:classical}.

\begin{proposition}[Good models; Weil pairs]\label{pf:tlift:pairex}
Under the assumptions of \cref{pf:tlift:weilex} (without $M$) a good model exists. Consequently Weil pairs
$(M,\chi_X,\chi_W,\omega)$ exist for every good model $M$ and all characters as in \cref{pf:tlift:weilex}.
\formallabel{PositivityTheory.ThKG.Tow.modelEx_k}\formallabel{PositivityTheory.ThKG.Tow.pairEx_of_k}
\end{proposition}

\begin{proof}
Rational Darboux coordinates exist (\lref{Tw.RatDarboux.exists_rational_darboux_of_isGood_k}), a Haar
measure on $\Uni(X)(\A)$ with a fundamental domain exists (\lref{PositivityTheory.ThKF.unimodularity_holds}),
vacuum vectors exist at every place (\lref{Tw.VacAny.exists_isVac_any_darboux_k}), and
$\one_{\OO_v^\kappa}$ is a vacuum vector for almost all $v$ (\lref{Tw.VacIndic.eventually_isVac_indicator_k}).
Take $\tliftvac_v=\one_{\OO_v^\kappa}$ whenever it is a vacuum vector, and some vacuum vector otherwise; then
(M1)--(M4) hold. The second assertion combines this with \cref{pf:tlift:weilex}.
\end{proof}

\begin{proposition}[Compact quotients]\label{pf:tlift:cpt}
Let $Y$ be a definite space. There is a compact $\Omega\subset\Uni(Y)(\A)$ with
$\Uni(Y)(\A)=\Uni(Y)(F)\,\Omega$; $\Uni(Y)(F)$ is countable; Haar measures on $\Uni(Y)(\A)$ are right invariant;
fundamental domains for $\Uni(Y)(F)$ exist and have finite Haar measure. In particular continuous left
$\Uni(Y)(F)$-invariant functions on $\Uni(Y)(\A)$ are bounded.
\formallabel{GlobalTheory.ThK.Dbl.cptAll_of_k}
\end{proposition}

\begin{proof}
This is reduction theory \cite{Borel1969,PlatonovRapinchuk1994}: the formal proof combines the topology of
the adelic group (\lref{PositivityTheory.ThKEst.grpTop_of}), Siegel sets for unitary groups
(\lref{UpperTheory.Lit.siegelSetUnitary_of_core}, \lref{PositivityTheory.ThKEst.siegelAll_of}),
finiteness of the volume (\lref{PositivityTheory.ThKEst.finVolU}) and unimodularity
(\lref{PositivityTheory.ThKF.unimodularity_holds}). Since $Y$ is definite, the Witt index of $Y$, which is at most
$\min(r,s)$, vanishes (\lref{PositivityTheory.ThKEst.epsOf_of_k}); hence the torus part of the Siegel set is trivial
and the Siegel set is the compact set $\Omega$.
\end{proof}

\begin{lemma}[Anisotropy]\label{pf:tlift:aniso}
If $Y$ is definite, then $h_{J_Y}(x,x)=0$ implies $x=0$ for $x\in E^m$.
\formallabel{GlobalTheory.ThK.Dbl.aniso_of_k}
\end{lemma}

\begin{proof}
Apply $\tau_w$ for one infinite place $w$: $\tau_w(J_Y)$ (resp.\ $i\tau_w(J_Y)$) has signature $(m,0)$ or $(0,m)$,
so the hermitian form it defines is definite and $\tau_w(x)=0$.
\end{proof}

\begin{proposition}[Theta series]\label{pf:tlift:basic}
Let $(M,\chi_X,\chi_W,\omega)$ be a Weil pair. For $\Phi\in\tliftS(\A^\kappa)$, $\theta(g,h;\Phi)$ is continuous in
$(g,h)$, invariant under $\Uni(X)(F)\times\Uni(W)(F)$ on the left, and of moderate growth:
$|\theta(g,h;\Phi)|\le C(\mathrm{ht}(g)\mathrm{ht}(h))^k$ for some $C,k$, where $\mathrm{ht}$ is the height on
$\GL_m(\tliftAE)$ defined as the product over all places of local heights; for fixed finite part,
$\theta$ depends continuously on the archimedean factor. Moreover $\tliftPD(M)\subset\tliftS(\A^\kappa)$.
\formallabel{PositivityTheory.ThKF0.thetaBasic_k}\formallabel{GlobalTheory.ThK.Dbl.polSch_k}
\end{proposition}

\begin{proof}
The theta distribution $\Psi\mapsto\sum_{\xi\in F^\kappa}\Psi(\xi)$ is an admissible linear form
(\lref{GlobalTheory.T.theta_distribution_nonzero_admissible}), so continuity in $(g,h)$ and in the
archimedean factor follow from (W1). Invariance under rational points is
\lref{PositivityTheory.ThKF0.thetaSum_rational_pair_invariant}, and the growth estimate
(\lref{PositivityTheory.ThKF0.norm_theta_le_mul_ht_pow_of_isPair}) follows from the uniform decay
$|\omega(g,h)\Phi(x)|(1+\|x_\infty\|)^N\le C(\mathrm{ht}(g)\mathrm{ht}(h))^k$
(\lref{PositivityTheory.ThKF0.weil_schwartz_decay_bound}); the formal proof is
\lref{PositivityTheory.ThKF0.thetaBasic_of_k}. The last assertion holds because pure tensors with
finite data are Schwartz--Bruhat functions, by (M3).
\end{proof}

\begin{proposition}[Local actions]\label{pf:tlift:acts}
Let $(M,\chi_X,\chi_W,\omega)$ be a Weil pair and $v$ a finite place. For $g\in\Uni(X)(\A)$ and $h\in\Uni(W)(\A)$
supported at $v$ and $\phi\in\tliftS(F_v^\kappa)$ there is a unique $\phi_1\in\tliftS(F_v^\kappa)$ such that
$\omega(g,h)$ acts at $v$ by $\phi\mapsto\phi_1$; the same holds for $\rho(u)$, $u$ supported at $v$; and
every $\phi\in\tliftS(F_v^\kappa)$ is fixed in this sense by $\omega(x)$ for all $x$ supported at $v$ in some
neighbourhood of $1$.
\formallabel{PositivityTheory.ThKF0.actsEx_k}
\end{proposition}

\begin{proof}
Fix a factor $r_0$ and a point $x_0$ with $r_0(x_0)=1$; the local operator is
$\phi\mapsto(\omega(g,h)(\phi\otimes r_0))(\cdot,x_0)$ (\lref{PositivityTheory.ThKH.Locrep.locRepEx_k}); it
does not depend on the choices because $\omega(g,h)$ commutes with the Heisenberg operators supported
away from $v$, by (W1). Smoothness follows from the continuity in (W1), since the compact open subgroups
supported at $v$ form a basis of neighbourhoods of $1$ in the group of elements supported at $v$.
\end{proof}

\begin{proposition}[Doubling]\label{pf:tlift:dblx}
For a Weil pair $(M,\chi_X,\chi_W,\omega)$ there are doubling data $(e^{(2)},\omega_2,S,\ell)$ for
$(W,X,\chi_W)$ such that for $\Phi_1,\Phi_2\in\tliftS(\A^\kappa)$ one has $\Phi_1\boxtimes\bar\Phi_2\in\tliftS(\A^{\kappa_2})$
and, for $g,g'\in\Uni(X)(\A)$ and $h\in\Uni(W)(\A)$,
\begin{align}
  \chi_W(\det g')^{-1}\,\theta_2\bigl(\iota(g,g'),h;\Phi_1\boxtimes\bar\Phi_2\bigr)
    &=\theta(g,h;\Phi_1)\,\overline{\theta(g',h;\Phi_2)},\label{eq:tlift:dbltheta}\\
  \chi_W(\det g')^{-1}\,\omega_2(\iota(g,g'))(\Phi_1\boxtimes\bar\Phi_2)
    &=\omega(g,1)\Phi_1\boxtimes\overline{\omega(g',1)\Phi_2};\label{eq:tlift:dblact}
\end{align}
the coordinates $e^{(2)}$ restrict to $e$ on each copy of $M_{n\times d}$; and for every Haar measure $\nu$ on
$\A^\kappa$ there is $c'>0$ with $\ell(\Psi_1\boxtimes\bar\Psi_2)=c'\int\Psi_1\bar\Psi_2\,d\nu$ for
$\Psi_i\in\tliftS(\A^\kappa)$.
\formallabel{GlobalTheory.ThK.Dbl.dblX_k}
\end{proposition}

\begin{proof}
The data are constructed from the splitting (W3) of $\Uni(X)$ in \lref{GlobalTheory.ThK.Dbl.dbl_main},
which also gives \eqref{eq:tlift:dblact} and the identity
$\chi_W(\det g')^{-1}\theta_2(\iota(g,g'),1;\Psi_1\boxtimes\bar\Psi_2)=\theta(g,1;\Psi_1)\overline{\theta(g',1;\Psi_2)}$.
Applying the latter to $\Psi_i=\omega(1,h)\Phi_i$ and using that $S(h)(\Phi_1\boxtimes\bar\Phi_2)=
\omega(1,h)\Phi_1\boxtimes\overline{\omega(1,h)\Phi_2}$ and $\omega(g,1)\omega(1,h)=\omega(g,h)$ gives
\eqref{eq:tlift:dbltheta}. The last assertion is the uniqueness of Haar measure
(\lref{GlobalTheory.ThK.Dbl.lD_haar}).
\end{proof}

\begin{lemma}[Section law]\label{pf:tlift:law}
In the situation of \cref{pf:tlift:dblx}, for $\Psi\in\tliftS(\A^{\kappa_2})$ the function
$f_\Psi(x)=\ell(\omega_2(x)\Psi)$ on $G_2(X)(\A)$ satisfies
$f_\Psi(px)=\chi_W(\det t_p)\,|\det t_p|_{\tliftAE}^{d/2}f_\Psi(x)$ for $p\in P(\A)$.
\formallabel{GlobalTheory.ThK.Dbl.Bad.law_k}
\end{lemma}

\begin{proof}
$f_\Psi(px)=\ell(\omega_2(p)\,\omega_2(x)\Psi)$, and $\omega_2(x)\Psi\in\tliftS$; the Siegel rule, applied to $\ell$
(which is admissible and invariant under the $\rho_2(z\,{}^t\tliftdia)$), gives the factor
$\chi_W(\det t_p)|\det t_p|_{\tliftAE}^{d/2}$, $d$ being the dimension of $W$.
\end{proof}

\begin{proposition}[Unfolding]\label{pf:tlift:unfold}
Let $Y$ be a definite space, $\chi$ a unitary Hecke character and $\iota$ as above. Let $f$ be continuous on
$G_2(Y)(\A)$ with $f(px)=\chi(\det t_p)|\det t_p|^z_{\tliftAE}f(x)$ for some $z\in\C$, such that
$\tliftE(x)=\sum_{\gamma\in P(F)\backslash G_2(Y)(F)}f(\gamma x)$ converges for every $x$ and $g\mapsto f(\iota(1,g))$ is
integrable for a Haar measure $\mu$ with fundamental domain $D$. Then for every bounded continuous left
$\Uni(Y)(F)$-invariant $\varphi$ and every $h$,
\begin{equation}\label{eq:tlift:unfold}
  \int_D\tliftE(\iota(h,h'))\,\chi(\det h')^{-1}\varphi(h')\,d\mu(h')
  =\int_{\Uni(Y)(\A)}f(\iota(1,g))\,\chi(\det g)^{-1}\varphi(hg)\,d\mu(g).
\end{equation}
\formallabel{GlobalTheory.ThK.Dbl.unfold0_k}
\end{proposition}

\begin{proof}
Since $Y$ is anisotropic (\cref{pf:tlift:aniso}), $\gamma\mapsto P(F)\iota(1,\gamma)$ is a bijection from $\Uni(Y)(F)$
onto $P(F)\backslash G_2(Y)(F)$ (\lref{GlobalTheory.ThK.Dbl.Law.hasSum}), so
$\tliftE(\iota(h,h'))=\sum_{\gamma}f(\iota(h,\gamma h'))$. As $\iota(h,h)\in P(\A)$ has Levi component $h$ and
$|\det h|_{\tliftAE}=1$, the law gives $f(\iota(h,g))=\chi(\det h)f(\iota(1,h^{-1}g))$
(\lref{GlobalTheory.ThK.Dbl.Law.diag}). Since $\chi(\det\gamma)=1$ and $\varphi(\gamma h')=\varphi(h')$ for rational
$\gamma$, the integrand on the left of \eqref{eq:tlift:unfold} equals $\sum_\gamma k(\gamma h')$ with
$k(g)=\chi(\det h)\chi(\det g)^{-1}f(\iota(1,h^{-1}g))\varphi(g)$, an integrable function. Unfolding the integral
over the fundamental domain (\lref{GlobalTheory.ThK.Dbl.An.unfold}) gives $\int k\,d\mu$, and the
substitution $g\mapsto hg$ gives the right side.
\end{proof}

\begin{theorem}[Siegel--Weil formula; formal form]\label{pf:tlift:swstd}
Let the base data be good, $A$ a definite space of dimension $m'\ge1$ and $Y$ a space of dimension
$m\ge1$, one of them hermitian and the other skew-hermitian, $\chi$ a unitary Hecke character with
$\chi|_{\A^\times}=\eta^{m'}$, $(e^{(2)},\omega_2,S,\ell)$ doubling data for $(A,Y,\chi)$, $\tliftKfr$ the maximal compact
subgroup of $G_2(Y)(\A)$ attached to a frame of $Y$ (\cref{pf:tlift:frameconj}) with a height $\tliftht_0$,
and $\mu_A$ a Haar measure on $\Uni(A)(\A)$ with fundamental domain $D_A$. Let
$\Psi\in\tliftS(\A^{\kappa_2})$ be $\tliftKfr$-finite under $\omega_2$, $f(x)=\ell(\omega_2(x)\Psi)$ and $s_0=\frac{m'-m}2$.
Then for every $x$ the series $E(x,s;f)$ of \eqref{eq:tlift:eis} (with $\tliftht_0$ and $s_1=s_0$) converges
for $\Re s$ large and extends to a meromorphic function of $s\in\C$ which is analytic at $s_0$ with value
\[
  \kappa\cdot\frac1{\mu_A(D_A)}\int_{D_A}\theta_2(x,h;\Psi)\,d\mu_A(h),\qquad\kappa=1\text{ if }m<m',\quad\kappa=2\text{ if }m\ge m'.
\]
\formallabel{PositivityTheory.SiegelWeilFormula.Fin.siegelWeilStd_k}
\end{theorem}

This is the Siegel--Weil formula for the dual pair $(\Uni(A),G_2(Y))$ with $A$ anisotropic
\cite{Weil1965,KudlaRallis1988,Ichino2007}; it is used below only with $m'=d\ge2m=2n$, where $\kappa=1$. The published form,
the differences with the formal statement and the structure of the formal proof (more than $15{,}000$ lines) are
given in \cref{pf:tlift:classical}.

\begin{theorem}[Continuation of Siegel Eisenstein series; formal form]\label{pf:tlift:eisbl}
Let $Y$ be a hermitian or skew-hermitian space of dimension $m\ge1$, $\chi$ a unitary Hecke character with
$\chi|_{\A^\times}=\eta^c$ for some $c$, $\tliftKfr$ the maximal compact subgroup attached to a frame with a height,
and $f$ a section of exponent $s_1$ for $(\tliftKfr,\chi)$. Then $E(x,s;f)$ converges for $\Re s>t$ (some $t$) to a
function holomorphic in $s$, and for every $s_0\in\C$ there are $k\in\N$ and a neighbourhood $\mathcal V$ of $s_0$
such that $(s-s_0)^kE(\cdot,s;f)$ extends from $\mathcal V\cap\{\Re s>t\}$ to a weakly holomorphic function on
$\mathcal V$ with values in the functions of moderate growth.
\formallabel{PositivityTheory.RgEis.Hd.eisBL_k}
\end{theorem}

This is the meromorphic continuation of Eisenstein series \cite{Langlands1976}, \cite[IV]{MoeglinWaldspurger1995}
in the case of the Siegel parabolic of a quasi-split unitary group of rank $m$; the formal proof is
long and we do not reproduce it.

\begin{lemma}[Frames]\label{pf:tlift:frameconj}
Let $Y$ be a hermitian or skew-hermitian space of dimension $m$. A \emph{frame} of $Y$ is a matrix
$b\in\GL_{2m}(E)$ with ${}^tb\,\mathrm{diag}(J_Y,-J_Y)\,\bar b={0\ \ 1\choose\epsilon\ \ 0}$ for a sign $\epsilon$ and whose
first $m$ columns span the column space of $\tliftdia$; the \emph{maximal compact subgroup attached to $b$}
consists of the $k\in G_2(Y)(\A)$ such that $b^{-1}kb$ is unitary at every infinite place and $b^{-1}k^{\pm1}b$
have integral entries at every finite place. The standard maximal compact subgroup $\tliftK$ is conjugate
to the subgroup $\tliftKfr$ attached to some frame: $\tliftK=p\tliftKfr p^{-1}$ for some $p\in P(\A)$.
\formallabel{GlobalTheory.ThK.Dbl.frameConj_k}
\end{lemma}

\begin{proof}
The formal proof constructs the frame and the conjugating element explicitly, from bases of the
archimedean spaces adapted to the Cartan involutions (\lref{GlobalTheory.ThK.Dbl.Frame.exists_arch}),
adelic bases of the lattice $\widehat L_Y$ (\lref{GlobalTheory.ThK.Dbl.Frame.exists_ad_k}), and an element
$\vartheta\in\OO_E$ with $\vartheta+\bar\vartheta=1$, which exists because $E/F$ is unramified at all finite places
(\lref{GlobalTheory.exists_int_add_conj_eq_one_of_base_good_k}).
\end{proof}

\begin{proposition}[Siegel--Weil formula in the form used]\label{pf:tlift:sw}
Let $(M,\chi_X,\chi_W,\omega)$ be a Weil pair with $n,d\ge1$ and $W$ definite, $(e^{(2)},\omega_2,S,\ell)$ doubling
data as in \cref{pf:tlift:dblx}, $\tliftK$ the standard maximal compact subgroup of $G_2(X)(\A)$ with a height
$\tliftht$, $\mu_W$ a Haar measure on $\Uni(W)(\A)$ with fundamental domain $D_W$, and
$\Phi_1,\Phi_2\in\tliftPD(M)$. Put $s_0=\frac{d-n}2$ and $f(x)=\ell(\omega_2(x)(\Phi_1\boxtimes\bar\Phi_2))$. Then $f$
is a section of exponent $s_0$ for $(\tliftK,\chi_W)$, and $E(\cdot,\cdot;f)$ has a continuation $\tliftE$ on an
open set $\mathcal U\ni s_0$ with
\begin{equation}\label{eq:tlift:sw}
  \tliftE(s_0,x)=\kappa\cdot\frac1{\mu_W(D_W)}\int_{D_W}\theta_2\bigl(x,h;\Phi_1\boxtimes\bar\Phi_2\bigr)\,d\mu_W(h),
  \qquad\kappa=\begin{cases}1,&n<d,\\2,&n\ge d.\end{cases}
\end{equation}
\formallabel{GlobalTheory.ThK.Dbl.swFullX_of_k}
\end{proposition}

\begin{proof}
\emph{Section property} (\lref{GlobalTheory.ThK.Dbl.boxSecX_of_k}). Continuity follows from (W1), and
the law from \cref{pf:tlift:law} since $s_0+\frac n2=\frac d2$. For the $\tliftK$-finiteness it suffices, by
bilinearity, to treat $\Phi_i=f_i\otimes\phi^{(i)}$ with $f_i\in\tliftPol$. At the finite places
$\Phi_1\boxtimes\bar\Phi_2$ is fixed by an open subgroup of the compact group $\tliftK$
(\lref{GlobalTheory.ThK.Dbl.finLevel_k}, \cref{pf:tlift:stab2cpt}); at the infinite places $\tliftK$
preserves the span of the $D_{u_1}\cdots D_{u_j}f_0$ with $f_0$ Gaussian and $j$ bounded, which is
finite-dimensional, and $\tliftPol$ is contained in the union of these spaces
(\lref{GlobalTheory.ThK.Dbl.Box.box_kfin_k}; compare \cref{pf:tlift:polfin} below).

\emph{Continuation.} By \cref{pf:tlift:frameconj}, $\tliftK=p\tliftKfr p^{-1}$ with $\tliftKfr$ attached to a frame,
$p\in P(\A)$. The function $x\mapsto f(xp^{-1})$ is a section of exponent $s_0$ for $(\tliftKfr,\chi_W)$, and
$\tliftht_0(x)=|\det t_p|_{\tliftAE}\,\tliftht(xp^{-1})$ is a height for $\tliftKfr$. Applying
\cref{pf:tlift:eisbl} to these and translating back by $x\mapsto xp$ (which multiplies the terms of
\eqref{eq:tlift:eis} by the constant $|\det t_p|_{\tliftAE}^{s_0-s}$) gives a continuation $\tliftE$ of
$E(\cdot,\cdot;f)$ on an open set $\mathcal U$, such that for every $s_1\in\C$ some $(s-s_1)^k\tliftE(s,\cdot)$ is
holomorphic near $s_1$ (\lref{GlobalTheory.ThK.Dbl.eisContinues_of_k}).

\emph{Value at $s_0$.} Translating \cref{pf:tlift:swstd} from $\tliftKfr$ to $\tliftK$ in the same way
(\lref{GlobalTheory.ThK.Dbl.Box.sw_transl_k}) and applying it with $(A,Y)=(W,X)$ and
$\Psi=\Phi_1\boxtimes\bar\Phi_2$, which is $\tliftK$-finite by the first step, shows that for each $x$ the function
$s\mapsto E(x,s;f)$ is the restriction of a meromorphic function analytic at $s_0$ with value the right
side of \eqref{eq:tlift:sw} (\lref{GlobalTheory.ThK.Dbl.siegelWeilAniso_of_k}). This meromorphic
function coincides with $\tliftE(\cdot,x)$ on a half plane, hence on $\mathcal U\setminus\{s_0\}$ near $s_0$; so the
singularity of $\tliftE$ at $s_0$ (of order at most $k$) is removable, and defining $\tliftE(s_0,\cdot)$ by
\eqref{eq:tlift:sw} gives a continuation on $\mathcal U\cup\{s_0\}$; the required holomorphy, continuity and
growth near $s_0$ are deduced from those of $(s-s_0)^k\tliftE(s,\cdot)$
(\lref{GlobalTheory.ThK.Dbl.IsEisCont.update}).
\end{proof}

\begin{proposition}[Split places]\label{pf:tlift:split}
Let $(M,\chi_X,\chi_W,\omega)$ be a Weil pair, $v$ a split place and $t\colon E\to F_v$. There is a linear
bijection $\iota_t\colon\tliftS(F_v^\kappa)\to\tliftS(M_{n\times d}(F_v))$ intertwining the Heisenberg operators
supported at $v$ with translations and multiplications by characters (a \emph{linear model}); and for
every linear model $\iota_t$, if $g\in\Uni(X)(\A)$ and $h\in\Uni(W)(\A)$ are supported at $v$ and $\omega(g,h)$ acts at
$v$ by $\phi\mapsto\phi_1$, then
\[
  \iota_t\phi_1=\omega^{\chi_{W,t}^{-1},\chi_{X,t}}\bigl(\pi_t(g),\pi_t(h)\bigr)\,\iota_t\phi
\]
in the notation of \eqref{eq:tlift:omA}.
\formallabel{UpperTheory.TP.splitFree_k}
\end{proposition}

\begin{proof}
At a split place $\Uni(X)(F_v)\cong\GL_n(F_v)$ and $\Uni(W)(F_v)\cong\GL_d(F_v)$ via $\pi_t$, and
$M_{n\times d}(E\otimes_FF_v)\cong M_{n\times d}(F_v)\times M_{n\times d}(F_v)$ is a decomposition into two Lagrangians
stable under both groups. A linear model is an intertwiner from the Schr\"odinger model to the model
attached to this polarisation; it exists by \lref{UpperTheory.TP.exists_linW_k}. The formula is
\lref{UpperTheory.TP.split_formula_k}, deduced from the description of the action of $\omega$ in the
linear models (\lref{PositivityTheory.KCo.dictX_proof_k}). This is the classical description of the
Weil representation of a type II dual pair \cite{MoeglinVignerasWaldspurger1987,Minguez2008}.
\end{proof}

\begin{proposition}[Fixed vacuum]\label{pf:tlift:vacfix}
Let $(M,\chi_X,\chi_W,\omega)$ be a Weil pair, $(e^{(2)},\omega_2,S,\ell)$ doubling data as in \cref{pf:tlift:dblx},
$\tliftK$ the standard maximal compact subgroup of $G_2(X)(\A)$, and $v$ a finite place such that $\chi_W$ is
unramified at every place of $E$ above $v$. If $k\in\tliftK$ and $k-1$ is supported at $v$, then $\omega_2(k)$
acts at $v$ trivially on $\tliftvac_v\boxtimes\overline{\tliftvac_v}$.
\formallabel{GlobalTheory.ThK.Dbl.vacFix2_k}
\end{proposition}

\begin{proof}
Since $k^{\pm1}$ preserve the doubled lattice, $\omega_2(k)$ maps $\tliftvac_v\boxtimes\overline{\tliftvac_v}$ to a
multiple $c\,\tliftvac_v\boxtimes\overline{\tliftvac_v}$: by (M2) and the Heisenberg relation in (W1), both are
eigenvectors with the same eigencharacter for the Heisenberg operators of the doubled lattice, and such
eigenvectors are unique up to scalars (\lref{PositivityTheory.ThKG.Loc.Cone.Vf.dbl_eigen_k}). The formal
proof then shows $c=1$ by decomposing $k$, using its first block column, into an element of the Siegel
parabolic with integral entries and integral elements of an opposite unipotent subgroup
(\cref{pf:tlift:beta}); the factors are evaluated with the Siegel rule, in
which $\chi_W$ is unramified at $v$ and the Levi component is a unit.
\end{proof}

\begin{theorem}[The tool package]\label{pf:tlift:toolsthm}
\Cref{pf:tlift:pairex,pf:tlift:cpt,pf:tlift:basic,pf:tlift:acts,pf:tlift:dblx,pf:tlift:unfold,pf:tlift:sw,%
pf:tlift:split,pf:tlift:vacfix} hold, as well as \cref{pf:tlift:stab2cpt} and the Lie algebra actions of
\cref{pf:tlift:fock}.
\formallabel{UpperTheory.TP.tools_holds_k}
\end{theorem}

\begin{proof}
This collects the results of this subsection, \cref{pf:tlift:stab2cpt} below, and the two statements on
the infinitesimal action of the unitary groups on the Fock model, \lref{PositivityTheory.RL.lieFirst} and
\lref{PositivityTheory.ThKR.lieSecond2} (recalled in \cref{pf:tlift:fock}).
\end{proof}


% ---------------------------------------------------------------------------
\subsection{The Weil representation and the Siegel--Weil formula}\label{pf:tlift:classical}

This subsection contains the formal proofs of \cref{pf:tlift:weilex,pf:tlift:swstd}. Both are forms of classical
theorems. For each we state the published form, describe how the formal statement differs from it, and convert
the parts of the formal proof which are specific to the development. The classical cores (the construction of the
Weil representation in a lattice model, and the comparison of Fourier coefficients of Eisenstein series and theta
integrals) are cited with one-sentence paraphrases.

\subsubsection*{The Weil representation: published form and formal form}
\begin{littheorem}[{\cite{Weil1964,Kudla1994}; see also \cite{HarrisKudlaSweet1996}}]\label{pf:tlift:litweil}
Let $V$ be a nondegenerate hermitian space of dimension $m$ and $U$ a nondegenerate skew-hermitian space of
dimension $k$ over $E$, let $\psi$ be a nontrivial character of $\A/F$, and let $\chi_V,\chi_U$ be characters of
$\tliftAE^\times/E^\times$ with $\chi_V|_{\A^\times}=\eta^m$ and $\chi_U|_{\A^\times}=\eta^k$. The Weil representation of the
metaplectic cover of $\mathrm{Sp}(\mathrm{Res}_{E/F}(V\otimes_EU))(\A)$, realised on the Schwartz--Bruhat space of a
Lagrangian, restricts, through the splitting of the cover over $\Uni(V)(\A)\times\Uni(U)(\A)$ determined by
$(\chi_V,\chi_U)$, to a smooth unitary representation $\omega$ of $\Uni(V)(\A)\times\Uni(U)(\A)$ which satisfies the
Heisenberg relation $\omega(g)\rho(u)\omega(g)^{-1}=\rho(g u)$. If $U=U_1\oplus U_2$ is a complete polarisation and
$P\subset\Uni(U)$ is the stabiliser of $U_2$, with Levi elements $m(a)$, $a\in\GL(U_2)$, then in the model on the
Schwartz--Bruhat space of $(V\otimes_EU_2)(\A)$ the linear form $\phi\mapsto\phi(0)$ is invariant under the unipotent
radical of $P$ and satisfies $(\omega(m(a))\phi)(0)=\chi_V(\det a)\,|\det a|_{\tliftAE}^{m/2}\,\phi(0)$.
\end{littheorem}

\begin{remark}[The formal statement]\label{pf:tlift:remweil}
\Cref{pf:tlift:weilex} differs from \cref{pf:tlift:litweil} as follows.
\begin{enumerate}[label=(\alph*),nosep]
\item The formal notion of a Weil representation (\cref{pf:tlift:weil}) is axiomatic: no metaplectic group
  appears, every representation on functions on $\A^\kappa$ satisfying (W1)--(W3) qualifies, and no uniqueness is
  asserted.
\item The splitting characters enter only through the Siegel rules in (W2) and (W3): the representation of the
  doubled group is part of the data, and the Siegel rule prescribes the action of the Siegel parabolic on every
  admissible linear form which is invariant under the Heisenberg translations along the diagonal Lagrangian, as
  in the last assertion of \cref{pf:tlift:litweil} (the evaluation at $0$ being one such form).
\item The coordinates are arbitrary rational Darboux coordinates (those of a good model); the representation acts
  on functions on $\A^\kappa$, not on the Schwartz space of a Lagrangian adapted to the dual pair.
\item The hypotheses are those of good base data; the formal construction uses that $E/F$ is unramified at every
  finite place and that $\psi$ is standard.
\end{enumerate}
\end{remark}

The formal proof passes through a second realisation, the \emph{theta model}, in which the representation acts on
functions on the full symplectic space which are quasi-periodic under the rational points.

\begin{definition}[Theta functions; Weil data in the theta model]\label{pf:tlift:weildatumdef}
Let $J_S\in M_{n_S}(E)$ and $J_T\in M_{n_T}(E)$ be nondegenerate hermitian matrices and $\xi_0\in E^\times$ with
$\bar\xi_0=-\xi_0$. On $\mathcal W=M_{n_T\times n_S}(\tliftAE)$ let $\langle\cdot,\cdot\rangle_0$ be one half of the alternating
$\A$-bilinear form attached to $(J_S,J_T,\xi_0)$, and let $\Uni_{J_S}(\A)\times\Uni_{J_T}(\A)$ act by
$(g,h)\cdot w=h\,w\,g^{-1}$ (in the matrix realisations of the groups). The Heisenberg operators are
$(\rho_0(w_0)f)(w)=\psi(\langle w,w_0\rangle_0)\,f(w+w_0)$. A \emph{theta function} is a continuous function $f$ on $\mathcal W$
with $f(\gamma+w)=\psi(-\langle\gamma,w\rangle_0)f(w)$ for $\gamma\in M_{n_T\times n_S}(E)$, which is fixed by $\rho_0(u)$ for all $u$
with trivial archimedean components in some open subgroup, and which is smooth along archimedean directions. A
\emph{Weil datum} for $(J_S,J_T,\xi_0,\psi)$ with characters $(\chi_S,\chi_T)$ (Hecke characters with
$\chi_S|_{\A^\times}=\eta^{n_S}$, $\chi_T|_{\A^\times}=\eta^{n_T}$, $\psi$ standard) is a homomorphism $\omega$ from
$\Uni_{J_S}(\A)\times\Uni_{J_T}(\A)$ to the linear operators on functions on $\mathcal W$ which preserves theta
functions, satisfies $\omega(g,h)\rho_0(w_0)=\rho_0((g,h)\cdot w_0)\,\omega(g,h)$ on them, fixes their value at $0$
under rational elements, is continuous pointwise, is unitary for the $L^2$-norm on a fundamental domain of
$\mathcal W/M_{n_T\times n_S}(E)$, and satisfies the following \emph{splitting rule}: if $e\in\tliftAE$ is an idempotent with
$e\bar e=0$ and $e+\bar e$ the idempotent of a place $v$ of $F$, and $g,h$ are supported at $v$, then there is $r>0$
such that $\ell(\omega(g,h)f)=r\,\chi_T(x)^{-1}\chi_S(y)\,\ell(f)$ for every theta function $f$ and every linear form $\ell$
invariant under the $\rho_0(u)$ with $\bar e u=u$, where $x=e\det g+1-e$ and $y=e\det h+1-e$.
\formal{Formalism.IndefR.DatumF.IsWeilDatumOF}
\end{definition}

\begin{proposition}[Weil data exist]\label{pf:tlift:weildatum}
Let the base data be good, $J_S,J_T,\xi_0$ as in \cref{pf:tlift:weildatumdef}, and suppose there is a finitely
generated $\OO_E$-submodule $\Lambda\subset M_{n_T\times n_S}(E)$ of full rank which is self-dual for the pairing
$(M,M')\mapsto\Tr_{E/\Q}(\xi_0\,h_0(M,M'))$ attached to $(J_S,J_T)$. Then for all Hecke characters $\chi_S,\chi_T$ with
$\chi_S|_{\A^\times}=\eta^{n_S}$ and $\chi_T|_{\A^\times}=\eta^{n_T}$ there is a Weil datum with characters $(\chi_S,\chi_T)$.
\formallabel{PositivityTheory.ThKW.S.exists_isWeilDatumOF_k}
\end{proposition}

\begin{proof}
\emph{An auxiliary character.} There is a Hecke character $\eta_1$ of $E$ with $\eta_1|_{\A^\times}=\eta$, archimedean
exponents all equal to $1$, and unramified at the inert places
(\lref{PositivityTheory.ThKW.exists_eta_k}; such characters are classical, e.g.\ the unitary normalisation of a
Hecke character of type $z/|z|$ at every infinite place).

\emph{A Weil datum with characters $(\eta_1^{n_S},\eta_1^{n_T})$.} From the lattice $\Lambda$ one builds the lattice
representation: a representation of $\Uni_{J_S}(\A)\times\Uni_{J_T}(\A)$ on the span of the Heisenberg translates of
the theta functions attached to the vacuum vectors of the lattice models at the finite places
(\lref{UpperTheory.TP.W.latK}, using that theta functions are stable under rational Heisenberg translations,
\lref{GlobalTheory.T.theta_function_heisenberg_translate}). Its splitting factors are computed explicitly in the
lattice models, by lattice indices and signs, and are the characters $\eta_1^{n_S},\eta_1^{n_T}$
(\lref{UpperTheory.TP.W.datumK_k}). This is the classical construction of the Weil representation from lattice
models at the finite places (cf.\ \cite{Weil1964,MoeglinVignerasWaldspurger1987}).

\emph{Changing the characters.} Put $\nu_S=\chi_S\eta_1^{-n_S}$ and $\nu_T=\chi_T\eta_1^{-n_T}$. These are Hecke characters,
and they are trivial on $\A^\times$: on $\A^\times$, $\chi_S$ and $\eta_1^{n_S}$ both restrict to $\eta^{n_S}$, and similarly for $T$.
A Weil datum with characters $(\chi_1,\chi_2)$ can be twisted into one with characters $(\chi_1\nu_1,\chi_2\nu_2)$ for any
Hecke characters $\nu_1,\nu_2$ trivial on $\A^\times$ (\lref{GlobalTheory.T.weil_datum_splitting_characters}). Applied
with $(\nu_S,\nu_T)$ this gives a Weil datum with characters $(\eta_1^{n_S}\nu_S,\eta_1^{n_T}\nu_T)=(\chi_S,\chi_T)$.
\end{proof}

\begin{proposition}[From the theta model to the Schr\"odinger model]\label{pf:tlift:weilpairhp}
In the situation of \cref{pf:tlift:weildatum}, let $\omega_0$ be a Weil datum with characters $(\chi_S,\chi_T)$, let
$J_A={}^tJ_T$, and let $J_B\in M_{n_S}(E)$ and $\xi\in\tliftAE$ satisfy $\xi J_B=\xi_0J_S^{-1}$. Let $e$ be Darboux
coordinates for the alternating form built from $(\xi,J_A,J_B)$ which map $M_{n_T\times n_S}(E)$ into rational vectors.
Then there is a homomorphism $\omega$ from $\Uni_{J_A}(\A)\times\Uni_{J_B}(\A)$ to the operators on functions on $\A^\kappa$
which satisfies the conditions (W1)--(W3) of \cref{pf:tlift:weil} for these coordinates, for the alternating form built
from $(\xi,J_A,J_B)$, with $\Uni_{J_B}$ split by $\chi_T$ and $\Uni_{J_A}$ split by $\chi_S$.
\formallabel{PositivityTheory.ThKW.S.exists_isUWeilPair_cm_k}
\end{proposition}

\begin{proof}
\emph{Transport.} The theta transform $\Phi\mapsto\theta_\Phi$, $\theta_\Phi(w)=\sum_{\xi\in F^\kappa}(\rho(w)\Phi)(\xi)$, maps
Schwartz--Bruhat functions on $\A^\kappa$ to theta functions (\lref{GlobalTheory.TransportHeisF.isThetaFunOF_thetaTr_k}),
and $\omega_0$ is transported to a representation $\omega$ on functions on $\A^\kappa$ with
$\theta_{\omega(p)\Phi}=\omega_0(p)\theta_\Phi$ (\lref{GlobalTheory.T.weil_representation_hermitian_pair}).

\emph{Condition (W1).} $\omega$ preserves $\tliftS(\A^\kappa)$, is continuous in the sense of (W1)
(\lref{GlobalTheory.T.weil_representation_pair_continuous}, \lref{GlobalTheory.T.weil_representation_pair_smooth}),
is unitary, which follows from the unitarity of $\omega_0$ on a fundamental domain
(\lref{GlobalTheory.T.weil_representation_pair_unitary}), and satisfies the Heisenberg relation, which follows from
the Heisenberg equivariance of $\omega_0$ (\lref{GlobalTheory.T.weil_representation_pair_equivariant}).

\emph{Conditions (W2), (W3).} For the factor $\Uni_{J_A}$ there are doubled Darboux coordinates on
$\A^{\kappa\sqcup\kappa}$ and a representation of the unitary group of $\mathrm{diag}(J_A,-J_A)$ which satisfies (W1), the Siegel
rule with the character $\chi_S$ for the diagonal Lagrangian, and the compatibility
$\omega^{(2)}(\mathrm{diag}(g,1))(\Phi_1\boxtimes\Phi_2)=\omega(g,1)\Phi_1\boxtimes\Phi_2$ on Schwartz--Bruhat functions
(\lref{PositivityTheory.ThKW.S.exists_dbl_target_k}). The same statement for the factor $\Uni_{J_B}$, with the
character $\chi_T$, is deduced from it by exchanging the roles of the two spaces
(\lref{PositivityTheory.ThKW.dblS_of_dblT_k}).
\end{proof}

\begin{proof}[Proof of \cref{pf:tlift:weilex}]
Since $E$ contains $\sqrt{-D}$ with $D>0$ (G2), there is $\xi_0\in E^\times$ with $\bar\xi_0=-\xi_0$
(take $\xi_0=\sqrt{-D}$, whose images under all $\tau$ are purely imaginary). The Gram matrices $J_X$ and $J_W$ are invertible, since
$\tau_w(J_X)$ and $i\tau_w(J_W)$ have signatures $(r,s)$ with $r+s$ equal to the dimension. Put
$J_S=-\xi_0J_W^{-1}$ and $J_T={}^tJ_X$. Both are hermitian: ${}^t\bar J_W=-J_W$ gives
${}^t\overline{J_W^{-1}}=-J_W^{-1}$, hence ${}^t\bar J_S=-\bar\xi_0\,{}^t\overline{J_W^{-1}}=-\xi_0J_W^{-1}=J_S$. A lattice $\Lambda$ as in
\cref{pf:tlift:weildatum} is built from $L_X$ and $L_W$ (\lref{PositivityTheory.ThKW.latH_of}), and the restriction
conditions on $\chi_X,\chi_W$ have the form required there
(\lref{GlobalTheory.T.restricted_eta_power_character}). \Cref{pf:tlift:weildatum} gives a Weil datum for
$(J_S,J_T,\xi_0)$ with characters $(\chi_W,\chi_X)$. Finally ${}^tJ_T=J_X$ and
$\xi_0J_S^{-1}=\xi_0(-\xi_0)^{-1}J_W=-J_W$, so \cref{pf:tlift:weilpairhp} applies with $(J_A,J_B,\xi)=(J_X,J_W,-1)$ and the
Darboux coordinates of $M$, which are rational by (M1). It gives a Weil representation for $(M,\psi,\chi_X,\chi_W)$.
\end{proof}

\subsubsection*{The Siegel--Weil formula: published form and formal form}
For a hermitian space $V$ of dimension $m$ and Witt index $r$, a split skew-hermitian space of dimension $2k$ with
unitary group $\Uni(k,k)$ and Siegel parabolic $P$, a $K$-finite $\phi\in\tliftS(V^k(\A))$ and the Weil representation
$\omega$ of \cref{pf:tlift:litweil}, let $E(g,s,\phi)=\sum_{\gamma\in P(F)\backslash\Uni(k,k)(F)}\Phi(\gamma g,s)$ with the standard
section $\Phi(g,s)=(\omega(g)\phi)(0)\,|a(g)|^{s-s_0}$, $s_0=\frac{m-k}2$, and let
$I(g,\phi)=\int_{\Uni(V)(F)\backslash\Uni(V)(\A)}\theta(g,h;\phi)\,dh$ with the invariant probability measure.

\begin{littheorem}[{\cite{Weil1965,KudlaRallis1988,Ichino2007}}]\label{pf:tlift:litsw}
\begin{enumerate}[label=(\alph*),nosep]
\item \cite{Weil1965} If $m>2k$, then $E(g,s,\phi)$ converges absolutely at $s=s_0$ and $E(g,s_0,\phi)=I(g,\phi)$.
\item \cite{Ichino2007} If $V$ is anisotropic or $m-r>k$, then $E(g,s,\phi)$ is holomorphic at $s=s_0$ and
  $E(g,s_0,\phi)=\kappa\,I(g,\phi)$, where $\kappa=1$ if $m>k$ and $\kappa=2$ if $m\le k$.
\end{enumerate}
\end{littheorem}

Part (b) is the unitary analogue of the theorem of Kudla and Rallis for orthogonal--symplectic pairs
\cite{KudlaRallis1988}; see also \cite{Yamana2014} for all classical dual pairs.

\begin{remark}[The formal statement]\label{pf:tlift:remsw}
\Cref{pf:tlift:swstd} is the case of \cref{pf:tlift:litsw}(b) in which the space carrying the theta integral is
definite at every infinite place, hence anisotropic, in the following form.
\begin{enumerate}[label=(\alph*),nosep]
\item The section is $f(x)=\ell(\omega_2(x)\Psi)$ for the admissible linear form $\ell$ of doubling data
  (\cref{pf:tlift:dbl}); $\ell$ plays the role of the evaluation at $0$ in a model adapted to the Siegel parabolic,
  and its normalisation $\sum_y\ell(\rho_2(y\,|\,{-y})\Psi)=\sum_\xi\Psi(\xi)$ ties it to the theta distribution. The
  Weil representations are those of \cref{pf:tlift:weil}, in arbitrary rational Darboux coordinates.
\item The maximal compact subgroup and the height are attached to a frame, $\Psi$ is finite under this compact
  subgroup, and the value is normalised by the volume of a fundamental domain.
\item Both orders (the anisotropic space hermitian and the other skew-hermitian, or conversely) and all dimensions
  $m'=\dim A\ge1$, $m=\dim Y\ge1$ are covered; with $(m,k)$ of \cref{pf:tlift:litsw} replaced by $(m',m)$ the
  constant $\kappa$ is the same.
\item The continuation is asserted pointwise in $x$, as a meromorphic function of $s$ analytic at $s_0$.
\end{enumerate}
In this section the formula is used only for $m'=d\ge2n=2m$, that is, in Weil's range of \cref{pf:tlift:litsw}(a) when
$d>2n$, and on its boundary when $d=2n$.
\end{remark}

The formal proof of \cref{pf:tlift:swstd} is an induction on the rank of the doubled group which compares Fourier
coefficients along the Siegel parabolic. We describe its structure and convert the two statements on the
degenerate coefficients of Eisenstein series which carry the induction.

\begin{definition}[Siegel data; Fourier coefficients; lower-rank data]\label{pf:tlift:sdata}
A \emph{Siegel datum} consists of a definite space $A$ of dimension $m'$, a space $Y$ of dimension $m$, a unitary Hecke
character $\chi$, doubling data $(e^{(2)},\omega_2,S,\ell)$ for $(A,Y,\chi)$, a frame $b$ of $Y$ with its maximal compact
subgroup $\tliftKfr$ and a height $\tliftht_0$, and a Haar measure with fundamental domain on $\Uni(A)(\A)$, subject to the
conditions of \cref{pf:tlift:swstd}; $G=G_2(Y)$ and $P$ is its Siegel parabolic. For a section $f$ of exponent $s_1$
and $z\in\C$ put $\phi_z=\tliftht_0^{\,z-s_1}f$ and $E_z(x)=\sum_{\gamma\in P(F)\backslash G(F)}\phi_z(\gamma x)$.
\begin{enumerate}[label=(\alph*),nosep]
\item \emph{Coefficient data} $(n,\lambda,\mathcal G)$ consist of a parametrisation $\beta\mapsto n(\beta)$ of the unipotent
  radical $N(\A)$ of $P(\A)$, written in the frame $b$, by an additive group $\mathcal N(\A)$ of matrices (with rational points
  $\mathcal N(F)$ corresponding to $N(F)$), a Haar measure $\lambda$ on $\mathcal N(\A)$, and a fundamental domain $\mathcal G$ for $\mathcal N(F)$ with
  $\lambda(\mathcal G)=1$. For $T\in\mathcal N(F)$ the $T$-th \emph{Fourier coefficient} of a function $\phi$ on
  $G(\A)$ is $\mathrm{fc}_T(\phi)(x)=\int_{\mathcal G}\phi(n(\beta)x)\,\psi(-\Tr\tr(T\beta))\,d\lambda(\beta)$.
\item A \emph{lower-rank datum} consists of a Siegel datum for a space $Y'$ of dimension $m-1$ (with the same $A$), an
  index $i_0$, an injection of the other indices, and a continuous homomorphism $j\colon G_2(Y')(\A)\to G(\A)$ which in
  the frames is the identity in the $i_0$-th coordinates and the given inclusion in the others, compatible with
  the Weil representations.
\item \emph{Unipotent data} $(u,\nu,\mathcal U)$ consist of the map $c\mapsto u(c)$ from $\tliftAE^{m-1}$ to the Levi subgroup of $P$
  whose first block is $1+\sum_jc_j\epsilon_{i_0,j}$ in the frame, a Haar measure $\nu$, and a fundamental domain $\mathcal U$ for
  $E^{m-1}$ with $\nu(\mathcal U)=1$.
\item For $T\in\mathcal N(F)$ and $T'\in\mathcal N'(F)$ we write $T\rhd T'$ if the $i_0$-th row and column of $T$ vanish and
  $T'$ is $T$ with them removed.
\item The \emph{first orbit} consists of the $\gamma\in G(F)$ whose lower left block has vanishing $i_0$-th column in the
  frame. For every infinite place $w$ there is a family of \emph{dilations} $\iota_w(y)\in G(F_w)$ in the $i_0$-th
  coordinate, with parameter $\tau\in\R$.
\end{enumerate}
\formal{PositivityTheory.SiegelWeilFormula.SD_k}
\end{definition}

\begin{lemma}[Unfolding of the second orbit]\label{pf:tlift:unf}
Let a Siegel datum, a lower-rank datum, coefficient data $(n,\lambda,\mathcal G)$ and $(n',\lambda',\mathcal G')$ for $Y$ and $Y'$,
unipotent data $(u,\nu,\mathcal U)$, $T\rhd T'$, a section $f$ and $z\in\C$ be given. Let $w_1\in G(F)$ be an element which in
the frame is the identity except in the $i_0$-th coordinates, where it is $\bigl(\begin{smallmatrix}0&a\\ b&0\end{smallmatrix}\bigr)$ with
$a,b\in E^\times$; let $\delta\in E^\times$ and let $c=(c_1,t)\mapsto c_M(c)$ be a continuous map from $\tliftAE^{m-1}\times\A$ to
$P(\A)$ whose Levi part is $u(c_1)$ and whose unipotent part has the single entry $t\delta$ at the place $(i_0,i_0)$; and
let $\kappa$ be a Haar measure on $\tliftAE^{m-1}\times\A$. Assume that for every $x$ and every $\gamma\in G(F)$ the function
$(c,\beta)\mapsto\phi_z(\gamma\,n(\beta)\,u(c)\,x)\,\psi(-\Tr\tr(T\beta))$ is integrable on $\mathcal U\times\mathcal G$, and that the integrals of its absolute
value are summable over representatives $\gamma$ of $P(F)\backslash G(F)$. Let
\[
  S_2(x)=\sum_{\gamma}\int_{\mathcal U}\int_{\mathcal G}\phi_z\bigl(\gamma\,n(\beta)\,u(c)\,x\bigr)\,\psi(-\Tr\tr(T\beta))\,d\lambda(\beta)\,d\nu(c)
\]
(the \emph{second-orbit section}), the sum being over representatives $\gamma$ of the cosets outside the first orbit.
Then there is $K\in\C$ such that for all $x$
\[
  S_2(x)=K\sum_{\gamma'\in P'(F)\backslash G_2(Y')(F)}\int_{\mathcal G'}\Bigl(\int\phi_z\bigl(w_1\,c_M(c)\,j(\gamma'n'(\beta'))\,x\bigr)\,d\kappa(c)\Bigr)
  \psi(-\Tr\tr(T'\beta'))\,d\lambda'(\beta').
\]
\formallabel{PositivityTheory.SiegelWeilFormula.Act.C5.unf_k}
\end{lemma}

\begin{proof}
\emph{Orbits.} The cosets outside the first orbit are in bijection with the pairs $(\gamma',r)$, $\gamma'\in P'(F)\backslash G_2(Y')(F)$
and $r\in E^{m-1}\times F$, through $\gamma=p\,w_1\,c_M(r)\,j(\gamma')$ with $p\in P(F)$
(\lref{PositivityTheory.SiegelWeilFormula.Act.C5.second_orbit_bijection}); since $\phi_z$ is left
$P(F)$-invariant ($\chi$ and the height are trivial on $P(F)$), $p$ may be dropped.

\emph{Coordinates.} Let $\mathcal A$ be a fundamental domain of $\A/F$ of finite measure. There is an isomorphism
$\beta\mapsto(c(\beta),\beta')$ of $\mathcal N(\A)$ onto $(\tliftAE^{m-1}\times\A)\times\mathcal N'(\A)$ whose image measure is $c_P$ times the product
of the Haar measures, for a constant $c_P>0$. In these coordinates, for every function $\Psi$ on $G(\A)$ which is left
invariant under the rational elements fixing the $i_0$-th basis vector, the integral of
$\Psi(n(\beta)u(c)x)\psi(-\Tr\tr(T\beta))$ over $\mathcal U\times\mathcal G$ equals $c_P$ times the integral over
$\mathcal G'\times\mathcal U\times\mathcal U\times\mathcal A$ of $\Psi(s(c'')\,c_M(c',t)\,j(n'(\beta'))\,x)\,\psi(-\Tr\tr(T'\beta'))$, where $c''\mapsto s(c'')$ is an
explicit map into $G(\A)$; here the character becomes $\psi(-\Tr\tr(T'\beta'))$ because the $i_0$-th row and column of $T$
vanish. The same holds for nonnegative functions, and $\kappa$ is $c_K$ times the product of the Haar measures used
(\lref{PositivityTheory.SiegelWeilFormula.Act.C5.coordsC_k},
\lref{PositivityTheory.SiegelWeilFormula.Act.C5.coordsE_k},
\lref{PositivityTheory.SiegelWeilFormula.Act.C5.lam_map_k},
\lref{PositivityTheory.SiegelWeilFormula.Act.C5.kap_eq_k}). These are applied to the second-orbit terms, which
have the required invariance after the orbit decomposition.

\emph{Unfolding.} The same computation for $|\phi_z|$, together with the assumed absolute convergence, shows that the
double series over $(\gamma',r)$ of the absolute values of the integrals is finite. By Tonelli's and Fubini's theorems
the sum over $r\in E^{m-1}\times F$ and the integral over the fundamental domain $\mathcal U\times\mathcal A$ combine into the integral over
$\tliftAE^{m-1}\times\A$ with respect to $\kappa$, and the claim follows with $K=c_P/c_K$.
\end{proof}

\begin{proposition}[Splitting of the degenerate coefficients]\label{pf:tlift:ctsplit}
Let a Siegel datum, a lower-rank datum and the dilations $\iota_w$ be given, and let $f$ be a section of exponent $s_1$
for $(\tliftKfr,\chi)$ such that $E_z$ has a meromorphic continuation in the following sense: for every $s'\in\C$ there are an
open set containing $s'$ and a right half plane, a nonzero polynomial $p$ and a holomorphic, continuous, left
$G(F)$-invariant function $\tliftE$ on it with $\tliftE(z,\cdot)=p(z)E_z$ where the series converges locally uniformly. Then
there is $t\in\R$ such that for every $z$ with $\Re z>t$, all coefficient data for $Y$ and $Y'$, all unipotent data
and all $T\rhd T'$ there is a function $c_2$ on $G(\A)$ with
\begin{equation}\label{eq:tlift:ctsplit}
  \int_{\mathcal U}\mathrm{fc}_T(E_z)\bigl(u(c)x\bigr)\,d\nu(c)=\mathrm{fc}'_{T'}\Bigl(g\mapsto\sum_{\gamma'\in P'(F)\backslash G_2(Y')(F)}
  \phi_z\bigl(j(\gamma'g)\,x\bigr)\Bigr)(1)+c_2(x)
\end{equation}
for all $x$, and $c_2(\iota_w(y)x)=e^{\tau(m-2z)}c_2(x)$ for every infinite place $w$ and every dilation $y$ with
parameter $\tau$.
\formallabel{PositivityTheory.SiegelWeilFormula.Act.ctSplit_of_k}
\end{proposition}

\begin{proof}
\emph{Choices.} Fundamental domains $\mathcal G,\mathcal U,\mathcal G'$ with compact closure exist, and both sides of
\eqref{eq:tlift:ctsplit} do not depend on the choice of the fundamental domains
(\lref{PositivityTheory.SiegelWeilFormula.Act.cptDom_k},
\lref{PositivityTheory.SiegelWeilFormula.Act.indepDom_k}); so we may assume them relatively compact. There is $t$
such that for $\Re z>t$ the series defining $E_z$, integrated over $\mathcal U\times\mathcal G$ against the character, converges
absolutely (\lref{PositivityTheory.SiegelWeilFormula.Act.absConv_k}); then the left side of
\eqref{eq:tlift:ctsplit} is the sum over $\gamma\in P(F)\backslash G(F)$ of the integrals
$\int_{\mathcal U}\int_{\mathcal G}\phi_z(\gamma n(\beta)u(c)x)\psi(-\Tr\tr(T\beta))\,d\lambda\,d\nu$.

\emph{First orbit.} For $\Re z$ large the sum over the cosets in the first orbit equals the first term on the right of
\eqref{eq:tlift:ctsplit} (\lref{PositivityTheory.SiegelWeilFormula.Act.firstEq_k}).

\emph{Second orbit.} Let $c_2=S_2$ be the sum over the remaining cosets. The development provides $w_1$, $c_M$ and $\kappa$ as
in \cref{pf:tlift:unf} such that (i) $j(g)$ commutes with $\iota_w(y)$; (ii) $w_1\iota_w(y)=p\,w_1$ with $p\in P(\A)$ whose Levi
component $t_p$ satisfies $|\det t_p|_{\tliftAE}=e^{-2\tau}$ and $\chi(\det t_p)=1$; (iii) $c_M(c)\iota_w(y)=\iota_w(y)c_M(\sigma(c))$ for a
measurable bijection $\sigma$ with $\sigma_*\kappa=e^{2m\tau}\kappa$
(\lref{PositivityTheory.SiegelWeilFormula.Act.cLawEx_k}). By \cref{pf:tlift:unf}, $c_2(\iota_w(y)x)$ is $K$ times the
sum over $\gamma'$ and the integral over $\beta'$ of $\int\phi_z(w_1c_M(c)j(\gamma'n'(\beta'))\iota_w(y)x)\,d\kappa(c)$. Moving $\iota_w(y)$ to
the left with (i) and (iii) and substituting $c\mapsto\sigma^{-1}(c)$ gives the factor $e^{2m\tau}$; then (ii) and the law
$\phi_z(px)=\chi(\det t_p)|\det t_p|_{\tliftAE}^{z+m/2}\phi_z(x)$ (the section law with exponent $s_1$, multiplied by the
height to the power $z-s_1$) give the factor $e^{-2\tau(z+m/2)}$. Altogether $c_2(\iota_w(y)x)=e^{\tau(m-2z)}c_2(x)$
(\lref{PositivityTheory.SiegelWeilFormula.Act.secUnf_of_k},
\lref{PositivityTheory.SiegelWeilFormula.Act.secLaw_of_k}).
\end{proof}

The formal statement of \cref{pf:tlift:ctsplit} has two further hypotheses, two auxiliary propositions about Siegel
data and their Fourier coefficients, which are not used in its proof.

\begin{theorem}[Siegel--Weil formula from the splitting of the degenerate coefficients]\label{pf:tlift:swfromct}
If the conclusion of \cref{pf:tlift:ctsplit} holds for all Siegel data, then \cref{pf:tlift:swstd} holds.
\formallabel{PositivityTheory.SiegelWeilFormula.Fin.siegelWeilStd_of_ct_k}
\end{theorem}

\begin{proof}
The formal proof assembles the following statements, all proved in the development.
\begin{enumerate}[label=(\arabic*),nosep]
\item \emph{Poisson summation:} for the self-dual measure $\nu$ on the space of the doubled Schr\"odinger model and every
  Schwartz--Bruhat $f$, the sum over rational points of the Fourier transform of $f$ (for the pairing
  $(X,Z)\mapsto\psi(\pm\mathrm{tr}\,\mathrm{gram}(X,Z))$) equals the sum of $f$ over rational points
  (\lref{PositivityTheory.SiegelWeilFormula.TW.Twp.poisson_k}).
\item \emph{Automorphy of the doubled theta function:} the Weyl element of $G$ acts in the model adapted to the Siegel
  parabolic by this Fourier transform, the Siegel parabolic acts by the Siegel rule, and consequently
  $x\mapsto\theta_2(x,h;\Psi)$ is left $G(F)$-invariant
  (\lref{PositivityTheory.SiegelWeilFormula.Fin.tWeyl_of_k},
  \lref{PositivityTheory.SiegelWeilFormula.Fin.tAll_of_k}).
\item \emph{Finiteness:} a Schwartz--Bruhat function on the doubled space is finite under the maximal compact
  subgroup if and only if its image in the model adapted to the Siegel parabolic is a finite datum tensored with a
  polynomial Fock vector (\lref{PositivityTheory.SiegelWeilFormula.KFP.kFinPoly_of_k},
  \lref{PositivityTheory.SiegelWeilFormula.KK.polyKFin_of_k}).
\item \emph{Lower-rank data} exist for $m\ge2$ (\lref{PositivityTheory.SiegelWeilFormula.Ald.ldEx_of_k}).
\item \emph{Nondegenerate coefficients:} for nondegenerate $T$ and finite $\Phi$, the Jacquet integral
  $z\mapsto\int_{\mathcal N(\A)}\phi_z(w\,n(\beta)x)\psi(-\Tr\tr(T\beta))\,d\lambda(\beta)$ ($w$ a Weyl element) is holomorphic at $s_0$ with
  value $c\cdot\mathrm{orb}_T(\omega_2(x)\Phi)$, where $\mathrm{orb}_T$ is the average over a fundamental domain of $\Uni(A)(\A)$ of the
  sum of the adapted model of $\Phi$ over the rational $X$ with $\mathrm{gram}(X)=-2T$, and $c=1$ if $m<m'$, $c=2$ if $m=m'$, and
  $c=0$ (for vacuum vectors) if $m'<m$ (\lref{PositivityTheory.SiegelWeilFormula.KWhit.whitK_of_k},
  \lref{PositivityTheory.SiegelWeilFormula.KWhit.whitEq_of_k},
  \lref{PositivityTheory.SiegelWeilFormula.KWhit.whitLow_of_k}); and $\mathrm{orb}_T(\omega_2(x)\Phi)$ is the $T$-th Fourier
  coefficient of the theta integral (\lref{PositivityTheory.SiegelWeilFormula.Fin.thetaCoef_k}).
\item \emph{Induction step} (for $m'<m$): the identity $E(x,s_0)=2I(x)$ for vacuum vectors of a lower-rank datum
  implies the same identity for the datum itself; this step uses \cref{pf:tlift:ctsplit}
  (\lref{PositivityTheory.SiegelWeilFormula.AStep.stepOf_k}).
\item \emph{Conclusion:} from (1)--(6) the identity $E(x,s_0)=\kappa\,I(x)$ follows, first for vacuum vectors and then,
  since these generate the finite vectors, for all finite $\Psi$; the meromorphic continuation and the analyticity at
  $s_0$ are obtained along the way (\lref{PositivityTheory.SiegelWeilFormula.Atop.siegelWeilStd_of_k},
  \lref{PositivityTheory.SiegelWeilFormula.Fin.siegelWeilStd_of_k}).
\end{enumerate}
\end{proof}

\begin{proof}[Proof of \cref{pf:tlift:swstd}]
Combine \cref{pf:tlift:swfromct,pf:tlift:ctsplit}.
\end{proof}

% ---------------------------------------------------------------------------
\subsection{The doubled group: compactness, heights and integrability}\label{pf:tlift:dblgrp}

In this subsection the base data are good and $Y=(J,L,\vartheta,b)$ is a hermitian or skew-hermitian space of
dimension $m$ (of any signature); $G_2=G_2(Y)$, $P$, $\iota$ and $\tliftK$ are as in
\cref{pf:tlift:stab2}.

\begin{lemma}[Compactness of $\tliftK$]\label{pf:tlift:stab2cpt}
The standard maximal compact subgroup $\tliftK\subset G_2(\A)$ is compact.
\formallabel{GlobalTheory.ThK.Dbl.stab2Cpt_k}
\end{lemma}

\begin{proof}
Let $c=1$ if $Y$ is hermitian and $c=i$ if it is skew-hermitian, so that $H_w=c\,\tau_w(J)$ is a hermitian
matrix with Cartan involution $\vartheta_w$. For $k\in\tliftK$ the matrix $A=\tau_w(k)$ satisfies
${}^tA\,\mathrm{diag}(H_w,-H_w)\,\bar A=\mathrm{diag}(H_w,-H_w)$ and commutes with $\mathrm{diag}(\vartheta_w,-\vartheta_w)$, which is a
Cartan involution of $\mathrm{diag}(H_w,-H_w)$; so $A$ preserves the positive definite hermitian form
$z\mapsto\Re\,h_{\mathrm{diag}(H_w,-H_w)}(\mathrm{diag}(\vartheta_w,-\vartheta_w)z,z)$, and its entries are bounded by a constant
$R_w$ depending only on $H_w$ and $\vartheta_w$ (\lref{GlobalTheory.ThK.Dbl.cartan_bound2}). At the finite
places, $k$ maps $\widehat L\oplus\widehat L$ into itself; since $L$ is a lattice, there are compact sets
$C_{ij}\subset\tliftAE$ such that every matrix mapping $\widehat L$ into itself has its finite $(i,j)$-entry in
$C_{ij}$ (\lref{GlobalTheory.ThK.Dbl.St.fin_bound}). Writing an adele of $E$ as the sum of its archimedean
and finite parts, the entries of $k$ and of $k^{-1}\in\tliftK$ lie in fixed compact subsets of $\tliftAE$. The
group $\tliftK$ is closed in $G_2(\A)$ (commuting with a fixed matrix at infinity and mapping the closed set
$\widehat L\oplus\widehat L$ into itself are closed conditions), and $G_2(\A)$ is closed in
$\GL_{2m}(\tliftAE)$; a closed set of invertible matrices whose entries and whose inverses' entries lie in
compact sets is compact (\lref{GlobalTheory.ThK.Dbl.St.isCompact_of_units}).
\end{proof}

\begin{proposition}[Iwasawa decomposition]\label{pf:tlift:iwasawa}
$G_2(\A)=P(\A)\,\tliftK$.
\formallabel{GlobalTheory.ThK.Dbl.Hgt.iwasawa2_k}
\end{proposition}

\begin{proof}
Write $H=\mathrm{diag}(J,-J)$, so that $G_2=\Uni_H$, and $\tliftdia={1_m\choose1_m}$; then $P$ is the stabiliser of the Lagrangian
$\tliftdia$ (the column space of $\tliftdia$), and $x\in P(\A)\tliftK$ if and only if $x\tliftdia=k\tliftdia\,t$ for some $k\in\tliftK$ and some
invertible $t$, i.e.\ if and only if the Lagrangian $x\tliftdia$ is a $\tliftK$-translate of $\tliftdia$. Multiplying $J$ by a scalar (and
transposing) does not change the groups involved, so we may assume that $H$ is hermitian.

\emph{Finite places.} Let $v$ be finite, $\OO=\OO_E\otimes\OO_v$ and $M_v=\widehat L_v\oplus\widehat L_v$, where $\widehat L_v$ is the
$v$-component of $\widehat L$; let $K_v$ be the stabiliser of $M_v$ in $\Uni_H(F_v)$. In a frame of $\widehat L_v$ (an
$\OO$-basis, which exists because $\OO$ is a finite product of discrete valuation rings) the matrix of $lH$ is integral with
integral inverse for some $l\in E^\times$, by the self-duality of $L$. Let $y\in\Uni_H(F_v)$ and $\mathcal L=y\tliftdia\otimes F_v$.
Because $\OO$ is a product of discrete valuation rings, $\mathcal L\cap M_v$ is a direct summand of $M_v$ of rank $m$ (the columns
of $y\tliftdia$ can be rescaled by some $t\in\GL_m(E\otimes F_v)$ to an integral matrix with an integral left inverse). Choose
an $\OO$-basis $n_1,\dots,n_m$ of $\mathcal L\cap M_v$. Since $lH$ is perfect on $M_v$ and $\mathcal L$ is isotropic, there are
$f_1,\dots,f_m\in M_v$ with $lH(n_i,f_j)=\delta_{ij}$, and after subtracting from each $f_j$ a suitable combination of the
$n_k$ (whose coefficients are $w_0$ times the values of $lH$ on the $f$'s, for an element $w_0\in\OO_E$ with $w_0+\bar w_0=1$,
which exists because $E/F$ is unramified at all finite places,
\lref{GlobalTheory.exists_int_add_conj_eq_one_of_base_good_k}) the $f_j$ are isotropic as well. Thus $M_v$ has a
basis $(n_i,f_j)$ in which $lH$ has the standard antidiagonal form and whose first half spans $\mathcal L\cap M_v$; the same
construction gives such a basis whose first half spans $\tliftdia\cap M_v$. The element of $\Uni_H(F_v)$ which maps the second
basis to the first lies in $K_v$ and maps $\tliftdia$ to $\mathcal L$. Hence $K_v$ acts transitively on Lagrangians and
$\Uni_H(F_v)=K_vP(F_v)=P(F_v)K_v$, the second equality by taking inverses.

\emph{Infinite places.} Let $w$ be infinite and $\Theta=\mathrm{diag}(\vartheta_w,-\vartheta_w)$, a Cartan involution of $\tau_w(H)$. Then
$\C^{2m}=V_+\oplus V_-$ with $V_\pm$ the $\pm1$-eigenspaces of $\Theta$, on which $\tau_w(H)$ is positive, resp.\ negative,
definite. A Lagrangian $\mathcal L$ meets $V_\pm$ trivially, hence is the graph of a linear map $T\colon V_+\to V_-$, and
isotropy says $h(Tv,Tv)=-h(v,v)$, where $h$ is the form of $\tau_w(H)$. The group of $k\in\Uni_H(F_w)$ commuting with $\Theta$ is
$\Uni(V_+)\times\Uni(V_-)$, and it acts transitively on such graphs. Hence again $\Uni_H(F_w)=P(F_w)K_w$, $K_w$ the
centraliser of $\Theta$; in the formal proof this is done in frames adapted to the Cartan involutions
(\lref{GlobalTheory.ThK.Dbl.Frame.exists_arch}).

\emph{Assembly.} Let $x\in G_2(\A)$. The matrices $x$ and $x^{-1}$ preserve $M_v$ for all but finitely many $v$. At those
places put $p_v=1$, $k_v=x_v$; at the finitely many remaining finite places and at the infinite places decompose
$x_v=p_vk_v$ as above. Then $p=(p_v)\in P(\A)$, $k=(k_v)\in\tliftK$ (the conditions defining $\tliftK$ are local), and $x=pk$.
\end{proof}

\begin{proposition}[Iwasawa decomposition and heights]\label{pf:tlift:hgtex}
$G_2(\A)=P(\A)\tliftK$; there is a height $\tliftht$ for $\tliftK$ (\cref{pf:tlift:stab2}), and every height
is continuous and positive; and $\iota\colon\Uni(Y)(\A)\times\Uni(Y)(\A)\to G_2(\A)$,
$(h,h')\mapsto\mathrm{diag}(h,h')$, is a continuous homomorphism.
\formallabel{GlobalTheory.ThK.Dbl.hgtEx_of_k}\formallabel{GlobalTheory.ThK.Dbl.Hgt.hgtEx_core_k}
\end{proposition}

\begin{proof}
The decomposition is \cref{pf:tlift:iwasawa}. Given it, put
$\tliftht(pk)=|\det t_p|_{\tliftAE}$. This is well defined: if $pk=p'k'$ then
$p'^{-1}p=k'k^{-1}\in P(\A)\cap\tliftK$, and $p\mapsto|\det t_p|_{\tliftAE}$ is a continuous homomorphism on the
closed subgroup $P(\A)$, so its image on the compact group $P(\A)\cap\tliftK$ is a compact subgroup of
$\R_{>0}$, i.e.\ trivial. Continuity of $\tliftht$
is deduced from the same two facts, compactness of $\tliftK$ and continuity of $p\mapsto|\det t_p|_{\tliftAE}$.
Any two heights for $\tliftK$ coincide, since they agree on every $pk$. Positivity is clear and $\iota$ is
obviously a continuous homomorphism.
\end{proof}

\begin{lemma}[Integrability from a distribution bound]\label{pf:tlift:count}
Let $\mu$ be a measure on a measurable space $G$ and $f\colon G\to\R$ measurable with $0<f\le R$ for some
$R\ge1$. Suppose there are $C<\infty$ and $b\ge0$ with $\mu\{g:f(g)\ge2^{-j}\}\le C\,b^j$ for all $j\in\N$.
Then there is $t_0$ such that $f^t$ is $\mu$-integrable for every $t\ge t_0$.
\formallabel{GlobalTheory.ThK.Dbl.Hgt.integrable_of_count}
\end{lemma}

\begin{proof}
Choose $k\in\N$ with $2^{-k}<1/(b+1)$ and let $t\ge t_0=k$. The sets $P_j=\{2^{-j}\le f\le2R\cdot2^{-j}\}$,
$j\ge0$, cover $G$: given $g$, let $j$ be minimal with $2^{-j}\le f(g)$; if $j=0$ then $f(g)\le R\le 2R$,
and if $j\ge1$ then $f(g)<2^{-(j-1)}=2\cdot2^{-j}\le2R\cdot2^{-j}$. On $P_j$ we have
$f^t\le(2R)^t2^{-jt}$, and $\mu(P_j)\le Cb^j$. Hence
\[
  \int f^t\,d\mu\le\sum_{j\ge0}(2R)^t\,C\,(2^{-t}b)^j<\infty,
\]
because $2^{-t}b\le2^{-k}b<b/(b+1)<1$.
\end{proof}

The following counting lemma supplies the distribution bound for heights. Let $q$ denote the rank of a free
$\A$-module.

\begin{lemma}[Counting]\label{pf:tlift:countad}
Let $B$ be a ring which is an $\A$-algebra, with an $\A$-linear homeomorphism $c\colon B\to\A^\kappa$ onto a free module of finite rank
$q=|\kappa|$; let $H\subset B^\times$ be a closed subgroup with a Haar measure $\mu$, and let $\Omega\subset B$ and $\Omega_0\subset\A$ be compact.
Then there are $R\ge1$, $C<\infty$ and $b\ge0$ such that $|u|_{\tliftAE}\le R$ for every $u\in\tliftAE^\times$ with $\Nm u\in\Omega_0$, and for
every $j\in\N$ the set of $g\in H$ for which there is $u\in\tliftAE^\times$ with
\[
  \Nm u\in\Omega_0,\qquad|u|_{\tliftAE}\ge2^{-j},\qquad(\Nm u)\,g\in\Omega,\qquad(\Nm u)\,g^{-1}\in\Omega
\]
has $\mu$-measure at most $Cb^j$. One can take $b=2^{1+4q[F:\Q]}$.
\formallabel{GlobalTheory.ThK.Dbl.Hgt.count_ad}
\end{lemma}

\begin{proof}
Let $\rho\in\A^\times$ have components $\frac12$ at the infinite places and $1$ at the finite places, and let $\nu$ be a Haar measure on
$\A$. Multiplication by $\rho^{-1}$ multiplies $\nu$ by $\delta=2^{[F:\Q]}$, since $F$ is totally real; multiplication by a rational
number $m\ne0$ preserves $\nu$ (product formula).

\emph{Step 1: scaling.} There are $i_0\in\N$ and a compact set $C_0\subset\A$ such that (a) $m\rho^{2i}\in C_0$ for all $i\in\N$ and all
integers $0\le m\le2^i$, and (b) for $u\in\tliftAE^\times$ with $\Nm u\in\Omega_0$ one has $|u|_{\tliftAE}\le2^{i_0}$, and for every $j$ with
$|u|_{\tliftAE}\ge2^{-j}$ there are an integer $0<m\le2^{j+i_0}$ and $x\in C_0$ with $m\rho^{2(j+i_0)}=x\Nm u$. Indeed, by compactness there
are an integer $c\ge1$ and $R_0$ with $|a_w|\le R_0$ at the infinite places and $c\,a_v$ integral at the finite places for $a\in\Omega_0$. Let
$C_1=\{x\in\A:|x_w|\le1,\ |x_v|_v\le1\text{ for all }w,v\}$, $C_0=C_1\cup c^2C_1$, $R_1=\max(1,c^2R_0)$, and choose $i_0$ with
$R_1^{[F:\Q]}\le2^{i_0}$. (a) holds because $|m\rho^{2i}|_w=m4^{-i}\le1$ and $m$ is integral. For (b) put $\nu_u=\Nm(cu)=c^2\Nm u$; then
$|\nu_{u,w}|\le R_1$, $\nu_{u,v}$ is integral, and $|u|_{\tliftAE}=|\nu_u|_\A$ by the product formula, so
$|u|_{\tliftAE}\le R_1^{[F:\Q]}\le2^{i_0}$. If $|u|_{\tliftAE}\ge2^{-j}$, put $i=j+i_0$, so that $R_1^{[F:\Q]}\le2^i|\nu_u|_\A$, and let
$m=\prod_v|\nu_{u,v}|_v^{-1}$, the absolute norm of the ideal generated by the finite part of $\nu_u$; it is a positive integer with
$|m|_v\le|\nu_{u,v}|_v$ for every finite $v$. Since $\prod_w|\nu_{u,w}|\le R_1^{[F:\Q]}$ we get $m\le2^i$, and since
$|\nu_u|_\A\le|\nu_{u,w}|R_1^{[F:\Q]}$ for each $w$ we get $|\nu_{u,w}|\ge2^{-i}$. Hence $x'=m\rho^{2i}\nu_u^{-1}$ satisfies
$|x'_w|\le2^i4^{-i}2^i=1$ and $|x'_v|_v\le1$, so $x'\in C_1$ and $m\rho^{2i}=(c^2x')\Nm u$ with $c^2x'\in C_0$.

\emph{Step 2: packing.} If $V\subset\A$ is a compact neighbourhood of $0$, $Z\subset\A$ has finite measure and $X+V\subset Z$, then there is a
finite set $s\subset\A$ with $|s|\,\nu(V)\le\nu(Z)$ and $X\subset s+(V-V)$: take a maximal family of points $x_i\in X$ with pairwise
disjoint translates $x_i+V$; these lie in $Z$, which bounds their number, and by maximality every $x\in X$ satisfies
$x-x_i\in V-V$ for some $i$.

\emph{Step 3: covering.} Let $\Omega'\subset B$ and $W\subset\A$ be compact. There are compact sets $X\subset\A$ and $U\subset H$ such that for every
$l\in\A^\times$ and every finite set $s$ with $X\subset s+l^2W$,
\[
  \mu\{g\in H:lg\in\Omega',\ lg^{-1}\in\Omega'\}\le|s|^q\,\mu(U).
\]
Take for $X$ the union of the coordinate projections of $c(\Omega')$, and let $\Omega''$ be the compact set of the $1+\omega\,c^{-1}(w_1-w_2)$
with $\omega\in\Omega'$ and $w_1,w_2\in W^\kappa$, and $U=\{h\in H:h,h^{-1}\in\Omega''\}$, which is compact because $H$ is closed in $B^\times$. For
$y\in s^\kappa$ let $S(y)$ be the set of $g$ with $lg^{-1}\in\Omega'$ and $c(lg)_k-y_k\in l^2W$ for all $k$. The set on the left is covered by
the $|s|^q$ sets $S(y)$. If $g,g'\in S(y)$, then $c(lg')-c(lg)=l^2(w'-w)$ with $w,w'\in W^\kappa$, so $g'-g=l\,c^{-1}(w'-w)$ and
$g^{-1}g'=1+(lg^{-1})\,c^{-1}(w'-w)\in\Omega''$, and symmetrically $g'^{-1}g\in\Omega''$. Thus $S(y)\subset gU$, and $\mu(S(y))\le\mu(U)$ by
left invariance.

\emph{Step 4: conclusion.} Apply Step 3 with $\Omega'=C_0\Omega$ and $W=V-V$, $V$ a compact neighbourhood of $0$, and put
$Z=X+C_0C_0V$ and $A_0=\nu(Z)/\nu(V)$. Let $i\in\N$, $0<m\le2^i$ and $l=m\rho^{2i}$. Then $l\in C_0$ by Step 1(a), so $X+l^2V\subset Z$, and
Step 2 (with $l^2V$ in place of $V$) gives $s$ with $X\subset s+l^2W$ and $|s|\le\nu(Z)/\nu(l^2V)=A_0\delta^{4i}$, because $l^2V$ has
measure $\delta^{-4i}\nu(V)$. By Step 3,
\[
  \mu\{g:lg\in C_0\Omega,\ lg^{-1}\in C_0\Omega\}\le(A_0\delta^{4i})^q\mu(U).
\]
Now let $j\in\N$ and let $g,u$ be as in the statement, and put $i=j+i_0$. By Step 1(b), $m\rho^{2i}=x\Nm u$ with $0<m\le2^i$ and
$x\in C_0$, hence $m\rho^{2i}g=x(\Nm u)g\in C_0\Omega$, and likewise for $g^{-1}$. Summing the last estimate over the at most $2^i$
values of $m$, the set in question has measure at most
$2^i(A_0\delta^{4i})^q\mu(U)=A_0^q\mu(U)(2\delta^{4q})^{i_0}\cdot(2\delta^{4q})^j$. This is the claim with $R=2^{i_0}$,
$C=A_0^q\mu(U)(2\delta^{4q})^{i_0}$ and $b=2\delta^{4q}=2^{1+4q[F:\Q]}$.
\end{proof}

\begin{remark}[An explicit threshold]\label{pf:tlift:remdzero}
In \cref{pf:tlift:hgtint} below, \cref{pf:tlift:countad} is applied with $B=M_m(\tliftAE)$, a free $\A$-module of rank $q=2m^2$, so
$b=2^{1+8m^2[F:\Q]}$, and \cref{pf:tlift:count} holds with any integer $t_0$ such that $2^{t_0}>b+1$, for instance
$t_0=2+8m^2[F:\Q]$. Through \cref{pf:tlift:hgtfst} and the proof of \cref{pf:tlift:main} (where $m=n$ and $d_0=\max(N,2n)$
for an integer $N>2t_0$) this shows that $d_0=5+16n^2[F:\Q]$ is an admissible value of the threshold. The formal
statement only asserts the existence of $d_0$.
\end{remark}

\begin{proposition}[Integrability of heights]\label{pf:tlift:hgtint}
Let $\tliftht$ be a height for $\tliftK$ and $\mu$ a Haar measure on $\Uni(Y)(\A)$. There is $t_0$ such that
$g\mapsto\tliftht(\iota(1,g))^t$ is $\mu$-integrable for all $t\ge t_0$.
\formallabel{GlobalTheory.ThK.Dbl.hgtInt_of_k}
\end{proposition}

\begin{proof}
Let $g\in\Uni(Y)(\A)$ and write $\iota(1,g)=pk$ with $p\in P(\A)$, $k\in\tliftK$ and $t=t_p$. Then
$k^{-1}=\iota(1,g^{-1})p$, and since $p\tliftdia=\tliftdia t$,
\begin{equation}\label{eq:tlift:rows}
  k^{-1}\tliftdia=\iota(1,g^{-1})\tliftdia\,t={t\choose g^{-1}t}.
\end{equation}
Let $T(k),T'(k)\in M_m(\tliftAE)$ be the upper and lower blocks of $k^{-1}\tliftdia$; they depend continuously on
$k\in\tliftK$, and $\tliftht(\iota(1,g))=|\det t|_{\tliftAE}$ with $t=T(k)$, $g^{-1}t=T'(k)$. Put
$u=\det t\in\tliftAE^\times$ and $N=\Nm_{\tliftAE/\A}u=u\bar u$. Since $g^{-1}$ is unitary,
$\det(g^{-1})\overline{\det(g^{-1})}=1$, and the elementary identity
\[
  (\det t\,\overline{\det t})\,g=\overline{\det t'}\;t\,\mathrm{adj}(t'),\qquad
  (\det t\,\overline{\det t})\,g^{-1}=\overline{\det t}\;t'\,\mathrm{adj}(t)\qquad(t'=g^{-1}t)
\]
(\lref{GlobalTheory.ThK.Dbl.Hgt.graph_alg}; it follows from $t=g\,t'$, $M\,\mathrm{adj}(M)=\det M$ and
$\det t'=\det g^{-1}\det t$) shows that $N\,g$ and $N\,g^{-1}$ lie in the compact set $\Omega$ obtained as the
image of $\tliftK$ under the continuous maps $k\mapsto\overline{\det T'(k)}\,T(k)\,\mathrm{adj}(T'(k))$ and
$k\mapsto\overline{\det T(k)}\,T'(k)\,\mathrm{adj}(T(k))$ (with $\bar{\ }$ the conjugation of $\tliftAE$), while $N$ lies in
the compact set $\Omega_0=\{\Nm\det T(k):k\in\tliftK\}$.

By \cref{pf:tlift:countad} (applied to the closed subgroup
$\Uni(Y)(\A)$ of $\GL_m(\tliftAE)$, the compact sets $\Omega,\Omega_0$ and $\mu$) there are $R\ge1$, $C<\infty$ and
$b\ge0$ such that $|u|_{\tliftAE}\le R$ for every $u\in\tliftAE^\times$ with $\Nm u\in\Omega_0$, and for every $j$ the set of $g$
for which there is $u\in\tliftAE^\times$ with $\Nm u\in\Omega_0$, $|u|_{\tliftAE}\ge2^{-j}$, $(\Nm u)\,g\in\Omega$ and $(\Nm u)\,g^{-1}\in\Omega$ has
measure at most $Cb^j$. Applying it with $u=\det t$ shows $\tliftht(\iota(1,g))\le R$ and
$\mu\{g:\tliftht(\iota(1,g))\ge2^{-j}\}\le Cb^j$. The function $g\mapsto\tliftht(\iota(1,g))$ is continuous and
positive (\cref{pf:tlift:hgtex}), and \cref{pf:tlift:count} gives $t_0$.
\end{proof}

\begin{proposition}[The integrability exponent]\label{pf:tlift:hgtfst}
Let $X$ be a positive definite hermitian space, $\tliftK$ the standard maximal compact subgroup of
$G_2(X)(\A)$, $\tliftht$ a height for $\tliftK$ and $\iota$ as above. There is $t_0\in\R$ such that for every Haar
measure $\mu$ on $\Uni(X)(\A)$ and every $t\ge t_0$ the function $g\mapsto\tliftht(\iota(g,1))^t$ is $\mu$-integrable.
\formallabel{UpperTheory.TP.hgtFst_k}
\end{proposition}

\begin{proof}
Fix a Haar measure $\mu_0$ and let $t_0$ be given by \cref{pf:tlift:hgtint} for $\mu_0$. Every Haar measure is a
multiple of $\mu_0$, so $g\mapsto\tliftht(\iota(1,g))^t$ is integrable for every Haar measure $\mu$ and $t\ge t_0$. Next,
$\iota(g,1)=\iota(g,g)\,\iota(1,g^{-1})$, where $\iota(g,g)\in P(\A)$ has Levi component $g$, and
$|\det g|_{\tliftAE}=1$ for unitary $g$ (as $\det g\,\overline{\det g}=1$); hence
\begin{equation}\label{eq:tlift:swap}
  \tliftht(\iota(g,1))=\tliftht(\iota(1,g^{-1})).
\end{equation}
Haar measures on $\Uni(X)(\A)$ are right invariant (\cref{pf:tlift:cpt}), so the image of $\mu$ under $g\mapsto g^{-1}$
is a multiple of $\mu$ (\lref{GlobalTheory.ThK.Dbl.inv_scale}), and integrability is preserved under
$g\mapsto g^{-1}$.
\end{proof}

\subsubsection*{An integral decomposition}
The following lemma is used in the proof of \cref{pf:tlift:vacfix}. For a hermitian space $X$ and a finite place $v$ we say that
a matrix $a\in M_n(\tliftAE)$ is \emph{integral at $v$} if its $v$-component, written in an $\OO_E\otimes\OO_v$-basis of $\widehat L_{X,v}$,
has integral entries. We write $a^*$ for the adjoint of $a$ with respect to $J_X$, i.e.\ ${}^ta\,J_X=J_X\,\overline{a^*}$.

\begin{lemma}[Integral decomposition at a place]\label{pf:tlift:beta}
Let $X$ be a hermitian space, $v$ a finite place and $g\in G_2(X)(\A)$ with $g\tliftdia={x\choose y}$. Suppose that $x,y$ and some
$a_1,a_2$ are integral at $v$ with $a_1x+a_2y=1$, and that $x-1$ and $y-1$ are supported at $v$. Then there are $b$ integral at $v$
with $b+b^*=-1$ and $t'$ integral at $v$ with $\bigl(y+b(y-x)\bigr)t'=1$.
\formallabel{GlobalTheory.ThK.Dbl.Vf2.exists_beta_k}
\end{lemma}

\begin{proof}
Let $\OO=\OO_E\otimes\OO_v$; it is either a discrete valuation ring with residue field $k$ (inert $v$) or a product of two
discrete valuation rings with residue field $k\times k$ ($v$ split). Fix an $\OO$-basis of $\widehat L_{X,v}$; in it the Gram matrix $G$ of
$J_X$ is hermitian and integral with integral inverse, by self-duality. Let $\theta\in\OO_E$ with $\theta+\bar\theta=1$
(\lref{GlobalTheory.exists_int_add_conj_eq_one_of_base_good_k}) and $b_0=-\theta\cdot1$; then $b_0^*=-\bar\theta\cdot1$ and
$b_0+b_0^*=-1$. We look for $b=b_0+S$ with $S^*=-S$ (in the basis: ${}^tS\,G+G\,\bar S=0$). Put $M=y+b_0(y-x)$ and
$m=y-x$, written in the basis at $v$; then $y+b(y-x)=M+Sm$. Since $g$ preserves $\mathrm{diag}(J_X,-J_X)$, the columns of
$g\tliftdia$ span an isotropic subspace, which gives ${}^tM\,G\,\bar m+{}^tm\,G\,\bar M=0$ (\lref{GlobalTheory.ThK.Dbl.Vf2.lag_of}); and
$a_1x+a_2y=1$ gives $PM+Qm=1$ with $P=a_1+a_2$ and $Q=-a_1-(a_1+a_2)b_0$, both integral. It suffices to find $S$, $G$-skew with
integral entries, such that $M+Sm$ is invertible over $\OO$: then $t'=(M+Sm)^{-1}$ is integral and $b=b_0+S$ is as required.

Let $r$ be the reduction modulo the maximal ideal of one of the factors of $\OO$, onto a field $K$. An element $z\in\OO$ is a unit if
$r(z)\ne0$ and $r(\bar z)\ne0$. In the inert case the residue involution is nontrivial: there is $\kappa\in\OO$ with $r(\kappa-\bar\kappa)\ne0$,
and $r(z)=0$ implies $r(\bar z)=0$. In the split case there is $\eta\in\OO$ with $r(\eta)=1$ and $r(\bar\eta)=0$
(the residue fields being those of $E_w$, $w\mid v$). We show: if $r(M)$ is singular, there is an integral $G$-skew $S'$ with
$\mathrm{rank}\,r(M+S'm)>\mathrm{rank}\,r(M)$. Since $M+S'm$ again satisfies the two identities above (with $Q-PS'$ in place of $Q$), at
most $n$ such steps give $S$ with $r(\det(M+Sm))\ne0$. Then also $r(\overline{\det(M+Sm)})\ne0$: if $r(\overline{N})y_1=0$ for
$N=M+Sm$, the isotropy identity ${}^tN\,G\,\bar m+{}^tm\,G\,\bar N=0$ gives $r({}^tN)r(G)r(\bar m)y_1=0$, hence $r(\bar m)y_1=0$ since $r(N)$ and
$r(G)$ are invertible, and then the conjugate of $PN+(Q-PS)m=1$ gives $y_1=0$. So $\det(M+Sm)$ is a unit.

For the step, let $y_0\ne0$ with $r(M)y_0=0$, lifted to $\OO^n$. Some coordinate $i$ has $r((my_0)_i)\ne0$, since otherwise
$r(PM+Qm)y_0=0$. Let $e_i$ be the $i$-th basis vector, $w=\overline{G^{-1}e_i}$, so that $G\bar w=e_i$. For every $u\in\OO^n$ the
matrix $S_u=u\,{}^te_i-w\,{}^t(G\bar u)$ is $G$-skew. In the split case take $z\notin\mathrm{im}\,r(M)$ and $u=\eta z$; in the inert case
take $u=\kappa w$ and $z=r(\kappa-\bar\kappa)r(w)$. In both cases $r(S_u)=z\,{}^te_i$, so $r(M+S_um)=r(M)+z\,({}^te_i\,r(m))$, a rank one
perturbation of $r(M)$ which does not vanish on $y_0$; its rank exceeds that of $r(M)$ provided $z\notin\mathrm{im}\,r(M)$. In the inert
case this holds: if $r(M)x_0=z$, then $r(w-Mx_0')=0$ for $x_0'=r(\kappa-\bar\kappa)^{-1}x_0$, hence also $r(\overline{w-Mx_0'})=0$, and
pairing the isotropy identity with $y_0$ and $\bar x_0'$ gives
$(my_0)_i={}^t(my_0)G\,\overline{(w-Mx_0')}-{}^t(My_0)G\,\overline{(mx_0')}$, whose reduction vanishes, a contradiction.
\end{proof}


% ---------------------------------------------------------------------------
\subsection{The Rallis inner product}\label{pf:tlift:rallis}

From now on $X$ is a positive definite hermitian space of dimension $n\ge1$ and $W$ a skew-hermitian
space of dimension $d\ge1$ and signature $(d,0)$, $(M,\chi_X,\chi_W,\omega)$ is a Weil pair for $(X,W)$ with
$M=(e,(\tliftvac_v),\mu,D)$, $\tliftK$ is the standard maximal compact subgroup of $G_2(X)(\A)$ with a height $\tliftht$,
and $\iota(g,g')=\mathrm{diag}(g,g')$. We use the matrix coefficients \eqref{eq:tlift:coef}. A function $\varphi$
on $\Uni(X)(\A)$ which is continuous and left $\Uni(X)(F)$-invariant is bounded (\cref{pf:tlift:cpt}).

\begin{lemma}[Petersson coefficients]\label{pf:tlift:mc}
Let $\varphi$ be continuous and left $\Uni(X)(F)$-invariant. Then $\tliftM_\varphi$ is continuous and bounded, and
$\tliftM_\varphi(g^{-1})=\overline{\tliftM_\varphi(g)}$. If moreover $\varphi\ne0$, then $r=\tliftM_\varphi(1)>0$ and
$|\tliftM_\varphi(g)|\le r$ for all $g$.
\formallabel{UpperTheory.TP.coef_facts_pk}\formallabel{UpperTheory.matrix_coefficient_bounded_by_value_at_one}
\end{lemma}

\begin{proof}
Continuity and boundedness follow from dominated convergence, $\varphi$ being bounded by some $C$ and
$\mu(D)<\infty$. For the symmetry, the function $y\mapsto\varphi(y)\overline{\varphi(yg)}$ is left $\Uni(X)(F)$-invariant
and $Dg^{-1}$ is again a fundamental domain (right translation preserves $\mu$); integrating over $Dg^{-1}$
instead of $D$ gives $\tliftM_\varphi(g^{-1})=\int_D\varphi(yg^{-1})\overline{\varphi(y)}\,d\mu=\int_D\varphi(y)\overline{\varphi(yg)}\,d\mu$.
Next, $\tliftM_\varphi(1)=\int_D|\varphi|^2\,d\mu$; it is nonzero, since a continuous left invariant function with
vanishing integral of $|\varphi|^2$ over a fundamental domain vanishes identically
(\lref{GlobalTheory.ThK.Dbl.Cpt.eq_zero}). The bound $|\tliftM_\varphi(g)|\le\tliftM_\varphi(1)$ is the Cauchy--Schwarz
inequality together with $\int_D|\varphi(yg)|^2\,d\mu=\int_D|\varphi|^2\,d\mu$, which holds because $Dg$ is again a
fundamental domain and $\mu$ is right invariant.
\end{proof}

\begin{theorem}[Rallis inner product]\label{pf:tlift:rallisthm}
Let $t_0\in\R$ be such that $g\mapsto\tliftht(\iota(g,1))^t$ is $\mu$-integrable for all $t\ge t_0$, and suppose
$t_0<d/2$. Let $\nu$ be a Haar measure on $\A^\kappa$, $\Phi\in\tliftPD(M)$, and $\varphi$ continuous and left
$\Uni(X)(F)$-invariant. Then $g\mapsto\tliftZ_\Phi(g)\overline{\tliftM_\varphi(g)}$ is $\mu$-integrable, and
\[
  \tliftlift_M(\Phi,\varphi)=0\quad\Longrightarrow\quad\int_{\Uni(X)(\A)}\tliftZ_\Phi(g)\,\overline{\tliftM_\varphi(g)}\,d\mu(g)=0.
\]
\formallabel{UpperTheory.TP.rallis_k}
\end{theorem}

The integral is the doubling zeta integral at the point $s_0=\frac{d-n}2$, and the argument is the
classical one \cite{Rallis1984,PiatetskiShapiroRallis1987,Li1992}: by the Siegel--Weil formula the value of
the doubling integral of the Eisenstein series at $s_0$ is a multiple of the Petersson norm of the theta
lift; by unfolding it is the integral above.

\begin{proof}
\emph{Step 1: data.} By \cref{pf:tlift:cpt} for $W$ there are a Haar measure $\mu_W$ on $\Uni(W)(\A)$ and a
fundamental domain $D_W$ of finite measure. By \cref{pf:tlift:hgtex} the height $\tliftht$ is continuous and
positive. Let $(e^{(2)},\omega_2,S,\ell)$ be doubling data as in \cref{pf:tlift:dblx} and $c'>0$ the constant with
$\ell(\Psi_1\boxtimes\bar\Psi_2)=c'\int\Psi_1\bar\Psi_2\,d\nu$. By \cref{pf:tlift:basic}, $\Phi\in\tliftS(\A^\kappa)$. Let
$s_0=\frac{d-n}2$ and $f(x)=\ell(\omega_2(x)(\Phi\boxtimes\bar\Phi))$; by \cref{pf:tlift:sw}, $f$ is a section of exponent
$s_0$ for $(\tliftK,\chi_W)$ and $E(\cdot,\cdot;f)$ has a continuation $\tliftE$ on an open set $\mathcal U\ni s_0$ satisfying
\eqref{eq:tlift:sw}. Finally let $c\ne0$ be the constant with $\int F(g^{-1})\,d\mu=c\int F\,d\mu$ for integrable
$F$ (\lref{GlobalTheory.ThK.Dbl.inv_scale}).

\emph{Step 2: the section on the embedded group.} For $g\in\Uni(X)(\A)$,
\begin{equation}\label{eq:tlift:idf}
  f(\iota(1,g))\,\chi_W(\det g)^{-1}=c'\,\tliftZ_\Phi(g^{-1}).
\end{equation}
Indeed $\iota(1,g)=\iota(g,g)\,\iota(g^{-1},1)$ with $\iota(g,g)\in P(\A)$ of Levi component $g$ and
$|\det g|_{\tliftAE}=1$, so the section law gives $f(\iota(1,g))=\chi_W(\det g)f(\iota(g^{-1},1))$; and by
\eqref{eq:tlift:dblact} with $g'=1$, $\omega_2(\iota(g^{-1},1))(\Phi\boxtimes\bar\Phi)=\omega(g^{-1},1)\Phi\boxtimes\bar\Phi$, whose
image under $\ell$ is $c'\tliftZ_\Phi(g^{-1})$. In particular $\tliftZ_\Phi$ is continuous.

\emph{Step 3: bounds.} Since $\tliftK$ is compact and $f$ continuous, $|f(pk)|=|\det t_p|_{\tliftAE}^{d/2}|f(k)|$ gives
$|f(x)|\le C_f\,\tliftht(x)^{d/2}$. With \eqref{eq:tlift:idf} and
$\tliftht(\iota(1,g^{-1}))=\tliftht(\iota(g,1))$ (\eqref{eq:tlift:swap}) this gives
$|\tliftZ_\Phi(g)|\le c'^{-1}C_f\,\tliftht(\iota(g,1))^{d/2}$. By \cref{pf:tlift:mc}, $|\tliftM_\varphi|\le C_m$. Hence
\begin{equation}\label{eq:tlift:zbound}
  |\tliftZ_\Phi(g)\,\overline{\tliftM_\varphi(g)}|\le C\,\tliftht(\iota(g,1))^{d/2}.
\end{equation}
For $z\in\C$ let $I(z)=\int\tliftht(\iota(g,1))^{z-s_0}\tliftZ_\Phi(g)\overline{\tliftM_\varphi(g)}\,d\mu(g)$. By
\eqref{eq:tlift:zbound} the integrand is bounded by $C\tliftht(\iota(g,1))^{\Re z-s_0+d/2}$, which is integrable as soon as
$\Re z-s_0+d/2\ge t_0$. Hence $I$ is defined on the half plane $H=\{\Re z>t_0+s_0-d/2\}$, and it is holomorphic
there: on a disc of radius $\epsilon$ around a point of $H$ the integrand is dominated by
$C(\tliftht^{u-\epsilon}+\tliftht^{u+\epsilon})$ with $u-\epsilon\ge t_0$, and holomorphy follows by differentiation under the
integral. Since $t_0<d/2$, $s_0\in H$, and $I(s_0)$ is the integral
of the theorem; this proves the integrability assertion.

\emph{Step 4: unfolding.} For $s\in\mathcal U$ let
\[
  \mathcal Z(s)=\int_{D\times D}\tliftE\bigl(s,\iota(y_1,y_2)\bigr)\,\chi_W(\det y_2)^{-1}\,\varphi(y_2)\,\overline{\varphi(y_1)}\,d\mu(y_1)\,d\mu(y_2).
\]
The function $\tliftE(s,\iota(\cdot,\cdot))$ is left invariant under $\Uni(X)(F)^2$ (it is so on a half plane, where it is
given by the Eisenstein series, and the identity extends by analytic continuation) and continuous,
hence bounded on $D\times D$ locally uniformly in $s$, since $\Uni(X)(\A)=\Uni(X)(F)\Omega$ with $\Omega$ compact; so
$\mathcal Z$ is holomorphic on $\mathcal U$. By \cref{pf:tlift:hgtint} and \eqref{eq:tlift:swap}, $g\mapsto\tliftht(\iota(1,g))^t$ is
integrable for $t\ge t_0$, so for $\Re s$ large $g\mapsto f_s(\iota(1,g))$ is integrable ($|f_s|\le C_f\tliftht^{\Re s-s_0+d/2}$),
and the unfolding \cref{pf:tlift:unfold}, applied to $f_s$ (whose law has exponent $s-s_0+\frac d2$) and
integrated against $\overline{\varphi(y_1)}$ over $y_1=h\in D$ (Fubini), gives for $\Re s$ larger than some $t_1$
\[
  \mathcal Z(s)=\int_{\Uni(X)(\A)}f_s(\iota(1,g))\,\chi_W(\det g)^{-1}\,\tliftM_\varphi(g)\,d\mu(g).
\]
By \eqref{eq:tlift:idf}, $\tliftht(\iota(1,g))=\tliftht(\iota(g^{-1},1))$ and $\tliftM_\varphi(g)=\overline{\tliftM_\varphi(g^{-1})}$, the
integrand equals $c'\,\tliftht(\iota(g^{-1},1))^{s-s_0}\tliftZ_\Phi(g^{-1})\overline{\tliftM_\varphi(g^{-1})}$, and the substitution
$g\mapsto g^{-1}$ gives
\begin{equation}\label{eq:tlift:zeta}
  \mathcal Z(s)=c\,c'\,I(s)\qquad(\Re s>t_1).
\end{equation}

\emph{Step 5: analytic continuation.} The function $\mathcal Z-cc'I$ is holomorphic on $\mathcal U\cap H$ and vanishes
for $\Re s>t_1$. Since $\mathcal U$ is open and contains a punctured neighbourhood of every point of $\C$, its
complement is closed and discrete, so $\mathcal U\cap H$ is connected; by the identity theorem
$\mathcal Z-cc'I$ vanishes on $\mathcal U\cap H$, in particular at $s_0$:
$\mathcal Z(s_0)=cc'I(s_0)$.

\emph{Step 6: Siegel--Weil.} By \eqref{eq:tlift:sw} and \eqref{eq:tlift:dbltheta},
\[
  \tliftE\bigl(s_0,\iota(y_1,y_2)\bigr)\chi_W(\det y_2)^{-1}
  =\frac{\kappa}{\mu_W(D_W)}\int_{D_W}\theta(y_1,h;\Phi)\,\overline{\theta(y_2,h;\Phi)}\,d\mu_W(h).
\]
The theta series is bounded on $\Uni(X)(\A)\times\Uni(W)(\A)$ (it is continuous and invariant under the rational
points of both compact quotients), so Fubini's theorem applies and
\begin{align*}
  \mathcal Z(s_0)&=\frac{\kappa}{\mu_W(D_W)}\int_{D_W}\Bigl(\int_D\theta(y_1,h;\Phi)\overline{\varphi(y_1)}\,d\mu(y_1)\Bigr)
  \Bigl(\int_D\varphi(y_2)\overline{\theta(y_2,h;\Phi)}\,d\mu(y_2)\Bigr)d\mu_W(h)\\
  &=\frac{\kappa}{\mu_W(D_W)}\int_{D_W}\bigl|\tliftlift_M(\Phi,\varphi)(h)\bigr|^2d\mu_W(h)
\end{align*}
(\lref{GlobalTheory.ThK.Dbl.An.at_value}). If $\tliftlift_M(\Phi,\varphi)=0$, then $\mathcal Z(s_0)=0$, and since $c,c'\ne0$,
$I(s_0)=0$.
\end{proof}


% ---------------------------------------------------------------------------
\subsection{Non-vanishing of the doubling integral}\label{pf:tlift:nonvan}

We keep the assumptions of \cref{pf:tlift:rallis}; in addition $E$ is the CM field of the statement, so
$F=E^+$ and the conjugation is the complex conjugation.

\subsubsection*{The standard groups}
\begin{lemma}[Identification with the standard group]\label{pf:tlift:brex}
Let $Y$ be a space of dimension $m$ with $J_Y=\delta\cdot1_m$, $\delta\in E^\times$. There is an isomorphism of
topological groups $e_Y\colon\Uni_m(\A)\to\Uni(Y)(\A)$ which maps $\Uni_m(E)$ onto $\Uni(Y)(F)$, $\Uni_m(E_\infty)$ onto
the subgroup generated by the elements supported at infinity, $\Uni_m(\A_f)$ onto the subgroup of elements with
trivial archimedean components, and, for every finite place $w$ of $E$ above $v$, $\Uni_m(F_v)$ onto the subgroup
generated by the elements supported at $v$; and which satisfies $\tau_w(e_Y(g))=\tau_w(g)$ for the embeddings
$\tau_w$ and $\pi_t(e_Y(g))=c_{w_t}(g)$ (under $F_v\cong E_{w_t}$) at the split places.
\formallabel{UpperTheory.TP.brEx}
\end{lemma}

\begin{proof}
Since $\delta\ne0$ is a scalar, ${}^tg\,\delta\,\bar g=\delta$ if and only if ${}^tg\,\bar g=1$, so $\Uni(Y)=\Uni_{1_m}$ as subgroups of
$\GL_m$. The map $e_Y$ is induced by the canonical isomorphism of topological rings between the adele ring
$\A_E$ of $E$ and $\A\otimes_FE$, which is compatible with the conjugations and with the decompositions into
archimedean, finite and local components; the listed compatibilities follow from these.
\end{proof}

We fix such isomorphisms $e=e_X$ and $e'=e_W$ (with $\delta=1$ for $X$). For a function $\varphi$ on $\Uni_n(\A)$ put
$\varphi^\flat=\varphi\circ e^{-1}$, a function on $\Uni(X)(\A)$.

\begin{definition}[Matching types]\label{pf:tlift:sametype}
Let $U$ be an automorphic representation of $\Uni_n$, $f_0\in\tliftS(F_\infty^\kappa)$ and $\varphi\in U$. We say that
$f_0$ and $\varphi$ \emph{have the same type} if there are a finite-dimensional irreducible representation
$(V,\rho)$ of $\Uni_n(E_\infty)$, linear maps $i\colon V\to\tliftS(F_\infty^\kappa)$ and $j\colon V\to U$ and $x_0\in V$ with
$i(x_0)=f_0$, $j(x_0)=\varphi$, $j(\rho(k)x)=R(k)j(x)$, and such that $\omega(e(k),1)$ acts at infinity by
$i(x)\mapsto i(\rho(k)x)$, for all $k\in\Uni_n(E_\infty)$ and $x\in V$.
\formal{UpperTheory.TP.SameType}
\end{definition}

\subsubsection*{Decomposition of the adelic group}
For a finite place $v$ let $\Uni(X)(F_v)$ be the subgroup of $\Uni(X)(\A)$ generated by the elements supported
at $v$, and $K_v$ the set of $g$ such that $g$ and $g^{-1}$ map the $v$-component of $\widehat L_X$ into itself
(\emph{integral at $v$}). Let $\tliftInt=\bigcap_vK_v$. For a finite set $S$ of finite places let $\tliftGinf$ be
the subgroup generated by the elements supported at infinity, $\tliftGS$ the subgroup generated by the
$\Uni(X)(F_v)$, $v\in S$, and $\tliftGoff$ the subgroup generated by the elements with trivial components at
infinity and at the places of $S$.

\begin{lemma}[Decomposition]\label{pf:tlift:decomp}
Let $S$ be a finite set of finite places.
\begin{enumerate}[label=(\alph*),nosep]
\item Multiplication $\tliftGinf\times\tliftGS\times\tliftGoff\to\Uni(X)(\A)$ is an isomorphism of topological groups;
  in particular the three subgroups commute with each other.
\item $\tliftGinf$ is compact, $\tliftGS$ and $\tliftGoff$ are closed, $\tliftInt$ is compact and open, and
  $\tliftGinf\subset\tliftInt$.
\item For $b\in\tliftGS$ there are $c_v\in\Uni(X)(F_v)$ ($v\in S$) such that $c_v^{-1}b$ has trivial component at $v$, and
  every subgroup containing all $c_v$, $v\in S$, contains $b$.
\item Every $g\in\Uni(X)(\A)$ lies in $K_v$ for almost all $v$.
\end{enumerate}
\formallabel{UpperTheory.TP.decomp_k}
\end{lemma}

\begin{proof}
Let $\epsilon_\infty\in\A$ be the idempotent with archimedean components $1$ and finite components $0$, $\epsilon_S$
the idempotent which is $1$ exactly at the places of $S$, and $\epsilon^S=1-\epsilon_\infty-\epsilon_S$; these are orthogonal
idempotents with sum $1$. An element $g$ lies in $\tliftGinf$, $\tliftGS$, $\tliftGoff$ if and only if
$\epsilon(g-1)=g-1$ for $\epsilon=\epsilon_\infty,\epsilon_S,\epsilon^S$ respectively; these are closed conditions. The map
$g\mapsto\bigl(1+\epsilon_\infty(g-1),\,1+\epsilon_S(g-1),\,1+\epsilon^S(g-1)\bigr)$ is a continuous inverse of the multiplication map,
because $(1+\epsilon(g-1))(1+\epsilon'(g'-1))=1+\epsilon(g-1)+\epsilon'(g'-1)$ for orthogonal $\epsilon,\epsilon'$, and each component
is again unitary since the $\epsilon$ are fixed by the conjugation. This proves (a) and the closedness. The set
$\tliftInt$ is the intersection of $\Uni(X)(\A)$ with the standard maximal compact subgroup of
$\GL_n$ stabilising $\widehat L_X$; it is compact because $\Uni(X)(F_\infty)$ is compact for positive definite $X$
and $\widehat L_X$ is compact, and open because the stabiliser of a lattice is open. Elements supported at infinity act trivially on $\widehat L_X$,
so $\tliftGinf\subset\tliftInt$, and $\tliftGinf$ is a closed subset of the compact set $\tliftInt$. For (c) take
$c_v=1+\epsilon_v(b-1)$, $\epsilon_v$ the idempotent of $v$; then $b=\prod_{v\in S}c_v$. For (d), the matrices $g$ and $g^{-1}$
have integral entries at almost all places.
\end{proof}

\subsubsection*{The tail estimate}
\begin{lemma}[Tail]\label{pf:tlift:tail}
Let $t\in\R$ be such that $g\mapsto\tliftht(\iota(g,1))^t$ is $\mu$-integrable. For every finite set $S_0$ of finite places
there is a finite set $S\supseteq S_0$ such that for every Haar measure $\mu_C$ on $\tliftGoff$
\begin{equation}\label{eq:tlift:tail}
  \int_{\tliftGoff\setminus\tliftInt}\tliftht(\iota(c,1))^t\,d\mu_C(c)<\mu_C\bigl(\tliftGoff\cap\tliftInt\bigr).
\end{equation}
\formallabel{UpperTheory.TP.tail_k}
\end{lemma}

\begin{proof}
For $u\in\tliftInt$ we have $\iota(u,1)\in\tliftK$ (the Cartan involutions of $X$ are $1$, and $u$ preserves $\widehat L_X$), hence
$\tliftht(\iota(gu,1))=\tliftht(\iota(g,1)\iota(u,1))=\tliftht(\iota(g,1))$. Thus $f(g)=\tliftht(\iota(g,1))^t$ is right $\tliftInt$-invariant,
hence locally constant ($\tliftInt$ is open), and $m=f\mu$ is a finite measure. For finite $S$ let
$A(S)=\{g:g\in K_v\text{ for }v\in S\}\setminus\tliftInt$. The sets $A(S)$ are measurable, decrease as $S$ grows, and
$\bigcap_SA(S)=\emptyset$ (an element of all of them would lie in every $K_v$, hence in $\tliftInt$). Hence
$m(A(S))\to0$ along the directed set of finite $S$, and since $\mu(\tliftInt)>0$ ($\tliftInt$ is open and contains $1$)
there is $S\supseteq S_0$ with $m(A(S))<\mu(\tliftInt)$.

Let $\mathcal D=\tliftGinf\tliftGS\cong\tliftGinf\times\tliftGS$, so $\Uni(X)(\A)=\mathcal D\times\tliftGoff$ by \cref{pf:tlift:decomp}, and choose Haar
measures $\mu_{\mathcal D}$ on $\mathcal D$ and $\mu_C$ on $\tliftGoff$; then $\mu_{\mathcal D}\times\mu_C$ is a multiple of $\mu$. Let
$T=\mathcal D\cap\tliftInt$, an open set of positive and finite measure. For $x\in\mathcal D$ and $c\in\tliftGoff$: $xc\in\tliftInt$ if and
only if $x\in\tliftInt$ and $c\in\tliftInt$, since integrality is a condition on the components at each finite place and
$x$, $c$ live at disjoint sets of places (the components of $x$ at places outside $S$ and of $c$ at places of $S$
are trivial). If $x\in T$ and $c\in\tliftGoff\setminus\tliftInt$, then $xc\in A(S)$ and $f(xc)=f(cx)=f(c)$. Hence
\[
  \mu_{\mathcal D}(T)\int_{\tliftGoff\setminus\tliftInt}f\,d\mu_C\le(\mu_{\mathcal D}\times\mu_C)\text{-}\!\!\int_{A(S)}f
  <(\mu_{\mathcal D}\times\mu_C)(\tliftInt)=\mu_{\mathcal D}(T)\,\mu_C(\tliftGoff\cap\tliftInt),
\]
where the strict inequality is $m(A(S))<\mu(\tliftInt)$ rescaled. Dividing by $\mu_{\mathcal D}(T)\in(0,\infty)$ gives
\eqref{eq:tlift:tail}.
\end{proof}

\subsubsection*{An abstract non-vanishing lemma}
\begin{lemma}[Abstract non-vanishing]\label{pf:tlift:absne}
Let $(\mathcal A,\mu_A)$, $(\mathcal B,\mu_B)$, $(\mathcal C,\mu_C)$ be $s$-finite measure spaces, $H\subset\mathcal B$ measurable with
$0<\mu_B(H)<\infty$, and $K\subset\mathcal C$ measurable with $\mu_C(K)<\infty$. Let $f\colon\mathcal A\times\mathcal B\times\mathcal C\to\C$ be integrable
for $\mu_A\times\mu_B\times\mu_C$, and let $m\colon\mathcal A\times\mathcal C\to\C$, $s,c_0\colon\mathcal C\to\C$, $w\colon\mathcal C\to[0,\infty]$,
$\kappa\in\C^\times$ and $r>0$ satisfy:
\begin{enumerate}[label=(\roman*),nosep]
\item $f(x,b,c)=0$ for $b\notin H$, and $f(x,b,c)=s(c)\,m(x,c)$ for $b\in H$;
\item $\int m(x,c)\,d\mu_A(x)=\kappa\,c_0(c)$ for every $c$;
\item $s(c)=1$ and $c_0(c)=r$ for $c\in K$; $|c_0(c)|\le r$ and $|s(c)|\le w(c)$ for all $c$;
\item $\int_{\mathcal C\setminus K}w\,d\mu_C<\mu_C(K)$.
\end{enumerate}
Then $\int f\,d(\mu_A\times\mu_B\times\mu_C)\ne0$.
\formallabel{UpperTheory.TP.absNe}
\end{lemma}

\begin{proof}
By Fubini and (i), (ii), the integral of $f$ over $x$ is $\one_H(b)\,\kappa\,s(c)c_0(c)$, and $c\mapsto s(c)c_0(c)$ is
integrable (otherwise $\mu_B$-almost every $b$ would lie outside $H$, contradicting $\mu_B(H)>0$). Hence
$\int f=\mu_B(H)\,\kappa\int_{\mathcal C}s\,c_0\,d\mu_C$, and it suffices to show $\int_{\mathcal C}s\,c_0\ne0$. By (iii),
$\int_Ks\,c_0=r\mu_C(K)$, while $|\int_{\mathcal C\setminus K}s\,c_0|\le r\int_{\mathcal C\setminus K}w<r\mu_C(K)$ by (iii), (iv).
\end{proof}

\subsubsection*{Bump vectors}
\begin{definition}[Bump vectors]\label{pf:tlift:bump}
Let $v$ be a finite place. A pair $(H,\phi)$ consisting of a subgroup $H\subset\Uni(X)(F_v)$ and
$\phi\in\tliftS(F_v^\kappa)$ is a \emph{bump} at $v$ if: $\phi\ne0$; $H\subset K_v$; $H$ contains $\Uni(X)(F_v)\cap O$ for some
open neighbourhood $O$ of $1$; $\omega(g,1)$ acts at $v$ trivially on $\phi$ for $g\in H$; and for every
$g\in\Uni(X)(F_v)\setminus H$ there are $u$ supported at $v$ and $c\ne c'$ such that $\rho(u)$ and $\rho(gu)$ act at $v$ on
$\phi$ by the scalars $c$ and $c'$ respectively.
\formal{UpperTheory.TP.IsBump}
\end{definition}

The last condition makes the matrix coefficient vanish off $H$:

\begin{lemma}[Orthogonality]\label{pf:tlift:orth}
Let $\Phi\in\tliftS(\A^\kappa)$, $\Phi\ne0$, $g\in\Uni(X)(\A)$ and $u$ with $\rho(u)\Phi=c\Phi$ and $\rho(gu)\Phi=c'\Phi$, $c\ne c'$. Then
$\tliftZ_\Phi(g)=0$.
\formallabel{GlobalTheory.ThK.Dbl.Bad.orth_k}
\end{lemma}

\begin{proof}
By (W1), $\rho(gu)\,\omega(g,1)\Phi=\omega(g,1)\rho(u)\Phi=c\,\omega(g,1)\Phi$. The operator $\rho(gu)$ is unitary for the
$L^2$-norm of $\nu$ (as $|\psi|=1$ and $\nu$ is translation invariant), so
$\langle\rho(gu)A,\rho(gu)B\rangle=\langle A,B\rangle$ for $A=\omega(g,1)\Phi$, $B=\Phi$ gives $c\,\bar c'\,\tliftZ_\Phi(g)=\tliftZ_\Phi(g)$, and for
$A=B=\Phi$ it gives $|c'|^2=1$ (as $\langle\Phi,\Phi\rangle>0$). If $\tliftZ_\Phi(g)\ne0$ then $c\bar c'=1=c'\bar c'$, so $c=c'$.
\end{proof}

\begin{lemma}[Lattice data]\label{pf:tlift:latdatum}
Assume $2n\le d$, and let $v$ be a finite place. There are:
\begin{enumerate}[label=(\roman*),nosep]
\item an additive group $\Lambda$ of $n\times n$ matrices supported at $v$ which is closed under multiplication, contains
  $1_v$ (the idempotent of $v$ times the identity) and is stable under multiplication by integral adeles supported at
  $v$, such that $a\gamma\in\Lambda_v$ for $a\in\Lambda$ and $\gamma\in\Lambda_v$ ($\Lambda_v$ as in (M2));
\item $x_1\in\Lambda_v$ with $x_1J_W{}^t\bar x_1=0$, and a matrix $r$ with $x_1r=1_v$ and $\gamma r\in\Lambda$ for all $\gamma\in\Lambda_v$;
\item $\varpi\in\OO_v$ with $\varpi\ne0$ and $|\varpi|<1$ such that $1_v+\varpi p$ has a right inverse in $\Lambda$ for every $p\in\Lambda$;
\end{enumerate}
such that the sets $H_k=\{g\in\Uni(X)(F_v):g-1\in\varpi^k\Lambda,\ g^{-1}-1\in\varpi^k\Lambda\}$ satisfy: every neighbourhood of $1$ in
$\Uni(X)(\A)$ contains $H_k$ for all large $k$, and for every $k$ there is an open $O\ni1$ with $\Uni(X)(F_v)\cap O\subset H_k$.
\formallabel{UpperTheory.exists_weilLatticeDatum}
\end{lemma}

\begin{proof}
\emph{An isotropic splitting of $W$ at $v$.} Since $E/F$ is unramified at $v$, the place $v$ is split or inert. Because $d\ge2n$,
there are matrices $Y\in M_{n\times d}(\tliftAE)$ and $R\in M_{d\times n}(\tliftAE)$, supported at $v$, such that the rows of $Y$ lie in
the local lattice $L_{W,v}$ of $W$ at $v$, $YR=1_v$, $YJ_W{}^t\bar Y=0$, and $zR$ has integral coordinates for every $z\in L_{W,v}$
(\lref{UpperTheory.exists_isotropic_splitting_of_two_mul_le} at split places,
\lref{UpperTheory.exists_isotropic_splitting_at_inert_prime} at inert places): the rows of $Y$ span an
$n$-dimensional isotropic subspace of $W\otimes_FF_v$ which is a direct summand of the lattice.

\emph{A frame of $X$ at $v$.} Since $\OO_E\otimes\OO_v$ is a product of discrete valuation rings, the local lattice $L_{X,v}$ is
free; so there are matrices $b,c$ supported at $v$ with $bc=cb=1_v$, the columns of $b$ in $L_{X,v}$, and $cz$ integral for
$z\in L_{X,v}$. Let $\varpi$ be a uniformiser of $F_v$.

\emph{Definition and properties.} Let $\Lambda$ be the set of matrices $a$ supported at $v$ such that $cab$ is integral, and put
$x_1=bY$, $r=Rc$. Then: $c(aa')b=(cab)(ca'b)$ because $bc=1_v$, so $\Lambda$ is closed under multiplication; $c1_vb=1_v$; and
$\Lambda$ is stable under integral scalars. For $a\in\Lambda$ and $z\in L_{X,v}$, $az=b\,(cab)(cz)\in L_{X,v}$; applied to the generators
$z\,{}^ty$ of $\Lambda_v$ this gives (i). Next, $x_1=\sum_kb_{\cdot k}Y_{k\cdot}\in\Lambda_v$, $x_1J_W{}^t\bar x_1=b\,(YJ_W{}^t\bar Y)\,{}^t\bar b=0$,
$x_1r=b(YR)c=bc=1_v$, and for a generator $\gamma=z\,{}^ty$ of $\Lambda_v$, $c\gamma rb=(cz)({}^tyR)(cb)$ is integral; this gives (ii). Under
$a\mapsto cab$, $\Lambda$ is identified with the ring of integral $n\times n$ matrices over $\OO_E\otimes\OO_v$, in which
$1+\varpi p$ is invertible (its determinant is congruent to $1$ modulo $\varpi$); this gives (iii). Finally, under the same
identification the sets $\varpi^k\Lambda$ form a basis of neighbourhoods of $0$ in the matrices supported at $v$, and the
topology of $\Uni(X)(F_v)$ is that of matrices together with their inverses; this gives the last assertion.
\end{proof}

\begin{lemma}[Separation]\label{pf:tlift:separate}
Let $v$ be a finite place, $g\in\Uni(X)(\A)$ supported at $v$, and $y\in M_{n\times d}(\tliftAE)$ supported at $v$ with
$y\notin\Lambda_v+g^{-1}\Lambda_v$. Then there is $u\in\Lambda_v$ with $gu\in\Lambda_v$ and $\psi(\langle u,y\rangle)\ne1$.
\formallabel{UpperTheory.exists_separating_of_not_mem_lattice_coset}
\end{lemma}

\begin{proof}
The lattice $\Lambda_v$ is self-dual: a matrix $z$ supported at $v$ lies in $\Lambda_v$ if and only if $\psi(\langle z,\gamma\rangle)=1$ for all
$\gamma\in\Lambda_v$ (\lref{GlobalTheory.LatSelfDualAll.latAtV_selfDual_k}; this uses the self-duality of $L_X$ and
$L_W$). The Darboux coordinates at $v$ identify the matrices supported at $v$ with $F_v^\kappa\times F_v^\kappa$, and
$S=\Lambda_v+g^{-1}\Lambda_v$ with an $\OO_v$-submodule which contains $\Lambda_v$ and is bounded, since $\Lambda_v$ and $g^{-1}\Lambda_v$
are. Since the locally compact group $F_v^\kappa\times F_v^\kappa$ is identified with its Pontryagin dual by
$(w,s)\mapsto\psi(\langle w,s\rangle)$, the closed subgroup $S$ is the annihilator of its annihilator: if $y\notin S$, there is $w$ with
$\psi(\langle w,s\rangle)=1$ for all $s\in S$ and $\psi(\langle w,y\rangle)\ne1$. Then $u=w$ works:
$u\in\Lambda_v$ by self-duality, since $\Lambda_v\subset S$; and for $\gamma\in\Lambda_v$,
$\psi(\langle gu,\gamma\rangle)=\psi(\langle u,g^{-1}\gamma\rangle)=1$ because $g^{-1}\gamma\in S$ and the form is invariant, so $gu\in\Lambda_v$.
\end{proof}

\begin{proposition}[Existence of bumps]\label{pf:tlift:bumpany}
Assume $2n\le d$. For every finite place $v$ and every neighbourhood $N$ of $1$ in $\Uni(X)(\A)$ there is a bump
$(H,\phi)$ at $v$ with $H\subset N$.
\formallabel{UpperTheory.TP.bumpAny_k}
\end{proposition}

\begin{proof}
\emph{Lattice datum.} Let $(\Lambda,x_1,r,\varpi)$ be as in \cref{pf:tlift:latdatum} (this is where $2n\le d$ is used), and $H_k$ the
congruence subgroups defined there.

\emph{Construction.} By \cref{pf:tlift:acts} there is a neighbourhood $U_1$ of $1$ such that $\omega(g,1)$ fixes
$\tliftvac_v$ at $v$ for $g\in U_1$ supported at $v$; and $K_v$ is open. Choose $k_0$ with $H_k\subset N\cap U_1\cap K_v$ for
$k\ge k_0$, let $k=k_0+1$, $H=H_k$, $x_0=\varpi^{-k}x_1$ (supported at $v$) and $\phi$ the vector by which $\rho(x_0)$ acts
on $\tliftvac_v$; $\phi\ne0$ since $\rho(x_0)$ is invertible. $H$ is a subgroup ($(1+\varpi^ka)(1+\varpi^kb)=1+\varpi^k(a+b+\varpi^kab)$), it
contains $\Uni(X)(F_v)\cap O$ for an open $O\ni1$ since the $H_k$ are neighbourhoods of $1$, and $H\subset N\cap K_v$.

\emph{$H$ fixes $\phi$.} Let $g\in H$, $g-1=\varpi^ka$ with $a\in\Lambda$. Then $gx_0=x_0+ax_1$ with $ax_1\in\Lambda_v$. By (M2),
$\rho(ax_1)$ acts on $\tliftvac_v$ by $\psi(\frac12\Tr(\vartheta h(ax_1,ax_1)))=1$, because $h(ax_1,ax_1)$ is a multiple of
$x_1J_W{}^t\bar x_1=0$; for the same reason $\langle x_0,ax_1\rangle=0$, so $\rho(x_0+ax_1)=\rho(x_0)\rho(ax_1)$. As $\omega(g,1)$
fixes $\tliftvac_v$ (since $g\in U_1$), (W1) gives
\[
  \omega(g,1)\phi=\omega(g,1)\rho(x_0)\tliftvac_v=\rho(gx_0)\,\omega(g,1)\tliftvac_v=\rho(x_0)\rho(ax_1)\tliftvac_v=\rho(x_0)\tliftvac_v=\phi
\]
(all at $v$).

\emph{Eigenvectors.} For $\gamma\in\Lambda_v$, the Heisenberg relations and (M2) give that $\rho(\gamma)$ acts on $\phi$ by the
scalar $\psi(\langle\gamma,x_0\rangle)\,Q(\gamma)$, $Q(\gamma)=\psi(\frac12\Tr(\vartheta h(\gamma,\gamma)))$.

\emph{Separation.} Let $g\in\Uni(X)(F_v)\setminus H$ and $y=g^{-1}x_0-x_0$. If $y=a+g^{-1}b$ with $a,b\in\Lambda_v$, then using
$x_1r=1_v$, $\gamma r\in\Lambda$ and the invertibility of $1_v+\varpi p$ in $\Lambda$ one solves for $g-1$ and $g^{-1}-1$ in $\varpi^k\Lambda$,
i.e.\ $g\in H$ (an explicit matrix identity, \lref{UpperTheory.TP.BA.cong_key_pk}); this is excluded. Hence
$y\notin\Lambda_v+g^{-1}\Lambda_v$, and by \cref{pf:tlift:separate} there is $u\in\Lambda_v$ with $gu\in\Lambda_v$ and $\psi(\langle u,y\rangle)\ne1$.
Now $Q(gu)=Q(u)$ and
$\langle gu,x_0\rangle=\langle u,g^{-1}x_0\rangle$ by invariance, so the eigenvalues $c=\psi(\langle u,x_0\rangle)Q(u)$ and
$c'=\psi(\langle gu,x_0\rangle)Q(gu)$ of $\rho(u)$ and $\rho(gu)$ on $\phi$ satisfy $c'/c=\psi(\langle u,y\rangle)\ne1$.
\end{proof}

\subsubsection*{The doubling integral on a bump datum}
\begin{proposition}[Bump data]\label{pf:tlift:datum}
Let $S$ be a finite set of finite places such that $\chi_W$ is unramified at every place of $E$ above every
$v\notin S$. Let $\phi$ be finite data with bumps $(H_v,\phi_v)$ at the places $v\in S$ and $\phi_v=\tliftvac_v$ for $v\notin S$,
let $\mathcal H$ be the subgroup generated by the $H_v$, $v\in S$, and let $0\ne f_0\in\tliftPol$ and $\Phi=f_0\otimes\phi$. Then
$\tliftZ_\Phi(1)\ne0$, and for $x\in\tliftGinf$, $b\in\tliftGS$, $c\in\tliftGoff$:
\begin{enumerate}[label=(\alph*),nosep]
\item if $b\notin\mathcal H$ then $\tliftZ_\Phi(xbc)=0$;
\item there is $s=s(c)\in\C$ with $|s|=\tliftht(\iota(c,1))^{d/2}$, $s=1$ if $c\in\tliftInt$, and
  $\tliftZ_\Phi(x'b'c)=s\,\tliftZ_\Phi(x')$ for all $x'\in\tliftGinf$ and $b'\in\mathcal H$.
\end{enumerate}
\formallabel{UpperTheory.TP.datum_k}
\end{proposition}

\begin{proof}
$\Phi$ is a nonzero Schwartz--Bruhat function, so $\tliftZ_\Phi(1)=\int|\Phi|^2d\nu>0$.

(a) Write $b=\prod_{v\in S}c_v$ as in \cref{pf:tlift:decomp}(c). If $b\notin\mathcal H$, some $c_v$ is not in $H_v$, and the bump
property gives $u$ supported at $v$ and $c\ne c'$ with $\rho(u)\phi_v=c\phi_v$, $\rho(c_vu)\phi_v=c'\phi_v$. Since $x$, $c$ and the
$c_{v'}$, $v'\ne v$, have trivial component at $v$, $(xbc)u=c_vu$; as $\rho(u)$ and $\rho(c_vu)$ act only at $v$,
$\rho(u)\Phi=c\Phi$ and $\rho((xbc)u)\Phi=c'\Phi$, and \cref{pf:tlift:orth} gives $\tliftZ_\Phi(xbc)=0$.

(b) Let $(e^{(2)},\omega_2,S_W,\ell)$ be doubling data as in \cref{pf:tlift:dblx}, with $\ell(\Psi_1\boxtimes\bar\Psi_2)=c_0'\int\Psi_1\bar\Psi_2\,d\nu$,
and $F(x)=\ell(\omega_2(x)(\Phi\boxtimes\bar\Phi))$. By \eqref{eq:tlift:dblact},
$F(\iota(g,1))=c_0'\,\tliftZ_\Phi(g)$. If $c\in\tliftInt$, put $p=1$, $t=1$, $k=\iota(c,1)\in\tliftK$; otherwise apply the Iwasawa
decomposition to the components of $\iota(c,1)$ outside $S$ and at infinity, obtaining $\iota(c,1)=pk$ with $p\in P(\A)$
of Levi component $t$ and $k\in\tliftK$, where $p$ and $k$ have trivial components at infinity and at the places of
$S$ (cut the decomposition of $\iota(c,1)$ with the idempotent of these places). Put $s=\chi_W(\det t)|\det t|_{\tliftAE}^{d/2}$; then
$|s|=|\det t|^{d/2}_{\tliftAE}=\tliftht(\iota(c,1))^{d/2}$, and $s=1$ if $c\in\tliftInt$.

Next, $\omega_2(k)$ fixes $\Psi=\Phi\boxtimes\bar\Phi$. Let $v\notin S$ and let $k'\in\tliftK$ with $k'-1$ supported at $v$.
Since $\chi_W$ is unramified above $v$, $\omega_2(k')$ acts at $v$ trivially on the factor
$\tliftvac_v\boxtimes\overline{\tliftvac_v}$ of $\Psi$ (\cref{pf:tlift:vacfix}), so $\omega_2(k')\Psi=\Psi$. For a finite set $T$ of
finite places let $k_T$ be the element which agrees with $k$ at the places of $T$ and is trivial elsewhere (cut
$k-1$ with the idempotent of $T$). Then $k_T\in\tliftK$, and $k_T$ is the product of the elements $k_{\{v\}}$,
$v\in T\setminus S$, each of which lies in $\tliftK$ and is supported at a single place $v\notin S$ (recall that $k$ is
trivial at infinity and at $S$); hence $\omega_2(k_T)\Psi=\Psi$. As $T$ grows, $k_T\to k$, and for every point $z$ the map
$g\mapsto(\omega_2(g)\Psi)(z)$ is continuous by (W1), the evaluation at $z$ being an admissible linear form.
Hence $\omega_2(k)\Psi=\Psi$. Let $y\in\tliftGinf\tliftGS$; then $y$ commutes with $c$ and $\iota(y,1)$
commutes with $k$, as they live at disjoint sets of places. Hence $\iota(yc,1)=\iota(cy,1)=p\,k\,\iota(y,1)=p\,\iota(y,1)\,k$ and, by
the section law (\cref{pf:tlift:law}),
\[
  c_0'\tliftZ_\Phi(yc)=F(p\,\iota(y,1)\,k)=s\,\ell\bigl(\omega_2(\iota(y,1))\omega_2(k)(\Phi\boxtimes\bar\Phi)\bigr)=s\,F(\iota(y,1))=s\,c_0'\,\tliftZ_\Phi(y).
\]
Finally, for $x'\in\tliftGinf$ and $b'\in\mathcal H$: $\omega(b',1)\Phi=\Phi$, because each $H_v$ fixes $\phi_v$ at $v$; so
$\tliftZ_\Phi(x'b'c)=s\,\tliftZ_\Phi(x'b')=s\,\langle\omega(x',1)\omega(b',1)\Phi,\Phi\rangle=s\,\tliftZ_\Phi(x')$.
\end{proof}

\subsubsection*{The archimedean integral}
\begin{proposition}[Schur orthogonality step]\label{pf:tlift:schur}
Let $U$ be an automorphic representation of $\Uni_n$, and let $0\ne f_0\in\tliftS(F_\infty^\kappa)$ and $0\ne\varphi\in U$ have
the same type (\cref{pf:tlift:sametype}). Let $\phi$ be finite data, $\Phi=f_0\otimes\phi$, and $\mu_A$ a Haar measure on the
compact group $\tliftGinf$. There is $\kappa\ne0$ such that for every $c$ with trivial archimedean components
\[
  \int_{\tliftGinf}\tliftZ_\Phi(x)\,\overline{\tliftM_{\varphi^\flat}(xc)}\,d\mu_A(x)=\kappa\,\tliftZ_\Phi(1)\,\overline{\tliftM_{\varphi^\flat}(c)}.
\]
\formallabel{UpperTheory.TP.schurStep_k}
\end{proposition}

\begin{proof}
Let $(V,\rho,i,j,x_0)$ be as in \cref{pf:tlift:sametype}, and transport $\rho$ to $\tliftGinf$ via $e$ (\cref{pf:tlift:brex}),
writing $\pi(x)=\rho(e^{-1}x)$. Since $\pi$ is irreducible and $\ker j$ is invariant, $j$ is injective ($j(x_0)=\varphi\ne0$).
For $a\in V$ let $\Psi(a)=i(a)\otimes\phi$ and $J(a)=j(a)^\flat$. Then $\omega(x,1)\Psi(a)=\Psi(\pi(x)a)$ and
$J(\pi(x)a)=R(x)J(a)$ for $x\in\tliftGinf$. Consider the sesquilinear forms
\[
  B_1(a,b)=\int\Psi(a)\,\overline{\Psi(b)}\,d\nu,\qquad B(a,b)=\int_DJ(a)\,\overline{J(b)}\,d\mu.
\]
Both are $\pi$-invariant: $B_1$ by unitarity of $\omega$ (W1), $B$ by right invariance of $\mu$ and the fact that
$Dx$ is a fundamental domain; and $B$ is
hermitian and positive definite since $J$ is injective and $J(a)$ is continuous and left invariant. By
Schur's lemma for the irreducible finite-dimensional $\pi$, $B_1=c_1B$ for some $c_1\in\C$. Fix $c$ and let $u=R(c)\varphi^\flat$; by the
Riesz representation for $B$ there is $u_0\in V$ with $\int_Du\,\overline{J(b)}\,d\mu=B(u_0,b)$ for all $b$. For $x\in\tliftGinf$ one
gets $\tliftM_{\varphi^\flat}(xc)=\int_Du(yx)\overline{J(x_0)(y)}\,d\mu(y)=B(\pi(x)u_0,x_0)$ (substitute $y\mapsto yx^{-1}$ and use
invariance), and $\tliftZ_\Phi(x)=B_1(\pi(x)x_0,x_0)$ since $\omega(x,1)\Phi=\Psi(\pi(x)x_0)$.
Schur's orthogonality relations for the compact group $\tliftGinf$ and the invariant inner product $B$ give
\[
  \int_{\tliftGinf}B(\pi(x)a,b)\,\overline{B(\pi(x)a',b')}\,d\mu_A(x)=\frac{\mu_A(\tliftGinf)}{\dim V}\,B(a,a')\,\overline{B(b,b')}.
\]
With $a=x_0$, $a'=u_0$, $b=b'=x_0$ and $B_1=c_1B$ this yields the claim with
$\kappa=\mu_A(\tliftGinf)/\dim V$, using $B(x_0,u_0)=\overline{B(u_0,x_0)}=\overline{\tliftM_{\varphi^\flat}(c)}$ and
$c_1B(x_0,x_0)=\tliftZ_\Phi(1)$.
\end{proof}

\subsubsection*{Non-vanishing}
\begin{theorem}[Non-vanishing of the lift]\label{pf:tlift:nonvanthm}
Assume $2n\le d$, and let $t_0<d/2$ be such that $g\mapsto\tliftht(\iota(g,1))^t$ is $\mu$-integrable for all $t\ge t_0$. Let
$U$ be an automorphic representation of $\Uni_n$, and $0\ne f_0\in\tliftPol$, $0\ne\varphi\in U$ of the same type. Then
there are finite data $\phi$ such that $\tliftlift_M(f_0\otimes\phi,\varphi^\flat)\ne0$.
\formallabel{UpperTheory.TP.nonVan_k}
\end{theorem}

\begin{proof}
\emph{Choices.} $\varphi^\flat$ is continuous and left $\Uni(X)(F)$-invariant, and nonzero. Let $r=\tliftM_{\varphi^\flat}(1)>0$, so that
$|\tliftM_{\varphi^\flat}|\le r$ (\cref{pf:tlift:mc}). Since $\varphi$ is an automorphic form, the $k\in\Uni_n(\A_f)$ with
$R(k)\varphi=\varphi$ form an open subset of $\Uni_n(\A_f)$ containing $1$; transporting by $e$ we obtain an open
neighbourhood $N$ of $1$ in $\Uni(X)(\A)$ such that $\varphi^\flat(yg)=\varphi^\flat(y)$ for all $y$ and all $g\in N$ with trivial
archimedean components. Since $\tliftInt$ is compact and the
groups $G^{S}$ are closed, decrease with $S$ and have intersection $\{1\}$, there is a finite $S_0$ such that every
$c\in G^{S_0}\cap\tliftInt$ lies in $N$ (the sets $(\tliftInt\setminus\mathrm{int}(N))\cap G^S$ form a directed family of closed
subsets of a compact set with empty intersection, so one of them is empty). Let $R$ be a finite set of finite places outside of which
$\chi_W$ is unramified ($\chi_W$ is continuous, hence unramified at almost all places). The function
$\tliftht(\iota(g,1))^{d/2}$ is integrable as $t_0<d/2$; let $S\supseteq S_0\cup R$ be as in \cref{pf:tlift:tail} for $t=d/2$. For
$v\in S$ choose a bump $(H_v,\phi_v)$ at $v$ with $H_v\subset N$ (\cref{pf:tlift:bumpany}), put $\phi_v=\tliftvac_v$ for $v\notin S$,
and let $\Phi=f_0\otimes\phi\in\tliftPD(M)$ and $\mathcal H$ the group generated by the $H_v$.

\emph{Suppose the lift vanishes.} By \cref{pf:tlift:rallisthm} (applicable since $t_0<d/2$),
$F(g)=\tliftZ_\Phi(g)\overline{\tliftM_{\varphi^\flat}(g)}$ is integrable and $\int F\,d\mu=0$. Choose Haar measures $\mu_A,\mu_B,\mu_C$ on
$\tliftGinf,\tliftGS,\tliftGoff$; by \cref{pf:tlift:decomp}(a), $f(x,b,c)=F(xbc)$ is integrable for $\mu_A\times\mu_B\times\mu_C$ and its
integral is a constant multiple of $\int F\,d\mu=0$ (the image of $\mu_A\times\mu_B\times\mu_C$ is a Haar measure on
$\Uni(X)(\A)$). We now verify the hypotheses of
\cref{pf:tlift:absne} with $\mathcal A=\tliftGinf$, $\mathcal B=\tliftGS$, $\mathcal C=\tliftGoff$, $H=\mathcal H$, $K=\tliftGoff\cap\tliftInt$,
\[
  m(x,c)=\tliftZ_\Phi(x)\overline{\tliftM_{\varphi^\flat}(xc)},\quad c_0(c)=\overline{\tliftM_{\varphi^\flat}(c)},\quad
  w(c)=\tliftht(\iota(c,1))^{d/2},
\]
$s(c)$ as in \cref{pf:tlift:datum}(b) and $\kappa$ replaced by $\kappa\,\tliftZ_\Phi(1)$ with $\kappa$ from \cref{pf:tlift:schur}; this
gives $\int f\ne0$, a contradiction.

\emph{The measure of $\mathcal H$.} $\mathcal H$ is open in $\tliftGS$, since it contains $\Uni(X)(F_v)\cap O_v$ for open $O_v\ni1$ and every
$b$ near $1$ in $\tliftGS$ has components $c_v$ near $1$;
so $\mu_B(\mathcal H)>0$. And $\mathcal H\subset\tliftInt\cap\tliftGS$, which is compact, so $\mu_B(\mathcal H)<\infty$. Similarly
$\mu_C(K)<\infty$.

\emph{(i).} If $b\notin\mathcal H$ then $F(xbc)=0$ by \cref{pf:tlift:datum}(a). If $b\in\mathcal H$, then
$\tliftZ_\Phi(xbc)=s(c)\tliftZ_\Phi(x)$ by \cref{pf:tlift:datum}(b); and since $b$ commutes with $c$ and $b\in\mathcal H\subset N$,
$\varphi^\flat(yxbc)=\varphi^\flat(yxcb)=\varphi^\flat(yxc)$, so $\tliftM_{\varphi^\flat}(xbc)=\tliftM_{\varphi^\flat}(xc)$. Hence $f(x,b,c)=s(c)m(x,c)$.

\emph{(ii)} is \cref{pf:tlift:schur}: $\int m(x,c)\,d\mu_A(x)=\kappa\tliftZ_\Phi(1)c_0(c)$, and $\kappa\tliftZ_\Phi(1)\ne0$ by
\cref{pf:tlift:datum}.

\emph{(iii).} For $c\in K$: $s(c)=1$ by \cref{pf:tlift:datum}(b), and $c\in G^{S_0}\cap\tliftInt\subset N$, so
$\varphi^\flat(yc)=\varphi^\flat(y)$ and $c_0(c)=\overline{\tliftM_{\varphi^\flat}(1)}=r$. For all $c$, $|c_0(c)|\le r$ and $|s(c)|=w(c)$.

\emph{(iv)} is \eqref{eq:tlift:tail} with $t=d/2$.
\end{proof}


% ---------------------------------------------------------------------------
\subsection{The archimedean pairing}\label{pf:tlift:real}

We keep the assumptions of \cref{pf:tlift:nonvan}, and assume moreover $J_X=1_n$, $J_W=\delta\cdot1_d$, $n\le d$,
and that $\chi_X,\chi_W$ have archimedean exponents $m_X,m_W$.

\subsubsection*{Polynomials and the Fock model}
Let $\tliftPP=\C[M_{n\times d}]$ be the polynomial ring in the entries $Z_{ij}$ of an $n\times d$ matrix $Z$. For $a,b\in\Z$,
$g\in\GL_n(\C)$, $h\in\GL_d(\C)$ and $P\in\tliftPP$ put
\[
  (L_a(g)P)(Z)=\det(g)^a\,P({}^tg\,Z),\qquad(R_b(h)P)(Z)=\det(h)^b\,P(Z\,h).
\]
For $1\le s\le n$ let $\Delta_s(Z)$ be the determinant of the upper left $s\times s$ block of $Z$, and for $\mu\in\N^n$ with
$\mu_1\ge\dots\ge\mu_n$ put $\Delta^\mu=\prod_{s=1}^n\Delta_s^{\mu_s-\mu_{s+1}}$ ($\mu_{n+1}=0$). The polynomial $\Delta^\mu$ is a
highest weight vector for both actions (it is the joint highest weight vector of the pair
$(\GL_n,\GL_d)$ \cite{Howe1989}).

The development provides a Fock model for the archimedean part of $\omega$: frames $A_w$ of $\tau_w(J_X)$ and $C_w$ of
$i\,\tau_w(J_W)$ adapted to the Cartan involutions (\lref{UpperTheory.TP.RP.frames}); an injective linear map
$\Phi_F\colon\bigotimes_w\tliftPP\to\tliftS(F_\infty^\kappa)$ (tensor product over the infinite places of $F$) through which
the derivatives $D_u$ act by explicit first order operators determined by $\tau_w(u)$ and the frames
(\lref{GlobalTheory.RealLift.Main.i_fockex}), and whose image is the polynomial Fock space $\tliftPol$
(\lref{GlobalTheory.RealLift.Main.i_polrange}); and the infinitesimal actions of the two unitary groups:
for a one-parameter subgroup $\gamma$ of $\Uni(X)(F_\infty)$ with generators $Y_w$ (written in the frames),
\[
  \frac{d}{dt}\Big|_{t=0}\omega(\gamma(t),1)\Phi_F(T)=\Phi_F\Bigl(\sum_w\bigl(\mathcal L(Y_w)+\tfrac{m_W(w)}2\tr Y_w\bigr)_wT\Bigr),
\]
where $\mathcal L(Y)$ is the derivation of $\tliftPP$ induced by left substitution plus $\frac d2\tr Y$ (so that
$\mathcal L(Y)1=\frac d2\tr(Y)\,1$) and $(\cdot)_w$ means acting on the $w$-th factor; and the analogous statement for
$\Uni(W)(F_\infty)$ with right substitution, $\frac n2\tr$ and $m_X$
(\lref{PositivityTheory.RL.lieFirst}, \lref{PositivityTheory.ThKR.lieSecond2}). These are the
classical Fock model formulas for the compact dual pair $(\Uni(n),\Uni(d))$ \cite{KashiwaraVergne1978,Howe1989}.
The frames have the property that $\tau_w(g)A_w=A_w\exp(Y)$ implies $\tau_w(g)=\exp(Y)$, and similarly for $C_w$
(\lref{UpperTheory.TP.RP.frames}).

\subsubsection*{From the Lie algebra to the group}
For a linear operator $L$ on a complex vector space $V$ and a linear operator $U$ we say that $U=\exp(L)$ \emph{weakly} if
$\sum_{k\ge0}\lambda(L^kv)/k!$ converges to $\lambda(Uv)$ for every $v\in V$ and every linear form $\lambda$ on $V$; $U$ is then determined by $L$.
If $L$ is \emph{locally finite} (every vector lies in a finite-dimensional $L$-stable subspace), then $\exp(L)$ exists and is
given on each finite-dimensional stable subspace by the usual exponential series. For $Y\in M_n(\C)$ let $D_Y$ be the derivation
of $\tliftPP$ with $(D_YP)(Z)=\frac d{dt}\big|_{t=0}P({}^t\exp(tY)Z)$, so that $\mathcal L(Y)=D_Y+\frac d2\tr Y$; and for
$Y\in M_d(\C)$ let $D'_Y$ be the derivation with $(D'_YP)(Z)=\frac d{dt}\big|_{t=0}P(Z\exp(tY))$, and $\mathcal R(Y)=D'_Y+\frac n2\tr Y$.

\begin{lemma}[Exponentials of the infinitesimal actions]\label{pf:tlift:expfirst}
Let $c\in\Z$ and $\kappa\in\C$.
\begin{enumerate}[label=(\alph*),nosep]
\item If $Y\in M_n(\C)$ and $2c=d+2\kappa$, then $L_c(\exp Y)=\exp(\mathcal L(Y)+\kappa\tr Y)$ weakly.
\item If $Y\in M_d(\C)$ and $2c=n+2\kappa$, then $R_c(\exp Y)=\exp(\mathcal R(Y)+\kappa\tr Y)$ weakly.
\end{enumerate}
\formallabel{UpperTheory.TP.RP.expFirst}\formallabel{UpperTheory.TP.RP.expSecond}
\end{lemma}

\begin{proof}
(a) Since $2c=d+2\kappa$, $\mathcal L(Y)+\kappa\tr Y=D_Y+c\tr Y$. The derivation $D_Y$ preserves the finite-dimensional space
$\tliftPP_{\le k}$ of polynomials of degree at most $k$, and so does the substitution $S(g)P(Z)=P({}^tgZ)$. On $\tliftPP_{\le k}$ the map
$t\mapsto S(\exp tY)$ is a continuous one-parameter group of linear maps whose derivative at $t=0$ is $D_Y$, by the chain rule;
hence $S(\exp tY)=\exp(tD_Y)$ on $\tliftPP_{\le k}$ for all $t$, and so $S(\exp Y)=\exp(D_Y)$ weakly on $\tliftPP$. Since $c\tr Y$ is a
scalar commuting with $D_Y$, $\exp(D_Y+c\tr Y)=e^{c\tr Y}S(\exp Y)=\det(\exp Y)^cS(\exp Y)=L_c(\exp Y)$, using $\det\exp Y=e^{\tr Y}$.
(b) is proved in the same way with right substitution.
\end{proof}

In the next two lemmas let $W'$ be a space of dimension $N\ge1$ with frames $C_w$ at the infinite places in which
$b\,\tau_w(J_{W'})$ becomes the identity for a constant $b$, let $\omega'$ be a homomorphism from $\Uni(W')(\A)$ to the linear operators on
functions on $\A^\kappa$ such that $\lambda\circ\omega'(g)$ is admissible for every admissible $\lambda$, let
$\Phi_F\colon\bigotimes_w\tliftPP'\to\tliftS(F_\infty^\kappa)$ be injective and linear ($\tliftPP'$ a polynomial ring), let
$Y\mapsto\mathrm{act}(Y)$ be a linear map from $M_N(\C)$ to the operators on $\tliftPP'$ such that every $\mathrm{act}(Y)+z$ is locally finite,
and let $(c_w)$ be complex numbers. Assume:
\begin{itemize}[nosep]
\item[$(\ast)$] for every continuous homomorphism $\gamma\colon\R\to\Uni(W')(\A)$ with values supported at infinity and all
  $Y_w\in M_N(\C)$ with $\tau_w(\gamma(t))C_w=C_w\exp(tY_w)$ for all $t$ and $w$, the derivative at $t=0$ of $\omega'(\gamma(t))$ on
  $\Phi_F(T)$ is $\Phi_F\bigl(\sum_w(\mathrm{act}(Y_w)+c_w\tr Y_w)_wT\bigr)$, in the sense that
  $\frac d{dt}\big|_{t=0}\lambda\bigl(\omega'(\gamma(t))(\Phi_F(T)\otimes\phi)\bigr)=\lambda\bigl(\Phi_F(\cdots)\otimes\phi\bigr)$ for every finite
  part $\phi$ and every admissible $\lambda$.
\end{itemize}
By a general exponentiation statement for locally finite infinitesimal actions
(\lref{GlobalTheory.T.weakExponential_action_on_fockVectors}), $(\ast)$ implies that $\omega'(\gamma(1))$ acts at infinity by
$\Phi_F(T)\mapsto\Phi_F(UT)$, where $U=\exp\bigl(\sum_w(\mathrm{act}(Y_w)+c_w\tr Y_w)_w\bigr)$ weakly.

\begin{lemma}[Parity]\label{pf:tlift:parity}
In this situation assume moreover that $\mathrm{act}(Y)1=\frac D2\tr(Y)\cdot1$ for all $Y$, for some $D\in\N$, and that $c_w=\frac{m(w)}2$ with
$m(w)\in\Z$. Then $D+m(w)$ is even for every infinite place $w$.
\formallabel{UpperTheory.TP.RP.FockModelActions.parity_of_pk}
\end{lemma}

\begin{proof}
Fix $w$ and let $Y_0=\mathrm{diag}(2\pi i,0,\dots,0)\in M_N(\C)$, which is skew-hermitian with $\exp Y_0=1$ and $\tr Y_0=2\pi i$; let $Y_w=Y_0$ and
$Y_{w'}=0$ for $w'\ne w$. Since the $Y_{w'}$ lie in the Lie algebra of the unitary group of the identity form, there is a
continuous homomorphism $\gamma\colon\R\to\Uni(W')(\A)$ with values supported at infinity and $\tau_{w'}(\gamma(t))C_{w'}=C_{w'}\exp(tY_{w'})$
(\lref{GlobalTheory.T.unitaryLieElements_generate_oneParamLines}). Then $\tau_{w'}(\gamma(1))C_{w'}=C_{w'}$ for all $w'$, and since the
$C_{w'}$ are invertible and $\gamma(1)$ is supported at infinity, $\gamma(1)=1$. By $(\ast)$ and the exponentiation statement,
$\omega'(\gamma(1))=1$ acts by $\Phi_F(T)\mapsto\Phi_F(UT)$, so $\Phi_F(UT_0)=\Phi_F(T_0)$ for $T_0=\otimes_{w'}1$, and $UT_0=T_0$ by injectivity.
On the other hand $T_0$ is an eigenvector of $L=\sum_{w'}(\mathrm{act}(Y_{w'})+c_{w'}\tr Y_{w'})_{w'}$ with eigenvalue
$\mu=\bigl(\frac D2+\frac{m(w)}2\bigr)2\pi i$, so $UT_0=e^\mu T_0$ by the exponential series. Hence $e^\mu=1$, $\mu\in2\pi i\Z$, and
$\frac{D+m(w)}2\in\Z$.
\end{proof}

\begin{lemma}[Group action from the Lie algebra action]\label{pf:tlift:actsofk}
In the same situation assume that $e\colon\Uni_N(\A)\to\Uni(W')(\A)$ is as in
\cref{pf:tlift:brex}, the Cartan involutions satisfy $\vartheta_wC_w=C_w$, and $\tau_w(g)C_w=C_w\exp(Y)$ implies $\tau_w(g)=\exp(Y)$.
Let $U^w(g)$, $g\in\GL_N(\C)$, be operators on $\tliftPP'$ with $U^w(\exp Y)=\exp(\mathrm{act}(Y)+c_w\tr Y)$ weakly for all $Y\in M_N(\C)$.
Then for every $k\in\Uni_N(E_\infty)$ the operator $\omega'(e(k))$ acts at infinity by
$\Phi_F(T)\mapsto\Phi_F\bigl((\otimes_wU^w(\tau_w(k)))T\bigr)$.
\formallabel{UpperTheory.TP.RP.FockModelActions.acts_of_k_pk}
\end{lemma}

\begin{proof}
The element $e(k)$ is supported at infinity and commutes with the Cartan involutions, which are trivial. Every such element
lies on a one-parameter subgroup: there are $\gamma$ as in $(\ast)$ and $Y_w$ with $\gamma(1)=e(k)$
(\lref{GlobalTheory.T.maximalCompact_oneParameterLines_of_lieLines}; this is the surjectivity of the exponential map of
the compact connected groups $\Uni(N)$). Then $\tau_w(e(k))=\tau_w(k)$ (\cref{pf:tlift:brex}) and $\tau_w(k)C_w=C_w\exp(Y_w)$, hence
$\tau_w(k)=\exp(Y_w)$. By $(\ast)$ and the exponentiation statement, $\omega'(e(k))$ acts by $\Phi_F(T)\mapsto\Phi_F(UT)$ with
$U=\exp(\sum_w L_w)$, $L_w=(\mathrm{act}(Y_w)+c_w\tr Y_w)_w$. The operators $L_w$ act on different tensor factors and are locally
finite; on a tensor product of finite-dimensional stable subspaces they commute, so $\exp(\sum_wL_w)=\prod_w\exp(L_w)$ there,
and $\exp(L_w)$ is $\exp(\mathrm{act}(Y_w)+c_w\tr Y_w)=U^w(\tau_w(k))$ on the $w$-th factor. Hence $U=\otimes_wU^w(\tau_w(k))$, by uniqueness of weak exponentials.
\end{proof}

\begin{proposition}[Fock model for the definite pair]\label{pf:tlift:fock}
There are an injective linear map $\Phi_F\colon\bigotimes_w\tliftPP\to\tliftS(F_\infty^\kappa)$ with image $\tliftPol$ and integers
$a_w,b_w$ with $2a_w=d+m_W(w)$ and $2b_w=n+m_X(w)$ such that for $T\in\bigotimes_w\tliftPP$:
\begin{enumerate}[label=(\alph*),nosep]
\item for $k\in\Uni_n(E_\infty)$, $\omega(e(k),1)$ acts at infinity by $\Phi_F(T)\mapsto\Phi_F\bigl((\otimes_wL_{a_w}(\tau_w(k)))T\bigr)$;
\item for $h\in\Uni_d(E_\infty)$, $\omega(1,e'(h))$ acts at infinity by $\Phi_F(T)\mapsto\Phi_F\bigl((\otimes_wR_{b_w}(\tau_w(h)))T\bigr)$.
\end{enumerate}
\formallabel{UpperTheory.TP.RP.fockOf_k}
\end{proposition}

\begin{proof}
Take $\Phi_F$ as above. The infinitesimal action of $\Uni(X)(F_\infty)$ satisfies $(\ast)$ with $\omega'=\omega(\cdot,1)$,
$\mathrm{act}=\mathcal L$ and $c_w=\frac{m_W(w)}2$, and $\mathcal L(Y)1=\frac d2\tr(Y)\,1$; so $a_w=\frac{d+m_W(w)}2\in\Z$ by \cref{pf:tlift:parity}.
By \cref{pf:tlift:expfirst}(a) with $c=a_w$ and $\kappa=\frac{m_W(w)}2$, the operators $U^w(g)=L_{a_w}(g)$ satisfy
$U^w(\exp Y)=\exp(\mathcal L(Y)+\frac{m_W(w)}2\tr Y)$, and \cref{pf:tlift:actsofk} gives (a). The same argument for $\Uni(W)(F_\infty)$, with
$\mathcal R$, $c_w=\frac{m_X(w)}2$, $D=n$ and \cref{pf:tlift:expfirst}(b), gives $b_w\in\Z$ and (b).
\end{proof}

\subsubsection*{Joint highest weight vectors}
\begin{lemma}[Equivariant maps hitting $\Delta^\mu$]\label{pf:tlift:joint}
Let $n\le d$ and $\mu\in\N^n$ non-increasing.
\begin{enumerate}[label=(\alph*),nosep]
\item If $\lambda\in\Z^n$ is dominant, $a\in\Z$ and $\lambda_k=\mu_k+a$ for all $k$, there is a linear map
  $i\colon V_\lambda\to\tliftPP$ with $i(V_\lambda(g)x)=L_a(g)i(x)$ for $g\in\GL_n(\C)$ and $\Delta^\mu\in i(V_\lambda)$.
\item If $\lambda'\in\Z^d$ is dominant, $b\in\Z$ and $\lambda'_k=\mu_k+b$ for $k\le n$, $\lambda'_k=b$ for $k>n$, there is a linear
  map $i\colon V_{\lambda'}\to\tliftPP$ with $i(V_{\lambda'}(h)x)=R_b(h)i(x)$ for $h\in\GL_d(\C)$ and $\Delta^\mu\in i(V_{\lambda'})$.
\end{enumerate}
\formallabel{UpperTheory.TP.RP.exists_act1_equivariant_nablaSq_map_hitting_joint}%
\formallabel{UpperTheory.TP.RP.exists_act2_equivariant_nablaSq_map_hitting_joint}
\end{lemma}

\begin{proof}
In the model of \cref{pf:tlift:aut}(c), $V_\lambda$ consists of polynomials on $M_n(\C)$, $g$ acts by
$\det(g)^{\lambda_n}$ times left substitution by ${}^tg$, and the highest weight vector is the product of leading
principal minors $\prod_s\Delta_s^{(\lambda_s-\lambda_n)-(\lambda_{s+1}-\lambda_n)}$ (with the convention $\lambda_{n+1}=\lambda_n$).

(a) Let $\sigma$ send a polynomial $P$ on $M_n$ to $Z\mapsto P(Z_0)$, $Z_0$ the first $n$ columns of $Z$, and put
$i(P)=\sigma(P)\,\Delta_n^{\lambda_n-a}$ (note $\lambda_n-a=\mu_n\ge0$). Since $\sigma$ commutes with left substitution and
$\Delta_n({}^tgZ)=\det(g)\Delta_n(Z)$, $i$ intertwines $\det^{\lambda_n}\cdot(\text{left substitution})$ with $L_a$; the highest weight
vector goes to $\prod_s\Delta_s^{\mu_s-\mu_{s+1}}=\Delta^\mu$, because $\lambda_s-\lambda_{s+1}=\mu_s-\mu_{s+1}$.

(b) Let $\sigma'$ send a polynomial $P$ on $M_d$ to $Z\mapsto P(\tilde Z)$, where the first $n$ columns of
$\tilde Z\in M_d(\C)$ form ${}^tZ$ and the others vanish, and put $i(P)=\sigma'(P)\,\Delta_n^{\lambda'_d-b}$ (the exponent is $0$ if
$n<d$ and $\mu_n$ if $n=d$). Left substitution by ${}^th$ on $\tilde Z$ corresponds to $Z\mapsto Zh$, and if $n=d$,
$\Delta_n(Zh)=\det(h)\Delta_n(Z)$; so $i$ intertwines with $R_b$. The highest weight vector of $V_{\lambda'}$ is sent to
$\Delta^\mu$ by the same computation.
\end{proof}

\begin{lemma}[Archimedean types of $U$]\label{pf:tlift:formex}
Let $U$ be an automorphic representation of $\Uni_n$ with weights $\lambda_w$ at $\tau_w$. Then
$V=\bigotimes_wV_{\lambda_w}$, with $\rho(k)=\otimes_wV_{\lambda_w}(\tau_w(k))$ for $k\in\Uni_n(E_\infty)$, is a finite-dimensional
irreducible representation, and there is an injective linear map $j\colon V\to U$ with $j(\rho(k)x)=R(k)j(x)$.
\formallabel{UpperTheory.TP.RP.formEx}
\end{lemma}

\begin{proof}
For each $w$, the subgroup $\Uni_n(F_w)$ commutes with the elements of $\Uni_n(\A)$ trivial at $w$, so the
$V_{\lambda_w}$-isotypic part of $U$ for $\Uni_n(F_w)$ is $\Uni_n(\A)$-invariant; it is nonzero because $U$ has weight
$\lambda_w$ at $\tau_w$, hence equal to $U$ by irreducibility (\lref{GlobalTheory.DefU.iSup_range_weight_eq_top}).
The subgroups $\Uni_n(F_w)$ commute with each other and generate $\Uni_n(E_\infty)$
(\lref{UpperTheory.archimedeanUnitary_le_iSup_embU_and_commute}); for a representation of such a group which is
$V_{\lambda_w}$-isotypic for each factor there is a nonzero multilinear intertwiner $\prod_wV_{\lambda_w}\to U$
(\lref{UpperTheory.exists_multilinear_intertwiner_of_commuting_subgroups}), i.e.\ a nonzero equivariant
$j\colon V\to U$. The external tensor product of the irreducible representations $V_{\lambda_w}$ of the factors
(\lref{GlobalTheory.DefU.isIrreducible_sqRep}) is irreducible, so the invariant subspace $\ker j$ is $0$.
\end{proof}

\begin{definition}[Archimedean type of a Fock vector for $\Uni(W)$]\label{pf:tlift:hastype}
Let $f_0\in\tliftS(F_\infty^\kappa)$, $w$ an infinite place and $\lambda'\in\Z^d$ dominant. We say that $f_0$ \emph{has type
$\lambda'$ at $w$} if there is a linear map $i\colon V_{\lambda'}\to\tliftPol$ with $f_0\in i(V_{\lambda'})$ such that $\omega(1,e'(k))$ acts at
infinity by $i(x)\mapsto i(V_{\lambda'}(\tau_w(k))x)$ for all $k\in\Uni_d(F_w)$ and $x\in V_{\lambda'}$.
\formal{UpperTheory.TP.HasType}
\end{definition}

\begin{proposition}[The archimedean pairing]\label{pf:tlift:realpair}
Let $U$ be an automorphic representation of $\Uni_n$ with weights $\lambda_w$ at $\tau_w$ satisfying
\eqref{eq:tlift:dom}. There are $0\ne f_0\in\tliftPol$ and $0\ne\varphi\in U$ of the same type
(\cref{pf:tlift:sametype}), and for every infinite place $w$ a dominant $\lambda'_w\in\Z^d$ with
$\tliftLW(\lambda_w,m_X(w),m_W(w),\lambda'_w)$ such that $f_0$ has type $\lambda'_w$ at $w$.
\formallabel{UpperTheory.TP.realPair_k}
\end{proposition}

\begin{proof}
Let $\Phi_F,a_w,b_w$ be as in \cref{pf:tlift:fock}. By \eqref{eq:tlift:dom}, $2\lambda_{w,k}\ge d+m_W(w)=2a_w$, so
$\mu_w=\lambda_w-a_w\in\N^n$, and $\mu_w$ is non-increasing. By \cref{pf:tlift:joint}(a) there are $i_w\colon V_{\lambda_w}\to\tliftPP$,
intertwining $V_{\lambda_w}$ with $L_{a_w}$, and $x_w\in V_{\lambda_w}$ with $i_w(x_w)=\Delta^{\mu_w}\ne0$. Let
$(V,\rho)$ and $j$ be as in \cref{pf:tlift:formex}, and put
\[
  T_0=\otimes_w\Delta^{\mu_w},\qquad f_0=\Phi_F(T_0),\qquad x_0=\otimes_wx_w,\qquad\varphi=j(x_0).
\]
Then $f_0\in\tliftPol$, $f_0\ne0$ because $\Phi_F$ is injective and $T_0\ne0$, and $\varphi\in U$, $\varphi\ne0$ because $j$ is injective and
$x_0\ne0$. With $i=\Phi_F\circ(\otimes_wi_w)$ we have $i(x_0)=f_0$ and, by \cref{pf:tlift:fock}(a), $\omega(e(k),1)$ acts at
infinity by $i(x)\mapsto i(\rho(k)x)$; together with the equivariance of $j$ this shows that $f_0$ and $\varphi$ have the
same type.

Fix $w$ and let $\lambda'_w=(\mu_w,0,\dots,0)+b_w\in\Z^d$; it is dominant since $\mu_w\ge0$ is non-increasing, and
$2\lambda'_{w,k}=2\mu_{w,k}+2b_w=2\lambda_{w,k}-d-m_W(w)+n+m_X(w)$ for $k\le n$, $2\lambda'_{w,k}=n+m_X(w)$ for $k>n$; this is
\eqref{eq:tlift:liftwt}. By \cref{pf:tlift:joint}(b) there are $i'\colon V_{\lambda'_w}\to\tliftPP$ intertwining $V_{\lambda'_w}$ with
$R_{b_w}$ and $x'$ with $i'(x')=\Delta^{\mu_w}$. Let $\mathrm{up}_w(P)$ be the pure tensor with $P$ at $w$ and $\Delta^{\mu_{w'}}$
at $w'\ne w$, and $i=\Phi_F\circ\mathrm{up}_w\circ i'$. Then $i(x')=f_0$, $i$ takes values in $\tliftPol$, and for $k\in\Uni_d(F_w)$,
$\tau_{w'}(k)=1$ for $w'\ne w$, so by \cref{pf:tlift:fock}(b) $\omega(1,e'(k))$ acts at infinity by
$i(x)\mapsto\Phi_F(\mathrm{up}_w(R_{b_w}(\tau_w(k))i'(x)))=i(V_{\lambda'_w}(\tau_w(k))x)$. Thus $f_0$ has type $\lambda'_w$ at $w$.
\end{proof}


% ---------------------------------------------------------------------------
\subsection{Theta lifts are automorphic forms}\label{pf:tlift:autsec}

We keep the assumptions of \cref{pf:tlift:nonvan}. For $\Phi\in\tliftPD(M)$ and an automorphic form $\varphi$ on $\Uni_n$
put
\begin{equation}\label{eq:tlift:Theta}
  \tliftlift(\Phi,\varphi)=\tliftlift_M(\Phi,\varphi^\flat)\circ e'\colon\Uni_d(\A)\to\C.
\end{equation}
(In the formal statement $\tliftlift$ is defined for all functions $\varphi$ on $\Uni_n(\A)$, by first projecting
$\varphi$ linearly onto the space of continuous left $\Uni_n(E)$-invariant functions; this makes $\tliftlift$
sesquilinear on all functions and does not change it on automorphic forms.)

\begin{lemma}[Archimedean operators]\label{pf:tlift:archop}
Let $k\in\Uni(W)(\A)$ be supported at infinity and $f\in\tliftS(F_\infty^\kappa)$. There is $f_1\in\tliftS(F_\infty^\kappa)$ with
$\omega(1,k)(f\otimes\phi)=f_1\otimes\phi$ for every locally constant compactly supported function $\phi$ on $\A_f^\kappa$.
\formallabel{PositivityTheory.ThKOwn.KPol.exists_arch_operator}
\end{lemma}

\begin{proof}
For $\phi$ as in the statement, $\omega(1,k)(f\otimes\phi)$ is a Schwartz--Bruhat function by (W1). For $\xi\in F_\infty^\kappa$ let $T_\xi\phi$ be
the function $y\mapsto\omega(1,k)(f\otimes\phi)(\xi,y)$ on $\A_f^\kappa$; it is locally constant with compact support, and $T_\xi$ is linear.
The translations $\phi\mapsto\phi(\cdot+x)$ and the multiplications by the characters $\phi\mapsto\psi(\langle\cdot,y\rangle)\phi$, for $x,y$ with
trivial archimedean components, are induced by Heisenberg operators $\rho(u)$ with $u$ having trivial archimedean components;
for such $u$, $u\,{}^tk=u$ because $k$ has trivial finite components, so $\omega(1,k)$ commutes with $\rho(u)$ by (W1). Hence $T_\xi$
commutes with all translations and all multiplications by characters, and is therefore a scalar $c(\xi)$
(\lref{GlobalTheory.CompanionTheory.SW.SchurChar.schur_translations_characters}). Choose $\phi_0$ with $\phi_0(0)=1$ and let
$f_1(\xi)=\omega(1,k)(f\otimes\phi_0)(\xi,0)$; then $f_1=c$, $f_1$ is a Schwartz function (a slice of a Schwartz--Bruhat function),
and $\omega(1,k)(f\otimes\phi)(\xi,y)=c(\xi)\phi(y)=f_1(\xi)\phi(y)$.
\end{proof}

For a Haar measure $\mu$ on $\A^\kappa$ and $\Phi_1,\Phi_2\in\tliftS(\A^\kappa)$ let
$c(\Phi_1,\Phi_2)(z)=\int(\rho(z)\Phi_1)\,\overline{\Phi_2}\,d\mu$ ($z\in M_{n\times d}(\tliftAE)$) be their Heisenberg matrix coefficient.

\begin{lemma}[Smoothness at the finite places]\label{pf:tlift:nhdsfix}
Let $\phi$ be a locally constant compactly supported function on $\A_f^\kappa$. There is a neighbourhood $U$ of $1$ in $\Uni(X)(\A)$ such that
$\omega(u,1)(f\otimes\phi)=f\otimes\phi$ for every $u\in U$ with trivial archimedean components and every $f\in\tliftS(F_\infty^\kappa)$.
\formallabel{UpperTheory.exists_nhds_one_weil_fixes_pure_tensor}
\end{lemma}

\begin{proof}
If all $f\otimes\phi$ vanish there is nothing to prove; otherwise fix $f_0$ and $x_0$ with $\Phi_0(x_0)\ne0$, $\Phi_0=f_0\otimes\phi$. We use
three facts on the coefficients.
\begin{enumerate}[label=(\roman*),nosep]
\item By (W1) (unitarity and the Heisenberg relation),
  $c(\omega(u,1)\Phi_1,\omega(u,1)\Phi_2)(z)=c(\Phi_1,\Phi_2)(u^{-1}z)$.
\item There is a compact set $C$ of matrices with trivial archimedean components such that $c(f\otimes\phi,f'\otimes\phi)(z)=0$ unless
  $z\in C+M_{n\times d}(E_\infty)$, for all $f,f'$ (\lref{PositivityTheory.ThKEst.coefBd_of_k}): the translation part of $\rho(z)$ moves the
  support of $\phi$, and the character part produces a Fourier transform of a locally constant compactly supported function.
\item The coefficients determine the functions: if $c(\Phi_1,\Phi_2)=c(\Phi_3,\Phi_4)$, then $\Phi_1(x)\overline{\Phi_2(y)}=\Phi_3(x)\overline{\Phi_4(y)}$
  for all $x,y$ (\lref{PositivityTheory.ThKEst.fourierInj_of}).
\end{enumerate}
The coefficient $c(f\otimes\phi,f'\otimes\phi)(z)$ is the product of an archimedean coefficient of $f,f'$ at the archimedean part of $z$ and of
the coefficient $c(\phi,\phi)$ at the finite part, which is locally constant; as $C$ is compact, there is a symmetric neighbourhood
$U_1$ of $1$ such that $c(f\otimes\phi,\Phi_0)(u^{-1}z)=c(f\otimes\phi,\Phi_0)(z)$ for all $f$, all $u\in U_1$ with trivial archimedean components and
all $z\in C+M_{n\times d}(E_\infty)$. Let $U_2$ be a neighbourhood of $1$ on which $|\omega(u,1)\Phi_0(x_0)/\Phi_0(x_0)-1|<\frac12$ (continuity, (W1)),
let $K'\subset U_1\cap U_2$ be a compact open subgroup of the group of elements with trivial archimedean components (such
subgroups form a basis of neighbourhoods of $1$ there),
and let $U=\tliftGinf K'$, an open neighbourhood of $1$ (\cref{pf:tlift:decomp}).

Let $u\in U$ have trivial archimedean components, so $u\in K'$. For every $f$ and $z$, $c(\omega(u,1)(f\otimes\phi),\omega(u,1)\Phi_0)(z)=
c(f\otimes\phi,\Phi_0)(u^{-1}z)=c(f\otimes\phi,\Phi_0)(z)$: by (i), and then by the choice of $U_1$ (applied to $u$, or to $u^{-1}$ and
$u^{-1}z$) if $z$ or $u^{-1}z$ lies in $C+M_{n\times d}(E_\infty)$, and by (ii) (both sides vanish) otherwise. By (iii),
$\omega(u,1)(f\otimes\phi)(x)\,\overline{\omega(u,1)\Phi_0(x_0)}=(f\otimes\phi)(x)\,\overline{\Phi_0(x_0)}$ for all $x$; for $f=f_0$, $x=x_0$ this shows
$\omega(u,1)\Phi_0(x_0)\ne0$. Hence $\omega(u,1)$ acts on all $f\otimes\phi$ by the same scalar $\lambda_u$, and by the homomorphism property
$\omega(u^k,1)$ acts by $\lambda_u^k$. As $u^k\in K'\subset U_2$ for all $k$, $|\lambda_u^k-1|<\frac12$ for all $k\in\N$, and this forces $\lambda_u=1$
(\lref{PositivityTheory.ThKH.Locrep.eq_one_of_pow_near}).
\end{proof}

\begin{lemma}[Finiteness of archimedean orbits]\label{pf:tlift:polfin}
For $f\in\tliftPol$ there is a finite-dimensional subspace $N_f\subset\tliftS(F_\infty^\kappa)$ which contains every $f_1$ such
that $\omega(1,k)$ acts at infinity by $f\mapsto f_1$ for some $k\in\Uni(W)(\A)$ supported at infinity. Consequently,
for every $\Phi\in\tliftPD(M)$ the vectors $\omega(1,k)\Phi$, $k$ supported at infinity, lie in a finite-dimensional
subspace of $\tliftS(\A^\kappa)$.
\formallabel{GlobalTheory.ThK.Dbl.polFin_k}%
\formallabel{UpperTheory.exists_finiteDimensional_weil_archimedean_orbit}
\end{lemma}

\begin{proof}
Since $W$ is definite, its Cartan involutions are $\vartheta_{W,w}=1$, so every $k$ supported at infinity commutes with
them. Each such $k$ acts at infinity through an operator $A_k$ on $\tliftS(F_\infty^\kappa)$
(\cref{pf:tlift:archop}), and by (W1)
$A_k\,D_u=D_{u\,{}^tk}\,A_k$ for archimedean $u$ (\lref{PositivityTheory.ThKOwn.KPol.inter}). Let $\mathcal N$ be the set of
$f$ for which a finite-dimensional $N_f$ with $A_kf\in N_f$ for all $k$ exists; it is a subspace. It contains the
Gaussians: $A_k$ maps a Gaussian with signs $\varepsilon$ to a Gaussian with the same signs, because the defining relation
of \cref{pf:tlift:pol} is preserved under $u\mapsto u\,{}^tk$ (\lref{PositivityTheory.ThKOwn.KPol.gauss_impl}), and the
Gaussians with given signs span a finite-dimensional space (\lref{GlobalTheory.ThK.Dbl.Dense.vacSp_fin}); there are
finitely many sign vectors. It is stable under each $D_u$: if $A_kf\in N_f$ for all $k$, then
$A_kD_uf=D_{u{}^tk}A_kf$ lies in the span of the $D_{u'}N_f$, $u'$ archimedean, which is finite-dimensional because
$u'\mapsto D_{u'}g$ is $\R$-linear on the finite-dimensional real space of archimedean $u'$
(\lref{GlobalTheory.ThK.Dbl.Dense.exists_dop_fin}). By minimality of $\tliftPol$, $\tliftPol\subset\mathcal N$. For the
second assertion, $\omega(1,k)(f\otimes\phi)=(A_kf)\otimes\phi$ for pure tensors with finite data
at infinity, so the span of the $N_f\otimes\phi$ works for $f\otimes\phi$, and sums of such
spaces work for $\tliftPD(M)$.
\end{proof}

\begin{lemma}[Stability of $\tliftPD$]\label{pf:tlift:polstab}
$\omega(x,y)\tliftPD(M)\subset\tliftPD(M)$ for all $x\in\Uni(X)(\A)$, $y\in\Uni(W)(\A)$.
\formallabel{UpperTheory.TP.polStab_k}
\end{lemma}

\begin{proof}
It suffices to treat $(x,1)$ and $(1,y)$. For elements supported at infinity this holds because the operators $A_k$
preserve $\tliftPol$ (\lref{PositivityTheory.ThKF.A.archOp_k}); for elements supported at a
finite place $v$, $\omega$ only changes the factor at $v$ of a pure tensor, which remains in $\tliftS(F_v^\kappa)$
(\cref{pf:tlift:acts}). Every $\Phi\in\tliftPD$ is fixed by
$\omega(u,1)$ (and by $\omega(1,u)$) for all $u$ in a neighbourhood of $1$ with trivial archimedean components
(\cref{pf:tlift:nhdsfix}, \lref{PositivityTheory.ThKEst.level_k}): the finite
data are fixed by a compact open subgroup, and at almost all places the vacuum vectors are fixed by the
integral points. A general $x$ is the product of its archimedean component, its components at a finite set $S$
of places, and an element $u$ which tends to $1$ as $S$ grows; for $S$ large, $u$ fixes $\Phi$, and the claim follows
from the previous cases.
\end{proof}

\begin{proposition}[Automorphy and equivariance of the lift]\label{pf:tlift:liftaut}
For $\Phi\in\tliftPD(M)$ and an automorphic form $\varphi$ on $\Uni_n$:
\begin{enumerate}[label=(\alph*),nosep]
\item $\tliftlift(\Phi,\varphi)$ is an automorphic form on $\Uni_d$;
\item $\tliftlift$ is linear in $\Phi$ and antilinear in $\varphi$;
\item for $g\in\Uni_n(\A)$ and $h\in\Uni_d(\A)$, $\Phi_1=\omega(e(g),e'(h))\Phi$ lies in $\tliftPD(M)$ and
  $\tliftlift(\Phi_1,R(g)\varphi)=R(h)\tliftlift(\Phi,\varphi)$.
\end{enumerate}
\formallabel{UpperTheory.TP.liftAut_k}\formallabel{UpperTheory.TP.aut_lift_k}
\end{proposition}

\begin{proof}
Write $\varphi^\flat=\varphi\circ e^{-1}$; it is continuous, left $\Uni(X)(F)$-invariant and bounded. All functions below are
integrable over $D$ by boundedness of $\theta$ in $g$ (\cref{pf:tlift:basic}, \cref{pf:tlift:cpt}).

(b) is clear from \eqref{eq:tlift:theta}.

(c) $\Phi_1\in\tliftPD$ by \cref{pf:tlift:polstab}. For $a\in\Uni(X)(\A)$, $b\in\Uni(W)(\A)$ we have
$\theta(g,h;\omega(a,b)\Phi)=\theta(ga,hb;\Phi)$, so
\begin{align*}
  \tliftlift_M(\omega(a,b)\Phi,\varphi^\flat(\,\cdot\,a))(h)&=\int_D\theta(ga,hb;\Phi)\overline{\varphi^\flat(ga)}\,d\mu(g)\\
  &=\int_{Da}\theta(g,hb;\Phi)\overline{\varphi^\flat(g)}\,d\mu(g)=\tliftlift_M(\Phi,\varphi^\flat)(hb),
\end{align*}
because $Da$ is again a fundamental domain, $\mu$ is right invariant and the integrand is left
$\Uni(X)(F)$-invariant. With $a=e(g)$, $b=e'(h)$ and
$(R(g)\varphi)^\flat=\varphi^\flat(\,\cdot\,e(g))$ this is (c).

(a) \emph{Continuity}: $(g,h)\mapsto\theta(g,h;\Phi)\overline{\varphi^\flat(g)}$ is continuous and bounded, so the integral over the
set $D$ of finite measure is continuous in $h$ by dominated convergence. \emph{Left invariance}: $\theta(g,\gamma h;\Phi)=\theta(g,h;\Phi)$ for
$\gamma\in\Uni(W)(F)$, and $e'$ preserves rational points (\cref{pf:tlift:brex}). \emph{Smoothness}: by
\lref{PositivityTheory.ThKEst.level_k} there is a neighbourhood $N'$ of $1$ such that $\omega(1,u)\Phi=\Phi$ for
$u\in N'$ with trivial archimedean components; for $k\in\Uni_d(\A_f)$ with $e'(k)\in N'$, (c) with $g=1$ gives
$R(k)\tliftlift(\Phi,\varphi)=\tliftlift(\omega(1,e'(k))\Phi,\varphi)=\tliftlift(\Phi,\varphi)$; so the stabiliser contains a
neighbourhood of $1$ in $\Uni_d(\A_f)$, and being a subgroup it is open. \emph{$\Uni_d(E_\infty)$-finiteness}: for
$k\in\Uni_d(E_\infty)$, $R(k)\tliftlift(\Phi,\varphi)=\tliftlift(\omega(1,e'(k))\Phi,\varphi)$ by (c), and $\omega(1,e'(k))\Phi$ lies in the
finite-dimensional space of \cref{pf:tlift:polfin} ($e'(k)$ is supported at infinity), so these translates lie
in the image of a finite-dimensional space under the linear map $\Psi\mapsto\tliftlift(\Psi,\varphi)$.
\end{proof}

% ---------------------------------------------------------------------------
\subsection{Projection and reading off the lift}\label{pf:tlift:read}

\begin{lemma}[Projection to an automorphic representation]\label{pf:tlift:projex}
Let $\psi\ne0$ be an automorphic form on $\Uni_d$. There are an automorphic representation $U'$ of $\Uni_d$ and a linear
map $\mathrm{pr}$ from the functions on $\Uni_d(\A)$ to themselves such that $\mathrm{pr}(\psi)\ne0$, and for every automorphic form
$f$: $\mathrm{pr}(f)\in U'$ and $\mathrm{pr}(R(x)f)=R(x)\mathrm{pr}(f)$ for all $x\in\Uni_d(\A)$.
\formallabel{UpperTheory.TP.projEx}
\end{lemma}

\begin{proof}
Let $\mathcal A$ be the space of automorphic forms on $\Uni_d$. There is a positive definite hermitian form on $\mathcal A$
invariant under right translations (the Petersson inner product over a fundamental domain;
\lref{GlobalTheory.T.exists_invariant_positive_petersson_form}), and with it every automorphic form lies in the
sum of the automorphic representations of $\Uni_d$ (\lref{UpperTheory.mem_iSup_automorphicReps_of_invariant_form}).
It remains to prove the following algebraic statement: if $G$ acts linearly on a
vector space $V$, $\mathcal A\subset V$ is a subspace contained in the sum of a family $\mathcal S$ of invariant irreducible subspaces
$Y\subset\mathcal A$, and $0\ne\psi\in\mathcal A$, then there are $U'\in\mathcal S$ and a linear $\mathrm{pr}\colon V\to V$ with $\mathrm{pr}(\psi)\ne0$, $\mathrm{pr}(\mathcal A)\subset U'$
and $\mathrm{pr}$ equivariant on $\mathcal A$.

By Zorn's lemma there is a maximal invariant subspace $C\subset\mathcal A$ with $\psi\notin C$ (a union of a chain of such
subspaces is again one). Some $U'\in\mathcal S$ is not contained in $C$, since $\psi\notin C$ lies in the sum of $\mathcal S$. Then
$C\cap U'=0$ by irreducibility, and $\psi\in C\oplus U'$ by maximality. Every $Y\in\mathcal S$ is contained in $C\oplus U'$: otherwise
$Y\not\subset C$, $(C\oplus U')\cap Y=0$ by irreducibility, and $\psi\in C+Y$ by maximality, say $\psi=c+y$; then $y=\psi-c\in(C\oplus U')\cap Y$,
so $y=0$ and $\psi\in C$, a contradiction. Hence $\mathcal A\subset C\oplus U'$. Choose a complement $D$ of $C\oplus U'$ in $V$ and let $\mathrm{pr}$
be the projection onto $U'$ along $C\oplus D$. Writing $\psi=c+u$ with $c\in C$, $u\in U'$, we get $\mathrm{pr}(\psi)=u\ne0$; and for
$f=c+u\in\mathcal A$, $\mathrm{pr}(R(x)f)=\mathrm{pr}(R(x)c+R(x)u)=R(x)u=R(x)\mathrm{pr}(f)$ since $C$ and $U'$ are invariant.
\end{proof}

\begin{lemma}[Isotypic components]\label{pf:tlift:absiso}
Let a group $G$ act linearly on $V$ by $\rho$, let $U\subset V$ be an invariant subspace without invariant subspaces other
than $0$ and $U$, let $H\subset G$ be a subgroup such that every $g\in G$ can be written $g=hz$ with $h\in H$ and $z$ commuting
with $H$, and let $(\tau,T)$ be an irreducible representation of $H$ with a nonzero $H$-equivariant linear map
$f\colon T\to V$ with values in $U$. Then:
\begin{enumerate}[label=(\alph*),nosep]
\item every $u\in U$ is a finite sum $\sum_ic_i(x_i)$ with $H$-equivariant linear maps $c_i\colon T\to V$ with values in
  $U$ and $x_i\in T$;
\item for every $0\ne u\in U$ there is a linear map $p\colon V\to T$ with $p(u)\ne0$ and $p(\rho(h)y)=\tau(h)p(y)$ for $h\in H$,
  $y\in U$.
\end{enumerate}
\formallabel{UpperTheory.TP.absIso}
\end{lemma}

\begin{proof}
Let $Z$ be the centraliser of $H$ and $\Sigma$ the span of the vectors $\rho(z)f(x)$, $z\in Z$, $x\in T$. Each $\rho(z)\circ f$ is
$H$-equivariant, and $\Sigma\subset U$. $\Sigma$ is $G$-invariant: for $g=hz'$ ($h\in H$, $z'\in Z$; then also $h^{-1}gz\in Z$ for
$z\in Z$), $\rho(g)\rho(z)f(x)=\rho(h^{-1}gz\,h)f(x)=\rho(h^{-1}gz)f(\tau(h)x)$. As $f\ne0$, $\Sigma\ne0$, so $\Sigma=U$; this is (a).
By (a), $U$ is contained in the $\tau$-isotypic component $I$ of $V$ restricted to $H$, which is isomorphic to a direct
sum of copies of $\tau$. For $0\ne u\in U$, some coordinate projection $I\to T$ in such a decomposition is $H$-equivariant
and does not vanish on $u$; extend it linearly from $I$ to $V$.
\end{proof}

\begin{proposition}[Reading off the lift]\label{pf:tlift:readlift}
Let $U$ be an automorphic representation of $\Uni_n$, $U'$ one of $\Uni_d$, and $\mathrm{pr}$ a linear map with $\mathrm{pr}(f)\in U'$ and
$\mathrm{pr}(R(x)f)=R(x)\mathrm{pr}(f)$ for automorphic forms $f$ on $\Uni_d$. Let $f_0\in\tliftPol$, $\varphi\in U$ and $\phi$ finite data with
$\mathrm{pr}(\tliftlift(f_0\otimes\phi,\varphi))\ne0$. Then:
\begin{enumerate}[label=(\alph*),nosep]
\item if $f_0$ has type $\lambda'$ at the infinite place $w$ (\cref{pf:tlift:hastype}), then $U'$ has weight $\lambda'$ at
  $\tau_w$;
\item for every split place $v$, $t\colon E\to F_v$ and local components $\tau$ of $U$ and $\tau'$ of $U'$ at $w_t$,
  $\tliftLP_t(\chi_X,\chi_W,\tau,\tau')$ holds.
\end{enumerate}
\formallabel{UpperTheory.TP.readLift_k}\formallabel{UpperTheory.TP.real_part_k}\formallabel{UpperTheory.TP.split_part_k}
\end{proposition}

\begin{proof}
(a) Let $i\colon V_{\lambda'}\to\tliftPol$ and $x'$ with $i(x')=f_0$ be as in \cref{pf:tlift:hastype}, and
$J(x)=\mathrm{pr}(\tliftlift(i(x)\otimes\phi,\varphi))$, a linear map $V_{\lambda'}\to U'$ (\cref{pf:tlift:liftaut}(a),(b)) with $J(x')\ne0$. For
$k\in\Uni_d(F_w)$, $\omega(1,e'(k))(i(x)\otimes\phi)=i(V_{\lambda'}(\tau_w(k))x)\otimes\phi$, since $\omega(1,e'(k))$ acts at infinity by
$i(x)\mapsto i(V_{\lambda'}(\tau_w(k))x)$; hence by \cref{pf:tlift:liftaut}(c) with $g=1$ and the equivariance of $\mathrm{pr}$,
$J(V_{\lambda'}(\tau_w(k))x)=R(k)J(x)$. As $V_{\lambda'}$ is irreducible, $U'$ has weight $\lambda'$ at $\tau_w$.

(b) Let $w=w_t$. The subgroup $\Uni_n(F_v)\cong\GL_n(E_w)$ and the elements of $\Uni_n(\A)$ trivial above $v$ satisfy the
hypothesis of \cref{pf:tlift:absiso}, and similarly for $\Uni_d$. Let $P(y)=\mathrm{pr}(\tliftlift(f_0\otimes\phi,y))$ for $y\in U$; it maps
$U$ to $U'$, additively. By \cref{pf:tlift:absiso}(a) applied to $\varphi\in U$ and the local component $\tau$, $\varphi$ is a sum of
values $c_i(x_i)$ of equivariant maps $c_i\colon\tau\to U$; by \cref{pf:tlift:absiso}(b) applied to $0\ne P(\varphi)\in U'$ and
$\tau'$, there is an equivariant $p\colon\{\text{functions on }\Uni_d(\A)\}\to\tau'$ with $p(P(\varphi))\ne0$. Hence some
$c=c_i$, $x_0=x_i$ satisfy $p(P(c(x_0)))\ne0$. Let $\iota_t$ be a linear model at $v$ (\cref{pf:tlift:split}) and, for
$\Psi\in\tliftS(M_{n\times d}(F_v))$, let $\phi^\Psi$ be the finite data obtained from $\phi$ by replacing $\phi_v$ with $\iota_t^{-1}\Psi$.
Define
\[
  \beta(\Psi)(x)=p\Bigl(\mathrm{pr}\bigl(\tliftlift(f_0\otimes\phi^\Psi,c(x))\bigr)\Bigr)\qquad(x\in\tau).
\]
Then $\beta(\Psi)$ is antilinear in $x$ (\cref{pf:tlift:liftaut}(b)) and $\beta$ is linear in $\Psi$ (pure tensors are linear in
each factor). $\beta(\iota_t\phi_v)(x_0)=p(P(c(x_0)))\ne0$, so $\beta\ne0$. For equivariance let $g\in\GL_n(F_v)$, $h\in\GL_d(F_v)$, and
$\hat g\in\Uni_n(F_v)$, $\hat h\in\Uni_d(F_v)$ with $c_w(\hat g)=g$, $c_w(\hat h)=h$ (under $F_v\cong E_w$). Then $e(\hat g)$, $e'(\hat h)$
are supported at $v$ and $\pi_t(e(\hat g))=g$, $\pi_t(e'(\hat h))=h$ (\cref{pf:tlift:brex}); if $\omega(e(\hat g),e'(\hat h))$ acts at
$v$ by $\iota_t^{-1}\Psi\mapsto\psi_1$, then by \cref{pf:tlift:split} $\iota_t\psi_1=\omega^{\chi_{W,t}^{-1},\chi_{X,t}}(g,h)\Psi$, i.e.
$\omega(e(\hat g),e'(\hat h))(f_0\otimes\phi^\Psi)=f_0\otimes\phi^{\Psi_1}$ with $\Psi_1=\omega^{\chi_{W,t}^{-1},\chi_{X,t}}(g,h)\Psi$. By
\cref{pf:tlift:liftaut}(c), the equivariance of $c$, $\mathrm{pr}$ and $p$:
\begin{align*}
  \beta(\Psi_1)(\tau(g)x)&=p\bigl(\mathrm{pr}(\tliftlift(\omega(e(\hat g),e'(\hat h))(f_0\otimes\phi^\Psi),R(\hat g)c(x)))\bigr)\\
  &=p\bigl(R(\hat h)\,\mathrm{pr}(\tliftlift(f_0\otimes\phi^\Psi,c(x)))\bigr)=\tau'(h)\,\beta(\Psi)(x).
\end{align*}
This is \eqref{eq:tlift:locpair}.
\end{proof}

% ---------------------------------------------------------------------------
\subsection{Proof of \texorpdfstring{\cref{pf:tlift:main}}{the theorem}}\label{pf:tlift:proof}

\begin{proposition}[Assembly]\label{pf:tlift:assembly}
\Cref{pf:tlift:main} follows from \cref{pf:tlift:toolsthm,pf:tlift:hgtex,pf:tlift:hgtfst,pf:tlift:brex,pf:tlift:realpair,%
pf:tlift:nonvanthm,pf:tlift:liftaut,pf:tlift:projex,pf:tlift:readlift}.
\formallabel{UpperTheory.TP.tLift_of_k}\formallabel{UpperTheory.TP.tLift_of_tools_k}
\end{proposition}

\begin{proof}
Let $E$, the base data, $n$ and $X$ be given. Let $\tliftK$ be the standard maximal compact subgroup of $G_2(X)(\A)$,
$\tliftht$ its height and $\iota$ the embedding (\cref{pf:tlift:hgtex}), and let $t_0$ be as in \cref{pf:tlift:hgtfst}. Choose
an integer $N>2t_0$ and put $d_0=\max(N,2n)$.

Let $d\ge d_0$ and $W,\delta,\chi_X,\chi_W,m_X,m_W,U,(\lambda_w)$ be as in the theorem. Then $2n\le d$ and $t_0<N/2\le d/2$. Also $\delta\ne0$:
if $\delta=0$ then $i\tau_w(J_W)=0$ could not have signature $(d,0)$ with $d\ge1$. By \cref{pf:tlift:pairex} there is a
Weil pair $(M,\chi_X,\chi_W,\omega)$ for $(X,W)$, and by \cref{pf:tlift:brex} there are $e=e_X$ and $e'=e_W$. By
\cref{pf:tlift:realpair} (applicable as $n\le2n\le d$) there are $0\ne f_0\in\tliftPol$, $0\ne\varphi\in U$ of the same type
and $\lambda'_w$ with $\tliftLW(\lambda_w,m_X(w),m_W(w),\lambda'_w)$ such that $f_0$ has type $\lambda'_w$ at $w$. By
\cref{pf:tlift:hgtfst}, $g\mapsto\tliftht(\iota(g,1))^t$ is $\mu$-integrable for $t\ge t_0$; as $2n\le d$ and $t_0<d/2$,
\cref{pf:tlift:nonvanthm} gives finite data $\phi$ with $\tliftlift_M(f_0\otimes\phi,\varphi^\flat)\ne0$. Since $e'$ is bijective,
$\psi=\tliftlift(f_0\otimes\phi,\varphi)\ne0$, and $\psi$ is an automorphic form on $\Uni_d$ by \cref{pf:tlift:liftaut}. By
\cref{pf:tlift:projex} there are $U'$ and $\mathrm{pr}$ with $\mathrm{pr}(\psi)\ne0$. By \cref{pf:tlift:readlift}, $U'$ has weight $\lambda'_w$ at
$\tau_w$ for every $w$, and $\tliftLP_t(\chi_X,\chi_W,\tau,\tau')$ holds for all split places and local components. Thus $U'$ is a
lift of $U$.
\end{proof}

\begin{proof}[Proof of \cref{pf:tlift:main}]
By \cref{pf:tlift:toolsthm} the tool package holds, \cref{pf:tlift:stab2cpt,pf:tlift:hgtex} give the heights, and
\cref{pf:tlift:hgtfst,pf:tlift:brex,pf:tlift:realpair,pf:tlift:nonvanthm,pf:tlift:liftaut,pf:tlift:projex,pf:tlift:readlift}
are proved above from the tool package; \cref{pf:tlift:assembly} concludes.
\end{proof}

